% concatenated from H2H2-CSC-hypersurface--20260623.tex
%%%%%%%%%%%%%%%%%%%%%%%%%%%%%%%%%%%%%%%%%%%%%%%%%%%%%%%%%%%%%%%
%%%%       Hypersurfaces with constant sectional curvature %%%%
%%%%          of $\mathbb{H}^2\times \mathbb{H}^2$         %%%%
%%%%                     2026-06-23                       %%%%
%%%%%%%%%%%%%%%%%%%%%%%%%%%%%%%%%%%%%%%%%%%%%%%%%%%%%%%%%%%%%%%
\documentclass[reqno]{amsart}%[a4paper,10pt,reqno]{amsart}

\usepackage[english]{babel}
\usepackage{amsmath,amssymb,amsthm}
\usepackage{hyperref}
\usepackage[centering]{geometry}
%\usepackage{showkeys}
\usepackage{xcolor}

%\textwidth=13.5cm
\newcommand{\metric}[2]{\ensuremath{\langle #1, #2\rangle}}  %% command for the metric  < , >
\newcommand{\nks}{\ensuremath{S^3\times S^3}}   %%  abbreviation for S3 x S3
\newcommand{\lcc}{\ensuremath{\tilde \nabla}}   %%  Levi Civita connection on S3 x S3
\renewcommand{\ne}{\ensuremath{\nabla^E}}   %%  Levi Civita connection on S3 x S3
%\renewcommand{\epsilon}{\varepsilon}            %%  changes epsilon
\newcommand{\ijk}{\ensuremath{\epsilon_{ijk}}}  %%  Levi-Civita symbol
\renewcommand{\aa}{\ensuremath{\alpha\cdot\alpha}}  %%  shorthand notation for \alpha\cdot\alpha, etc.
\newcommand{\ba}{\ensuremath{\beta\cdot\alpha}}     %%
\newcommand{\ab}{\ensuremath{\alpha\cdot\beta}}     %%
\newcommand{\bb}{\ensuremath{\beta\cdot\beta}}      %%
\newcommand{\dz}{\ensuremath{{\partial_{z}}}}
\newcommand{\dzb}{\ensuremath{{\partial_{\bar{z}}}}}
\newcommand{\du}{\ensuremath{{\partial_u}}}
\newcommand{\dv}{\ensuremath{{\partial_v}}}
\newcommand{\F}{\ensuremath{\mathcal{F}}}
\newcommand{\G}{\ensuremath{\mathcal{G}}}
\newcommand{\Z}{\ensuremath{\mathbb{Z}}}
\newcommand{\R}{\ensuremath{\mathbb{R}}}
\newcommand{\C}{\ensuremath{\mathbb{C}}}
\renewcommand{\H}{\ensuremath{\mathbb{H}}}


%%%%%%%%%%%%%%%%%%%%%%%%%%%
\newcommand{\red}{\textcolor{red}}
\newcommand{\blue}{\textcolor{blue}}
\newcommand{\magenta}{\textcolor{magenta}}
%%%%%%%%%%%%%%%%%%%%%%%%%%

\DeclareMathOperator{\re}{Re}    %% Real part
\DeclareMathOperator{\im}{Im}    %% Imaginary part
\DeclareMathOperator{\arccot}{arccot}    %% Imaginary part
\DeclareMathOperator{\csch}{csch}    %% Imaginary part

\newtheorem{theorem}{Theorem}[section]   %% definition of theorem environment
\newtheorem*{theorem*}{Theorem}          %% a theorem environment without numbering
\newtheorem{lemma}[theorem]{Lemma}
\newtheorem{proposition}[theorem]{Proposition}

\theoremstyle{definition}
\newtheorem{definition}[theorem]{Definition}
\newtheorem{corollary}[theorem]{Corollary}
\newtheorem{example}[theorem]{Example}
\newtheorem{exercise}{Exercise}
\newtheorem{question}{Question}
%\newtheorem*{remark}{Remark}
\newtheorem{remark}{Remark}[section]
\newtheorem*{acknow}{Acknowledgments}

\numberwithin{equation}{section}


\title[Hypersurfaces of $\mathbb{H}^2\times\mathbb{H}^2$ with constant sectional curvature]
{Hypersurfaces of $\mathbb{H}^2\times\mathbb{H}^2$ with constant sectional curvature}
\thanks{}

\author{Haizhong Li, Luc Vrancken, Xianfeng Wang \and Zeke Yao}

\subjclass[2020]{Primary 53C42; Secondary 53B25}
%\subjclass[2010]{Primary 53C42; Secondary 53B25}
%\subjclass[2010]{Primary 53C42; Secondary 53C24, 53B25, 53C40} % 53C55
\keywords{Riemannian product manifold, hypersurface with constant sectional curvature, minimal hypersurface, cosh-Gordon equation, sinh-Gordon equation}

\date{}
\begin{document}

\begin{abstract}
In this paper, we classify the hypersurfaces
of $\mathbb{H}^2\times\mathbb{H}^2$ with constant sectional curvature. 
In contrast to $\mathbb{S}^2\times\mathbb{S}^2$, the resulting examples for $\mathbb{H}^2\times\mathbb{H}^2$  exhibit more diversity, and we construct a special  example with non-constant product angle function. For  $\mathbb{S}^2\times\mathbb{S}^2$, however, the product angle function of any constant sectional curvature  hypersurface is identically zero. 
 %The main results in the paper can be regarded as the counterpart of [13] by the authors. However, due to the non-compactness of $\mathbb{H}^2\times\mathbb{H}^2$, there are
%essential differences between the hyperbolic case and the sphere case.
%The product angle function $C$ of any hypersurface of $\mathbb{S}^2\times\mathbb{S}^2$
%with constant sectional curvature is identically zero,
%while in the hyperbolic case, there exists an example which have constant sectional curvature and non-constant product angle function.
As a byproduct, we classify the hypersurfaces of  $\mathbb{H}^2\times\mathbb{H}^2$ with constant product angle function and constant mean curvature (or constant scalar curvature).
\end{abstract}

\maketitle

%%%%%%%%%%%%%%%%%%%%%%
\section{Introduction}\label{sect:1}


From the point of view of Riemannian geometry, the  classification problem of submanifolds with constant sectional curvature in different ambient spaces is an important and  attractive topic in differential geometry. 
When the ambient space is a real space form, the problem has a long history and has been well understood. It is remarkable that explicit examples of such 
submanifolds were constructed in \cite{DT2002} by  developing the classical Ribaucour transformation. We refer to the monograph \cite{DT2019} for more details of classification of  constant sectional curvature hypersurfaces in real space forms. 
In non-flat complex space forms and some other irreducible symmetric spaces, there are also some important classification results for hypersurfaces with constant sectional curvature or Einstein hypersurfaces (see, e.g.,  \cite{CR1982,I-R,K-R,Kon,M,N-P} and the references therein).
There are also classifications for constant sectional curvature  hypersurfaces of cylinders over space forms 
(see, e.g., \cite{AEG,CT2025,M-T} and the references  therein). 

  
Howerver, in general, if the curvature tensor of the ambient space is not simple, it is a very challenging problem to obtain
the classification of hypersurfaces with
constant sectional curvature, if one uses the Gauss and Codazzi
equations directly. In this aspect, the first three authors  introduced a new approach  in 2013,
the so-called {\it Tsinghua principle}, which applies the Codazzi equation
and the Ricci identity in an indirect way to obtain some nice linear equations involving
the components of the second fundamental form.
This new  approach has been  applied successfully to classsify submanifolds with constant sectional curvature in different ambient spaces (see, e.g., \cite{A-L-V-W,LMVVW,LVWY}).


 In recent years, the study of the submanifolds of $\mathbb{S}^2\times\mathbb{S}^2$ and
 $\mathbb{H}^2\times\mathbb{H}^2$ have become very active.
Both  $\mathbb{S}^2\times\mathbb{S}^2$ and
$\mathbb{H}^2\times\mathbb{H}^2$  are very typical 4-dimensional Hermitian symmetric manifolds.
When the codimension is  $2$, there have been many interesting results on minimal surfaces (cf. \cite{C-U,GVWX,LMVVW,T-U2}, etc.), and surfaces with parallel mean curvature vector   (cf. \cite{T-U1}, etc.) in $\mathbb{S}^2\times\mathbb{S}^2$ or
$\mathbb{H}^2\times\mathbb{H}^2$.
 When the codimension is $1$,  it is remarkable that the homogeneous hypersurfaces and the isoparametric hypersurfaces of $\mathbb{S}^2\times\mathbb{S}^2$ and $\mathbb{H}^2\times\mathbb{H}^2$ were classified in  \cite{Ur} and \cite{GMY}, respectively. There have also been some classification results on Hopf hypersurfaces and the hypersurfaces with isometric Reeb flow of  $\mathbb{S}^2\times\mathbb{S}^2$ and $\mathbb{H}^2\times\mathbb{H}^2$ (see \cite{GHMY} and \cite{ZGHY}, etc.).  In particular, the authors of the present paper classified the hypersurfaces of $\mathbb{S}^2\times\mathbb{S}^2$
 with constant sectional curvature in  \cite{LVWY}. It turns out that any such hypersurface is a parallel hypersurface of a minimal hypersurface in $\mathbb{S}^2\times\mathbb{S}^2$, and the aothors established a one-to-one correspondence between such minimal hypersurface and the solution to the famous
 ``sinh-Gordon equation''
 $
 (\frac{\partial^2}{\partial u^2}+\frac{\partial^2}{\partial v^2})h
 =-\tfrac{1}{\sqrt{2}}\sinh(\sqrt{2}h).
 $ 
 The main purpose of this paper is to classify the hypersurfaces of  $\mathbb{H}^2\times\mathbb{H}^2$ with constant sectional curvature. Due to the curvature properties and  non-compactness of $\mathbb{H}^2\times\mathbb{H}^2$, 
in contrast to  $\mathbb{S}^2\times\mathbb{S}^2$, the resulting examples for $\mathbb{H}^2\times\mathbb{H}^2$ exhibit more  diversity. 


Before stating our main results, we first recall some basic properties of the hypersurfaces in $\mathbb{H}^2\times\mathbb{H}^2$. Note that there is a natural  product
structure $P$ on $\mathbb{H}^2\times\mathbb{H}^2$ defined by $$P(v_1,v_2)=(v_1,-v_2)$$ for
any tangent vector fields $v_1, v_2$ on $\mathbb{H}^2$. 
For any orientable hypersurface $M$ of $\mathbb{H}^2\times\mathbb{H}^2$ with
a unit normal vector field $N$,  we can introduce an important function $C$ defined by
$$C=g(PN, N),$$ 
where $g$ denotes the standard product metric
on $\mathbb{H}^2\times\mathbb{H}^2$. %Denote $V:=PN-CN$ the tangential part of $PN$.
The geometry of hypersurfaces in $\mathbb{H}^2\times\mathbb{H}^2$ is closely related
to this function $C$. Hereafter, we always assume that $M$ is an orientable hypersurfaecs of $\mathbb{H}^2\times\mathbb{H}^2$, and we call $C$ the
{\it product angle function} of $M$.

As the first main result in this paper, we classify  the hypersurfaces of $\mathbb{H}^2\times\mathbb{H}^2$ with constant sectional curvature.
\begin{theorem}\label{thm:1.1}
Let $M$ be a hypersurface of $\mathbb{H}^2\times \mathbb{H}^2$ with constant sectional curvature $\kappa$.
Then $\kappa=-1/2$ and $M$ is either
\begin{enumerate}
\item[(1)]
an open part of Example \ref{exam:4.1ss1ss11ab1}; or

\item[(2)]
an open part of $M_{k_1,1/k_1}^{1/2}$ given in Example \ref{exam:3.4fg1} for some constant $0<k_1\leq1$; or

\item[(3)]
an open part of one of the hypersurfaces described by Theorems \ref{thm:5.10a1}--\ref{thm:5.13a1}.

\end{enumerate}
\end{theorem}

\begin{remark}\label{rem:1.1}
(i) Compared with $\mathbb{S}^2\times \mathbb{S}^2$, there is a hypersurface of $\mathbb{H}^2\times \mathbb{H}^2$ with constant sectional curvature and non-constant product angle function,
which is Case (1) in Theorem \ref{thm:1.1}. The examples in 
Case (2) and Case (3) of Theorem \ref{thm:1.1} have constant product angle function $C=0$.

(ii) In Theorem \ref{thm:5.10a1} (or Theorem \ref{thm:5.11a1}), for any function depending on one variable,
one can construct a hypersurface of $\mathbb{H}^2\times \mathbb{H}^2$ with constant sectional curvature
$-1/2$. In fact, for any function $\beta_1:I\rightarrow \mathbb{R}$ (or $\beta_2:I\rightarrow \mathbb{R}$), we construct a metric $g$ and two $(1,1)$-tensors $A,~\tilde{P}$ as \eqref{eqn:3.1ff3}--\eqref{eqn:3.1ff5}
(or \eqref{eqn:3.1ffff3}--\eqref{eqn:3.1ffff5}), then by the existence theorem
and uniqueness theorem in \cite{K,L-T-V}, up to an isometry of $\mathbb{H}^2\times\mathbb{H}^2$, there exists a unique hypersurface in $\mathbb{H}^2\times\mathbb{H}^2$ with constant sectional curvature
such that $C=0$ and its induced metric and shape operator  are given by $g$ and $A$,
respectively.

(iii) In Theorem \ref{thm:5.12a1}, for any function $h_3:(u,v)\to\mathbb{R}$ satisfying
the ``cosh-Gordon equation''
%
\begin{equation*}%\label{eqn:3.1fffff2}
\frac{\partial^2 h_3}{\partial u \partial v}
=4\sqrt{2}\cosh(h_3),
\end{equation*}
we construct a metric $g$ and two
$(1,1)$-tensors $A,~\tilde{P}$ (see \eqref{eqn:3.1fffff3}, \eqref{eqn:3.1fffff4} and \eqref{eqn:3.1fffff5}), then by the existence theorem
and uniqueness theorem in \cite{K,L-T-V}, up to an isometry of $\mathbb{H}^2\times\mathbb{H}^2$, there exists a unique hypersurface in $\mathbb{H}^2\times\mathbb{H}^2$ with constant sectional curvature
such that $C=0$ and its induced metric and shape operator  are given by $g$ and $A$,
respectively.

(iv) In Theorem \ref{thm:5.13a1}, for any non-zero function $h_4:(u,v)\to\mathbb{R}$ satisfying
the ``sinh-Gordon equation''
%
\begin{equation*}%\label{eqn:3.1ffffff2}
\frac{\partial^2 h_4}{\partial u \partial v}
=4\sqrt{2}\sinh(h_4),
\end{equation*}
we construct a metric $g$ and two
$(1,1)$-tensors $A,~\tilde{P}$ (see \eqref{eqn:3.1ffffff3}, \eqref{eqn:3.1ffffff4} and \eqref{eqn:3.1ffffff5}), then by the existence theorem
and uniqueness theorem in \cite{K,L-T-V}, up to an isometry of $\mathbb{H}^2\times\mathbb{H}^2$, there exists a unique hypersurface in $\mathbb{H}^2\times\mathbb{H}^2$ with constant sectional curvature
such that $C=0$ and its induced metric and shape operator  are given by $g$ and $A$,
respectively.

(v) In \cite{LVWY}, it was proved that any hypersurface of $\mathbb{S}^2\times\mathbb{S}^2$
with constant sectional curvature must be a parallel hypersurface of a minimal hypersurface with $C=0$, and there is 
a one-to-one correspondence between such minimal hypersurface and the solution to the
``sinh-Gordon equation''
$
(\frac{\partial^2}{\partial u^2}+\frac{\partial^2}{\partial v^2})h
=-\tfrac{1}{\sqrt{2}}\sinh(\sqrt{2}h).
$
However, this phenomenon is no longer true in $\mathbb{H}^2\times\mathbb{H}^2$. 



\end{remark}

As the second main result, we obtain a complete classification of
the hypersurfaces of $\mathbb{H}^2\times\mathbb{H}^2$
with constant product angle function and constant mean curvature.
\begin{theorem}\label{thm:1.2}
Let $M$ be a hypersurface of $\mathbb{H}^2\times \mathbb{H}^2$ with constant product angle function $C$. Then
$M$ has constant mean curvature if and only if either
%
\begin{enumerate}
\item[(1)]
$M$ is an open part of $M_\Gamma$ given in Example \ref{exam:3.1f1f2}, where $\Gamma$ is a curve of $\mathbb{H}^2$ with constant curvature; or

\item[(2)]
$M$ is an open part of $M_\tau$ given in Example \ref{exam:3.8f1f2} for some $\tau<-1$; or

\item[(3)]
$M$ is an open part of $M_{1,1}^c$ given in Example \ref{exam:3.4fg1} for some $c\in(0,1)$; or

\item[(4)]
$M$ is an open part of $M_{1,-1}^c$ given in Example \ref{exam:3.4fg1} for some $c\in(0,1)$; or

%\item[(5)]
%$M$ is an open part of $M_{1,1}^{1/2}$; or

\item[(5)]
$M$ is an open part of $M_{k_1,k_1}^{1/2}$ given in Example \ref{exam:3.4fg1} for some constant $k_1\neq1$; or

\item[(6)]
$M$ is an open part of a minimal hypersurface described by Theorem \ref{thm:3.1}.
\end{enumerate}
\end{theorem}

\begin{remark}\label{rem:1.2}
In Theorem \ref{thm:3.1}, we establish a one-to-one correspondence between the minimal hypersurface $M$ in Case (6) of Theorem \ref{thm:1.2} and the non-zero solution to the ``sin-Gordon equation'' 
\begin{equation*}%\label{eq:2.ab16}
(\frac{\partial^2}{\partial u^2}+\frac{\partial^2}{\partial v^2})h
=\tfrac{1}{\sqrt{2}}\sin(\sqrt{2}h).
\end{equation*}
In fact, we prove that the hypersurface $M$ in Theorem \ref{thm:3.1} admits coordinates $(u,v,t)$ and there is a non-zero function $h:(u,v)\to\mathbb{R}$ such that the components of the
second fundamental form of $M$ are given by \eqref{eq:2.ab15} with $h$ satisfying the above equation.
Conversely, for any non-zero solution $h:(u,v)\rightarrow \mathbb{R}$, we construct a metric $g$ and two
$(1,1)$-tensors $A,~\tilde{P}$ (see \eqref{eq:2.ab17}, \eqref{eqn:A} and \eqref{defP}), then by the existence theorem
and uniqueness theorem in \cite{K,L-T-V}, up to an isometry of $\mathbb{H}^2\times\mathbb{H}^2$, there exists a unique minimal hypersurface in $\mathbb{H}^2\times\mathbb{H}^2$
such that $C=0$ and its induced metric and shape operator  are given by $g$ and $A$,
respectively.
\end{remark}

\begin{remark}\label{rem:1.3}
The conditions of constant product angle function and constant mean curvature in Theorem \ref{thm:1.2} are necessary.

(i) In fact, we can define the map $\Phi:\Omega\subset \mathbb{R}^3\hookrightarrow\mathbb{H}^2\times\mathbb{H}^2$ as follows:
%
\begin{equation}\label{eqn:1.1a1a1}
\begin{aligned}
&p(t,r)=\cosh r\,(\cosh t,\sinh t,0)+\sinh r\,(0,0,1),\\
&q(t,s)=\cosh s\,(\cosh t,\sinh t,0)+\sinh s\,(0,0,1).
\end{aligned}
\end{equation}
Then hypersurface $\Phi(\Omega)$ is minimal, but its product angle function is $C=\frac{\cosh^2 s-\cosh^2 r}{\cosh^2 s+\cosh^2 r}$, which is not constant.

%
(ii) In Example \ref{exam:3.4fg1}, when $k_1(r)$ and $k_2(s)$ are non-constant,
then the hypersurface $M_{k_1,k_2}^c$ has constant product angle function $C=1-2c$,
but it does not have constant mean curvature.
\end{remark}







As the third main result, we obtain a complete classification of
the hypersurfaces of $\mathbb{H}^2\times\mathbb{H}^2$
with constant product angle function and constant scalar curvature.

\begin{theorem}\label{thm:1.3}
Let $M$ be a hypersurface of $\mathbb{H}^2\times \mathbb{H}^2$ with constant product angle function $C$. Then
$M$ has constant scalar curvature if and only if either
\begin{enumerate}
\item[(1)]
$M$ is an open part of $M_\Gamma$ given in Example \ref{exam:3.1f1f2},  where $\Gamma$ is a curve of $\mathbb{H}^2$ with constant curvature; or

\item[(2)]
$M$ is an open part of $M_\tau$ given in Example \ref{exam:3.8f1f2} for some $\tau<-1$; or

\item[(3)]
$M$ is an open part of $M_{1,1}^c$ given in Example \ref{exam:3.4fg1} for some $c\in(0,1)$; or

\item[(4)]
$M$ is an open part of $M_{1,-1}^c$ given in Example \ref{exam:3.4fg1} for some $c\in(0,1)$; or

%\item[(5)]
%$M$ is an open part of $M_{1,1}^{\frac{1}{2}}$; or

\item[(5)]
$M$ is an open part of $M_{k_1,1/k_1}^{1/2}$ given in Example \ref{exam:3.4fg1} for some constant $0<k_1<1$; or

\item[(6)]
$M$ is an open part of one of the hypersurfaces described by Theorems \ref{thm:5.10a1}--\ref{thm:5.13a1}.

\end{enumerate}
\end{theorem}

\begin{remark}\label{rem:1.4}
The conditions of constant product angle function and constant scalar curvature in Theorem \ref{thm:1.3} are necessary.

(i) In fact, Example \ref{exam:4.1ss1ss11ab1} has constant scalar curvature, but its product angle function $C$ is not constant.

%
(ii) In Example \ref{exam:3.4fg1}, when $k_1(r)$ and
$k_2(s)$ are non-constant, then the
hypersurface $M_{k_1,k_2}^c$ has constant product angle function $C=1-2c$, but
it does not have constant scalar curvature.
\end{remark}


The paper is organized as follows. In Sect. \ref{sect:2}, we  collect some basic properties of $\mathbb{H}^2\times \mathbb{H}^2$ and some preliminaries of the geometry of the hypersurfaces of $\mathbb{H}^2\times \mathbb{H}^2$. In Sect. \ref{sect:3}, we introduce some canonical examples of hypersurfaces in
$\mathbb{H}^2\times\mathbb{H}^2$. In particular, we construct a special  example with constant sectional curvature and non-constant product angle function (see Example \ref{exam:4.1ss1ss11ab1}). In Sect. \ref{sect:4}, we give a complete classification of the minimal hypersurfaces in $\mathbb{H}^2\times \mathbb{H}^2$ with $C=0$ (see Theorem \ref{thm:3.2aa}). %It turns out that the resulting hypersurfaces can only be 
%$M_{k_1,k_1}^{1/2}$ for some constant $k_1$ or one of the hypersurfaces described by Theorem \ref{thm:3.1}. 
In Sect. \ref{sect:5}, we present the proof of Theorem \ref{thm:1.1}. One of the key ingredients is to prove that the constant sectional curvature can only be $-\frac{1}{2}$, which is achieved by applying the so-called Tsinghua principle 
(see Proposition \ref{prop:4.5}). In Sect. \ref{sect:6}, we present the proofs of Theorems \ref{thm:1.2} 
and \ref{thm:1.3}. In Sect. \ref{sect:7}, we give a complete classification of the hypersurfaces 
in $\mathbb{S}^m\times \mathbb{S}^n$ ($m\geq3$, $n\geq2$), 
$\mathbb{H}^m\times \mathbb{H}^n$ ($m\geq3$, $n\geq2$), and $\mathbb{S}^m\times \mathbb{H}^n$ ($m,n\geq2$) with constant sectional curvature. 




\vskip 2mm

\textbf{Acknowledgments:} 
H. Li was supported by NSFC Grant No.12471047 and NSFC-FWO W2521103. 
L. Vrancken was supported by NSFC-FWO mobility project VS05726N. 
X. Wang was supported by Young Scientific and Technological Talents (Level Two) in Tianjin, Natural Science Foundation of Tianjin, China (Grant No. 25JCZDJC01110)  and the Fundamental Research Funds for the Central Universities (Grant No. 050-63261176). Z. Yao was supported by NSFC Grant No. 12401061.


%%%%%%%%%%%%%%%%%%%%%%%%%%%%%%%%%%%%%%%%%%%%%
\section{Preliminaries} \label{sect:2}
%%%%%%%%%%%%%%%%%%%%%%%%%%%%%%%%%%%%%%%%%%%%%
%%%%%%%%%%%%%%%%%%%%%%%%%%%%%%%%%%%%%%%%%%%%%
\subsection{The geometric structure on $\mathbb{H}^{2}\times \mathbb{H}^{2}$}~

Let $\mathbb{R}_{1}^{3}$ be the three-dimensional Minkowski space with the
Lorentzian metric $\langle\cdot,\cdot\rangle $.
The hyperbolic plane of curvature $-1$ can be defined as the following subset
of $\mathbb{R}_{1}^{3}$: 
$$
\mathbb{H}^{2}=\left\{\left(x_{1}, x_{2}, x_{3}\right) \in \mathbb{R}_{1}^{3}
\mid-x_{1}^{2}+x_{2}^{2}+x_{3}^{2}=-1  , x_{1}>0\right\}.
$$
The standard complex structure $J$ on $\mathbb{H}^{2}$ is defined by
$$J_xu=x\boxtimes u,$$
for all $x\in \mathbb{H}^{2}$ and all $u\in T_x\mathbb{H}^{2}$, where $\boxtimes$
is the Lorentzian cross product defined by
$$
(a_1,a_2,a_3)\boxtimes (b_1,b_2,b_3)=(a_3b_2-a_2b_3, a_3b_1-a_1b_3, a_1b_2-a_2b_1).
$$

Throughout the paper we will consider $\mathbb{H}^2 \times  \mathbb{H}^2$
embedded isometrically  in $\mathbb{R}_{1}^{3} \times \mathbb{R}_{1}^{3} \cong \mathbb{R}_{2}^{6}$,
with the induced Riemannian product metric which we denote by $g$.
We define two complex structures on $\mathbb{H}^2\times \mathbb{H}^2$ by
$$
J_1=(J,J),\quad J_2=(J,-J),
$$
which endow $\mathbb{H}^2 \times  \mathbb{H}^2$ with two structures of K\"ahler surface.
%It is clear that if $\text{Id}:\mathbb{H}^2\to \mathbb{H}^2$ is the identity map and
%$\mathcal{F}:\mathbb{H}^2\to \mathbb{H}^2$ is an anti-holomorphic isometry, then
%$$
%(\text{Id},\mathcal{F}) : \mathbb{H}^2\times \mathbb{H}^2\to \mathbb{H}^2\times \mathbb{H}^2
%$$
%is a holomorphic isometry from $(\mathbb{H}^2\times \mathbb{H}^2, g, J_1)$
%onto $(\mathbb{H}^2\times \mathbb{H}^2, g, J_2)$.
The isometry group of $\mathbb{H}^{2}\times \mathbb{H}^{2}$ is
$$
\operatorname{Iso}\left(\mathbb{H}^{2}\times \mathbb{H}^{2}\right)=\left\{\left(\begin{array}{cc}
A_{1} & 0 \\
0 & A_{2}
\end{array}\right),\left(\begin{array}{cc}
0 & B_1 \\
B_{2} & 0
\end{array}\right) \mid A_{1}, A_{2}, B_1, B_{2} \in \mathrm{O}^{+}(1,2)\right\},
$$
where $\mathrm{O}^{+}(1,2)$ denotes the ortochronous Lorentz group.

The product structure $P$ on $\mathbb{H}^2\times\mathbb{H}^2$ is defined by
$P: T(\mathbb{H}^2\times\mathbb{H}^2)\rightarrow T(\mathbb{H}^2\times\mathbb{H}^2)$ such that
$$
P(v_1,v_2)=(v_1,-v_2), \quad \forall\, v_1, v_2\in T\mathbb{H}^2.
$$
Obviously, we have $P=-J_1J_2=-J_2J_1$, $P^2=\mathrm{Id}$ and
%
\begin{equation}\label{eqn:2.1}
g(PY,Z)=g(Y,PZ), \quad \forall\,Y,Z\in T(\mathbb{H}^2\times\mathbb{H}^2).
\end{equation}
%
Moreover, $\bar{\nabla} P=0$, where $\bar{\nabla}$ is the Levi-Civita connection on
$\mathbb{H}^2\times\mathbb{H}^2$.

The curvature tensor $\bar{R}$ of $\mathbb{H}^2\times\mathbb{H}^2$
with the Riemannian product metric is given by
\begin{align*}
g(\bar{R}(U,Y)Z,W)=-\tfrac{1}{2}&\big\{g(Y,Z)g(U,W)
-g(U,Z)g(Y,W)\\
&+g(PY,Z)g(PU,W)-g(PU,Z)g(PY,W)\big\},
\end{align*}
where $U, Y, Z, W\in T(\mathbb{H}^2\times\mathbb{H}^2)$. Thus, $\mathbb{H}^2\times\mathbb{H}^2$
is an Einstein manifold with scalar curvature $-4$ and nonpositive sectional curvature.

\subsection{Hypersurfaces of $\mathbb{H}^{2}\times \mathbb{H}^{2}$}~

Let $M$ be an orientable hypersurface of $\mathbb{H}^2\times\mathbb{H}^2$ with
$N$ a unit normal vector field and still $g$ the induced metric on $M$.
Then, with respect to the product structure $P$, we define by
%
\begin{align*}
C&:=g(PN,N)=g(J_1 N,J_2 N),\\
X&:=PN-CN,
\end{align*}
%
the product angle function $C: M\rightarrow\mathbb{R}$ and a vector field
$X$ tangent to $M$. It is clear that $-1\leq C\leq 1$ and $\|X\|^2:=g(X,X)=1-C^2$.

For any tangential vector field $Y$ of $M$, acting by the product structure $P$,
we have the following decomposition
$$
PY=T Y+\mu(Y)N,
$$
where $T Y$ and $\mu(Y)N$ are the tangential
and normal parts of $PY$. Thus $T$ is a
tensorial field of type $(1,1)$, and $\mu$ is a $1$-form over $M$.
Moreover, $\mu(Y)=g(PY, N)$.

Let $\nabla$ be the Levi-Civita connection of the induced metric $g$ on $M$.
The Gauss and Weingarten formulae are given by 
%
\begin{align*}
\bar{\nabla}_Y Z=\nabla_Y Z+g(AY,Z)N, \quad \bar{\nabla}_YN=-AY,
\end{align*}
%
where $A$ is the shape operator of $M$.

Now, the Gauss and Codazzi equations of $M$ are given by
%
\begin{equation}\label{eqn:2.2}
\begin{aligned}
\operatorname{R}(U, Y)Z=&-\tfrac{1}{2}\big[g(Y,Z)U-g(U,Z)Y+g(TY,Z)TU
-g(TU,Z)TY\big]\\
&+g(AY,Z)AU-g(AU,Z)AY,
\end{aligned}
\end{equation}
%
\begin{equation}\label{eqn:2.3}
(\nabla_YA)Z-(\nabla_Z A)Y=-\tfrac{1}{2}\big[g(Y,X)TZ-g(Z,X)TY\big],
\end{equation}
%
where $U, Y, Z\in TM$, and $R$ denotes the curvature tensor of $M$ with respect
to the metric $g$.

It follows from \eqref{eqn:2.2} that the Ricci curvature tensor of $M$ is given by
%
\begin{equation}\label{eqn:2.4a11}
\begin{aligned}
g({\rm Ric}(Y),Z)%:&=\sum_{i=1}^3g(R(Y,E_i)E_i,Z)\\
=-\tfrac{1}{2}\big[g(Y,Z)-Cg(TY,Z)
+g(Y,X)g(Z,X)\big]+3Hg(AY,Z)-g(A^2Y,Z),
\end{aligned}
\end{equation}
%
where %$E_i\ (1\leq i\leq 3)$ is a local orthonormal frame field of $M$,
$H=\frac{1}{3}{\rm Tr}A$ denotes the mean curvature of $M$. %and we have used
%$$
%\sum_{i=1}^{3}g(PE_i,E_i)=\operatorname{tr}P-g(PN,N)=-C,
%$$
%
%$$
%\sum_{i=1}^{3}g(PY,E_i)g(PZ,E_i)=g(PY,PZ)
%-g(PY,N)g(PZ,N)=g(Y,Z)-g(Y,X)g(Z,X).
%$$
%Thus the Ricci curvature along a vector $X$ is given by
%
%\begin{equation}\label{eqn:2.5a11}
%{\rm Ric}(Y, Y)=-\tfrac{1}{2}\big[g(Y,Y)-Cg(TY,Y)
%+g(Y,X)^2\big]+Hg(AY,Y)-g(A^2Y,Y).
%\end{equation}
%
It follows that the scalar curvature $\rho$ of $M$ is given by
%
\begin{equation}\label{eqn:2.6a11}
\rho=-2+9H^{2}-\|A\|^{2}.
\end{equation}

Let $\nabla^2 A$ denote the second covariant derivative of $A$, i.e.,
$$
(\nabla^{2}A)(U,Y,Z):=\nabla_U[(\nabla_Y A)Z]-(\nabla_{\nabla_U Y}
A)Z-(\nabla_Y A){\nabla_U Z}.
$$
Then the Ricci identity states that
%
\begin{equation}\label{eqn:ric}
g((\nabla^{2}A)(U,Y,Z),W)-g((\nabla^{2}A)(Y,U,Z),W)
=-g( R(U,Y)Z,AW)-g( R(U,Y)W,AZ).
\end{equation}

%Notice that the product structure $P$ of $\mathbb{H}^2\times\mathbb{H}^2$ satisfies
%\eqref{eqn:2.1} and $\bar{\nabla}P=0$.
Now, by a similar proof as Lemma 2.1 of \cite{LVWY},
we can obtain the following Lemma \ref{lemma:2.1} which describes some properties
of the function $C$, the vector field $X$ and the covariant derivatives of $T$ and $\mu$.


\begin{lemma}\label{lemma:2.1}
Let $M$ be an orientable hypersurface of $\mathbb{H}^2\times\mathbb{H}^2$
with $A$ the shape operator associated
to the unit normal field $N$. Then the gradient of $C$, the
covariant derivatives of $X$, $T$ and $\mu$ are given by
%
\begin{equation}\label{eqn:2.4}
\begin{split}
&\nabla C=-2 AX,\ \ \nabla_{Y}X=C AY -T A Y,\\
&\left(\nabla_Y T\right) Z =g( AY,Z)X+\mu(Z) A Y,\\
&\left(\nabla_Y \mu\right) Z=C g( AY,Z)-g( TZ,AY).
\end{split}
\end{equation}
\end{lemma}
%
\begin{proof}
The formulae of the gradient of $C$ and the covariant derivative of $X$ can be directly found
in Lemma 2.1 of \cite{GMY}. For the rest two formulae, by the definitions of $T$, $\mu$
and the Gauss and Weingarten formulae, we get
%
\begin{align*}
\left(\nabla_{Y} T\right) Z &=(\bar{\nabla}_{Y} (T Z))^{\top}-T (\nabla_{Y} Z) 
=[\bar{\nabla}_{Y}(P Z-\mu(Z) N)]^{\top}-T (\nabla_{Y} Z) \\
&=[P \bar{\nabla}_{Y} Z-\mu(Z) \bar{\nabla}_{Y} N]^{\top}-T (\nabla_{Y} Z) 
=g(AY,Z)(P N)^{\top}+\mu(Z) A Y\\
&=g(AY,Z)X+\mu(Z) A Y,
\end{align*}
%
\begin{align*}
\left(\nabla_{Y} \mu\right) Z&= Y \mu(Z)-\mu\left(\nabla_{Y} Z\right) \\
&=g(\bar{\nabla}_{Y} PZ,N)+g(PZ,\bar{\nabla}_{Y} N)-g(P\bar{\nabla}_{Y} Z-g(AY,Z) PN,N)\\
&=-g(TZ,AY)+Cg(AY,Z),
\end{align*}
%
where $\cdot^\top$ denotes the tangential part.
\end{proof}

%\begin{remark}\label{rem:2.1}
%According to Lemma \ref{lemma:2.1}, if $C\neq\pm1$ is constant on $M$, then $V$ is a principal curvature
%vector field of $M$, and it satisfies $AV=0$.
%\end{remark}

\begin{lemma}\label{lemma:2.1aaa}
Let $M$ be an orientable hypersurface of $\mathbb{H}^2\times\mathbb{H}^2$
with $A$ the shape operator associated
to the unit normal field $N$. Then the Hessian of $C$,
the Laplacian of $C$ and the divergence of $X$ are given by
%
\begin{equation}\label{eqn:2.4aaa}
({\rm Hess}(C))(Y,Z)=-2g((\nabla_Y A)X,Z)-2Cg(A^2Y,Z)+2g(TAY,AZ).
\end{equation}
%
\begin{equation}\label{eqn:2.5aaa}
\triangle C=-6g(\nabla H,X)-2C\|A\|^2+2{\rm tr}(TA^2).
\end{equation}
%
\begin{equation}\label{eqn:2.6aaa}
{\rm div}X=3CH-{\rm tr}(TA).
\end{equation}
\end{lemma}
%
\begin{proof}
%
By using Lemma \ref{lemma:2.1}, the Gauss and Weingarten formulae, we get
\begin{align*}
({\rm Hess}(C))(Y,Z)&=Y(ZC)-(\nabla_YZ)C=-2Yg(AX,Z)+2g(AX,\nabla_YZ)\\
&=-2g(\nabla_Y{AX},Z)=-2g((\nabla_YA)X,Z)-2g(\nabla_Y X,AZ)\\
&=-2g((\nabla_Y A)X,Z)-2Cg(A^2Y,Z)+2g(TAY,AZ).
\end{align*}
%
\begin{align*}
\triangle C&=\sum_i({\rm Hess}(C))(e_i,e_i)
=\sum_i\{-2g((\nabla_{e_i} A)X,{e_i})-2Cg(A^2{e_i},{e_i})+2g(TA{e_i},A{e_i})\}\\
&=-2\sum_i\{g((\nabla_X A){e_i},{e_i})\}-2C\|A\|^2+2{\rm tr}(TA^2)
=-6g(\nabla H,X)-2C\|A\|^2+2{\rm tr}(TA^2).
\end{align*}
%
\begin{align*}
{\rm div}X&=\sum_ig(\nabla_{e_i}X,e_i)=\sum_ig(CAe_i-TAe_i,e_i)
=3CH-{\rm tr}(TA).
\end{align*}
\end{proof}

In \cite{GMY}, Gao, Ma and the fourth author classified the hypersurfaces with $C^2\equiv 1$ as follows.
%
\begin{lemma}[cf. Lemma 3.3 of \cite{GMY}]\label{lem:2.2}
Let $M$ be a hypersurface of $\mathbb{H}^2\times\mathbb{H}^2$ with $C^2\equiv 1$.
Then, up to isometries of $\mathbb{H}^2\times\mathbb{H}^2$, $M$ is locally the
product of a curve $\Gamma$ in $\mathbb{H}^2$ and an open
subset of $\mathbb{H}^2$.
\end{lemma}


\subsection{A canonical frame related to hypersurfaces of
$\mathbb{H}^2\times\mathbb{H}^2$}\label{sect:2.3}~

In this subsection, we assume that $|C|< 1$ holds on $M$,
and choose an appropriate local orthonormal frame fields $\{E_1, E_2, E_3\}$
of $M$ as follows:
%
\begin{equation}\label{eqn:2.5}
E_1=\frac{J_1N+J_2 N}{\sqrt{2(1+C)}},\
\ E_2=\frac{J_1N-J_2 N}{\sqrt{2(1-C)}},\ \ E_3=\frac{X}{\sqrt{1-C^{2}}}.
\end{equation}
%
This frame has the following properties:
%
\begin{equation}\label{eqn:2.6}
\left\{
\begin{aligned}
&P E_1=E_1,\ \ P E_2=-E_2,\ \ P E_3=\sqrt{1-C^{2}}N-C E_3,\ \ PN=CN+\sqrt{1-C^{2}}E_3.\\
&TE_1=E_1,\ \ TE_2=-E_2,\ \ TE_3=-CE_3,\ \ \mu(E_1)=\mu(E_2)=0,\ \ \mu(E_3)=\sqrt{1-C^{2}}.
\end{aligned}\right.
\end{equation}

We assume that
%
\begin{equation}\label{eqn:2.7}
\begin{aligned}	
&A E_1=b_1 E_1+b_2 E_2+b_3 E_3, \\
&A E_2=b_2 E_1+b_4 E_2+b_5 E_3,\\
&A E_3=b_3 E_1+b_5E_2+b_6 E_3,
\end{aligned}
\end{equation}
%
where $b_1$, $b_2$, $b_3$, $b_4$, $b_5$ and $b_6$ are smooth
functions on $M$. 
By \eqref{eqn:2.4a11}, under the frame $\{E_1,E_2,E_3\}$, the Ricci tensor is expressed by
%
\begin{equation}\label{eqn:2.8a1212}
\begin{aligned}
&{\rm Ric}(E_1)=(\tfrac{-1+C}{2}
+b_1b_4+b_1b_6-b_2^{2}-b_3^{2}) E_1+
(b_2b_6-b_3 b_5) E_2+(b_3b_4-b_2b_5) E_3, \\
&{\rm Ric}(E_2)=(b_2b_6-b_3 b_5) E_1
+(\tfrac{-1-C}{2}+b_1b_4+b_4b_6-b_2^{2}-b_5^{2}) E_2+(b_1b_5-b_2b_3) E_3,\\
&{\rm Ric}(E_3)=(b_3b_4-b_2b_5) E_1+(b_1b_5-b_2b_3) E_2+(-1+b_1b_6+b_4b_6-b_3^{2}-b_5^{2}) E_3.
\end{aligned}
\end{equation}
%	

When $C$ is constant which is not equal to $\pm1$ on $M$, we have the following lemma:

\begin{lemma}\label{lemma:2.3}
Let $M$ be a hypersurface of $\mathbb{H}^2\times \mathbb{H}^2$ with constant $C\neq\pm1$.
Then, under the local orthonormal frame fields $\{E_1, E_2, E_3\}$,
the following relations hold:
\begin{equation}\label{eqn:2.25}
\begin{aligned}
&b_3=b_5=b_6=0, \ \ A E_1=b_1 E_1+b_2 E_2,\ A E_2=b_2 E_1+b_4 E_2,\ A E_3=0,\\		
&\nabla_{E_1}E_1=b_1\sqrt{\tfrac{1-C}{1+C}}E_3,\
\nabla_{E_1}E_2=-b_2\sqrt{\tfrac{1+C}{1-C}}E_3,\ \nabla_{E_1}E_3=-b_1\sqrt{\tfrac{1-C}{1+C}}E_1
+b_2\sqrt{\tfrac{1+C}{1-C}}E_2,\\
&\nabla_{E_2} E_1=b_2\sqrt{\tfrac{1-C}{1+C}}E_3,\
\nabla_{E_2} E_2=-b_4\sqrt{\tfrac{1+C}{1-C}}E_3,\ \nabla_{E_2}E_3=-b_2\sqrt{\tfrac{1-C}{1+C}}E_1
+b_4\sqrt{\tfrac{1+C}{1-C}}E_2,\\
&\nabla_{E_3}E_1=\nabla_{E_3}E_2=\nabla_{E_3}E_3=0.
\end{aligned}
\end{equation}
Moreover, the Gauss and Codazzi equations are equivalent to the following equations:
\begin{equation}\label{eqn:2.26}
E_3b_1=-\tfrac{\sqrt{1-C^2}}{2}+b_1^{2}\sqrt{\tfrac{1-C}{1+C}}-b_2^2\sqrt{\tfrac{1+C}{1-C}},
\end{equation}
%
\begin{equation}\label{eqn:2.27}
E_3b_2=b_2\big(b_1\sqrt{\tfrac{1-C}{1+C}}-b_4\sqrt{\tfrac{1+C}{1-C}}\big),
\end{equation}
%
\begin{equation}\label{eqn:2.28}
E_3b_4=\tfrac{\sqrt{1-C^2}}{2}+b_2^{2}\sqrt{\tfrac{1-C}{1+C}}-b_4^2\sqrt{\tfrac{1+C}{1-C}},
\end{equation}
%
\begin{equation}\label{eqn:2.29}
E_1b_2-E_2b_1=0,	
\end{equation}
%
\begin{equation}\label{eqn:2.30}
E_1b_4-E_2b_2=0.
\end{equation}
\end{lemma}
%
\begin{proof}
By the assumption that $C$ is constant and $\nabla C=-2 AX$, we have $b_3=b_5=b_6=0$.

The product structures and the almost complex structures in
$\mathbb{S}^2\times\mathbb{S}^2$ and $\mathbb{H}^2\times\mathbb{H}^2$ share some same
properties, i.e., the product structures are parallel, and are symmetric with respect to the
product metric, the almost complex structures are parallel. Then the connections in \eqref{eqn:2.25}
can be obtained by a totally same way as (2.11) of \cite{LVWY}.

Because the Codazzi equation for a hypersurface of $\mathbb{H}^2\times\mathbb{H}^2$ 
is exactly opposite to the Codazzi equation for a hypersurface of $\mathbb{S}^2\times\mathbb{S}^2$, 
%Notice the opposite sign in the right side of the Codazzi equations for hypersurfaces of $\mathbb{S}^2\times\mathbb{S}^2$ and $\mathbb{H}^2\times\mathbb{H}^2$, 
thus \eqref{eqn:2.26}--\eqref{eqn:2.30} can be obtained by a similar way with
(2.12)--(2.16) of \cite{LVWY}.
\end{proof}

%==================================================
\section{Some canonical examples}\label{sect:3}
%


In this section, we introduce some canonical examples of hypersurfaces in
$\mathbb{H}^2\times\mathbb{H}^2$, where Example \ref{exam:3.1f1f2}, Example \ref{exam:3.8f1f2}
and Example \ref{exam:3.4fg1} have been appeared in \cite{GMY}.
%
\begin{example}\label{exam:3.1f1f2}
For any smooth curve $\Gamma$ of $\mathbb{H}^2$, one can define a hypersurface
in $\mathbb{H}^2\times\mathbb{H}^2$ by
$$
M_\Gamma:=\left\{(x,y)\in \mathbb{H}^2\times\mathbb{H}^2
~|~x\in \Gamma,\ y\in\mathbb{H}^2\right\}.
$$
\end{example}

Let $k_{\Gamma}$ be the curvature of $\Gamma$ in $\mathbb{H}^2$.
%Denote the unit normal vector field of $\Gamma$ in $\mathbb{H}^2$ by $N$. Then the unit normal vector field of $M_\Gamma$ is $(N,0)$. It is obvious that
Then, hypersurface $M_\Gamma$ has constant product angle function $C=1$, and its principal curvatures are $k_{\Gamma}$, $0$ and $0$.

\begin{lemma}[cf. Lemma 3.2 of \cite{GMY}]\label{lemma:3.1ff1}
For any smooth curve $\Gamma$ of $\mathbb{H}^2$, $M_\Gamma$ has constant principal curvatures
if and only if $\Gamma$ is an open part of a curve in $\mathbb{H}^2$ with constant curvature.
\end{lemma}

%Next example has been given by Example 3.8 of \cite{GMY}.

\begin{example}\label{exam:3.8f1f2}
For any given $\tau<-1$, we define $M_\tau=\{(p,q)\in
\mathbb{H}^2\times\mathbb{H}^2|\ \langle p,q\rangle=\tau\}$.
\end{example}


According to the calculations in Example 3.10 of \cite{GMY}, we know that
$M_\tau$ has three distinct constant principal curvatures
$\sqrt{\tfrac{\tau-1}{2(\tau+1)}}$, $\sqrt{\tfrac{\tau+1}{2(\tau-1)}}$ and $0$.
Its product angle function is $C=0$.
For any $\tau<-1$, $M_\tau$ cannot be a minimal hypersurface or a hypersurface of $\mathbb{H}^2\times\mathbb{H}^2$ with constant sectional curvature.


In the following, we define a class of hypersurfaces which generalizes the Example 3.4 of \cite{GMY}.

\begin{example}\label{exam:4.1}
For any two non-constant functions $f_1(t)$, $f_2(t)$ and any smooth curves $\gamma_1(r)$ and $\gamma_2(s)$
of $\mathbb{H}^2$ with $r,s$ being their arc length parameters, a hypersurface in
$\mathbb{H}^2\times\mathbb{H}^2$ can be defined by the immersion
$\Psi:\Omega\subset\mathbb{R}^3\rightarrow\mathbb{H}^2\times\mathbb{H}^2$:
$(t,r,s)\rightarrow(p(t,r),q(t,s))$
%
\begin{equation}\label{eqn:3.1ss1}
\begin{aligned}
&p(t,r)=\cosh(f_1(t))\gamma_1(r)+\sinh(f_1(t))N_1(r),\\
&q(t,s)=\cosh(f_2(t))\gamma_2(s)+\sinh(f_2(t))N_2(s),
\end{aligned}	
\end{equation}
%
where $N_1(r)$ and $N_2(s)$ are unit normal vector fields of $\gamma_1(r)$ and
$\gamma_2(s)$ in $\mathbb{H}^2$, respectively.
Let $k_1(r)$ and $k_2(s)$ be the curvatures of
curves $\gamma_1(r)$ and $\gamma_2(s)$ in $\mathbb{H}^2$, respectively.
%We call such hypersurface constructed by \eqref{eqn:3.1ss1} as $M_{\kappa,\tilde{\kappa}}^c$.

Since $\{\gamma_1(r), \frac{d\gamma_1(r)}{dr}, N_1(r)\}$ and $\{\gamma_2(s),
\frac{d\gamma_2(s)}{ds}, N_2(s)\}$ are two orthonormal frames
of $\mathbb{R}_1^{3}$ along curves $\gamma_1(r)$ and $\gamma_2(s)$ respectively, it follows that
\begin{equation}\label{eqn:3.2ss1}
\begin{aligned}
&\frac{d^2\gamma_1(r)}{dr^2}=\gamma_1(r)+k_1(r)N_1(r),\ \
\frac{dN_1(r)}{dr}=-k_1(r)\frac{d\gamma_1(r)}{dr},\\
&\frac{d^2\gamma_2(s)}{ds^2}=\gamma_2(s)+k_2(s)N_2(s),\ \
\frac{dN_2(s)}{ds}=-k_2(s)\frac{d\gamma_2(s)}{ds}.
\end{aligned} 	
\end{equation}%
Now, we have an orthonormal frame field $\left\{E_{1}, E_{2}, E_{3}\right\}$
of hypersurface $\Psi(t,r,s)$ defined as

\begin{equation*}
\begin{aligned}
E_1=&\frac{1}{\cosh(f_1(t))-\sinh(f_1(t))k_1(r)}
\frac{\partial}{\partial r}(p(t,r),q(t,s))
=(\frac{d\gamma_1(r)}{dr},0),\\
E_2=&\frac{1}{\cosh(f_2(t))-\sinh(f_2(t))k_2(s)}
\frac{\partial }{\partial s}(p(t,r),q(t,s))
=(0,\frac{d\gamma_2(s)}{ds}),
\end{aligned}	
\end{equation*}
\begin{equation*}
\begin{aligned}
E_3=&\frac{1}{\sqrt{(f_1'(t))^2+(f_2'(t))^2}}\frac{\partial}{\partial t}(p(t,r),q(t,s))\\
=&\frac{1}{\sqrt{(f_1'(t))^2+(f_2'(t))^2}}(f_1'(t)\sinh(f_1(t))\gamma_1(r)
+f_1'(t)\cosh(f_1(t))N_1(r),\\
&\qquad\qquad  f_2'(t)\sinh(f_2(t))\gamma_2(s)
+f_2'(t)\cosh(f_2(t))N_2(s)).\\
\end{aligned}	
\end{equation*}
The unit normal vector field $N$ of $\Psi(t,r,s)$ is given by
\begin{equation}\label{eqn:3.3ss1}
\begin{aligned}
N&=\frac{1}{\sqrt{(f_1'(t))^2+(f_2'(t))^2}}(f_2'(t)\sinh(f_1(t))\gamma_1(r)
+f_2'(t)\cosh(f_1(t))N_1(r),\\
&\qquad \qquad -f_1'(t)\sinh(f_2(t))\gamma_2(s)
-f_1'(t)\cosh(f_2(t))N_2(s)).
\end{aligned}
\end{equation}
It follows that $\Psi(t,r,s)$ has $C=g(PN,N)=\frac{(f_2'(t))^2-(f_1'(t))^2}{(f_1'(t))^2+(f_2'(t))^2}$.
By direct calculations, we get
\begin{equation}\label{eqn:3.4ss1}
\begin{aligned}
&AE_1=-\frac{f_2'(t)}{\sqrt{(f_1'(t))^2+(f_2'(t))^2}}\frac{\sinh(f_1(t))-k_1(r)\cosh(f_1(t))}
{\cosh(f_1(t))-k_1(r)\sinh(f_1(t))}E_1,\\
&AE_2=\frac{f_1'(t)}{\sqrt{(f_1'(t))^2+(f_2'(t))^2}}\frac{\sinh(f_2(t))-k_2(s)\cosh(f_2(t))}
{\cosh(f_2(t))-k_2(s)\sinh(f_2(t))}E_2,\\
&AE_3=-\frac{f_1'(t)f_2''(t)-f_1''(t)f_2'(t)}{((f_1'(t))^2+(f_2'(t))^2)^{\frac{3}{2}}}E_3.
\end{aligned}
\end{equation}
%That is to say:
%\begin{equation}\label{eqn:3.5ss1}
%\begin{aligned}
%&b_1=-\frac{f_2'(t)}{\sqrt{(f_1'(t))^2+(f_2'(t))^2}}\frac{\sinh(f_1(t))-\kappa(r)\cosh(f_1(t))}
%{\cosh(f_1(t))-\kappa(r)\sinh(f_1(t))},\ b_2=b_3=b_5=0,\\
%&b_4=\frac{f_1'(t)}{\sqrt{(f_1'(t))^2+(f_2'(t))^2}}\frac{\sinh(f_2(t))-\tilde{\kappa}(s)\cosh(f_2(t))}
%{\cosh(f_2(t))-\tilde{\kappa}(s)\sinh(f_2(t))},\ %b_6=-\frac{f_1'(t)f_2''(t)-f_1''(t)f_2'(t)}{((f_1'(t))^2+(f_2'(t))^2)^{\frac{3}{2}}}.
%\end{aligned}
%\end{equation}
\end{example}

%\begin{example}\label{exam:4.1ss1ss11ab1}
%For some constant $a$, when we take
%$$
%\begin{aligned}
%&f_1(t)=\ln\Big(\sqrt{2}\cosh(\frac{t}{\sqrt{2}})+\sqrt{\cosh(\sqrt{2}t)}\Big),\ \
%f_2(t)={\rm arcsinh}\Big(-\sqrt{2}\sinh(\frac{t}{\sqrt{2}})\Big)-a,\ \ t<0,\\
%&\gamma(r)=(\cosh(r),\sinh(r),0), \ \ N(r)=(0,0,1),\\
%&\tilde{\gamma}(s)=(\cosh(a),\sinh(a)\cos(\frac{s}{\sinh(a)}),\sinh(a)\sin(\frac{s}{\sinh(a)})), \ \ %{N}(s)=(\sinh(a),\cosh(a)\cos(\frac{s}{\sinh(a)}),\cosh(a)\sin(\frac{s}{\sinh(a)})),\\
%\end{aligned}	
%$$
%in \eqref{eqn:3.1ss1}, then we have the immersion:
%$\Phi_1:\Omega\subset\mathbb{R}^3\rightarrow\mathbb{H}^2\times\mathbb{H}^2$:
%$(t,r,s)\rightarrow(p(t,r),q(t,s))$, where $t<0$ and
%
%\begin{equation}\label{eqn:jr}
%\begin{aligned}
%&p(t,r)=\sqrt{\cosh(\sqrt{2}t)+1}(\cosh(r),\sinh(r),0)+\sqrt{\cosh(\sqrt{2}t)}(0,0,1),\\
%&q(t,s)=\sqrt{\cosh(\sqrt{2}t)}(1,0,0)+\sqrt{\cosh(\sqrt{2}t)-1}(0,\cos(s),\sin(s)),
%&p(t,r)=\Big(\sqrt{\cosh(\sqrt{2}t)+1}\cosh(r),\sqrt{\cosh(\sqrt{2}t)+1}\sinh(r),\sqrt{\cosh(\sqrt{2}t)}\Big),\\
%&q(t,s)=\Big(\sqrt{\cosh(\sqrt{2}t)},\sqrt{\cosh(\sqrt{2}t)-1}\cos(\frac{s}{\sinh(a)}),\sqrt{\cosh(\sqrt{2}t)-1}\sin(\frac{s}{\sinh(a)})\Big),
%\end{aligned}	
%\end{equation}
%It can be checked that such hypersurface constructed by \eqref{eqn:jr} has constant sectional curvature %$-\frac{1}{2}$, and
%$C={\rm sech}(\sqrt{2}t)$.
%\end{example}


\begin{remark}\label{rem:3.5fg1}
Hypersurfaces constructed by \eqref{eqn:3.1ss1} are curvature-adapted
hypersurface of $\mathbb{H}^2\times \mathbb{H}^2$ (see Proposition 3.1 of \cite{HLYZ}).
When taking $k_1(r)=k_2(s)=1$,
hypersurfaces constructed by \eqref{eqn:3.1ss1} are minimal if and only if the functions $f_1(t)$ and
$f_2(t)$ satisfy
$$
(f_2'(t)-f_1'(t))\big((f_1'(t))^2+(f_2'(t))^2\big)=f_1'(t)f_2''(t)-f_1''(t)f_2'(t).
$$
%Since $f_1(t)f_2(t)\neq0$, by the implicit function theorem we can assume
%$f_1(t)=t$. Then, above equation becomes $f_2''(t)=(f_2'(t)-1)\big((f_2'(t))^2+1\big)$.
When taking $k_1(r)=-k_2(s)=1$,
hypersurfaces constructed by \eqref{eqn:3.1ss1} are minimal if and only if the functions $f_1(t)$ and
$f_2(t)$ satisfy
$$
(f_2'(t)+f_1'(t))\big((f_1'(t))^2+(f_2'(t))^2\big)=f_1'(t)f_2''(t)-f_1''(t)f_2'(t).
$$
%Since $f_1(t)f_2(t)\neq0$, by the implicit function theorem we can assume
%$f_1(t)=t$. Then, above equation becomes $f_2''(t)=(f_2'(t)+1)\big((f_2'(t))^2+1\big)$.


\end{remark}

Now, we choose specific $f_1(t),f_2(t)$, $\gamma_1(r),N_1(r)$ and $\gamma_2(s),N_2(s)$ in \eqref{eqn:3.1ss1}, so that we can get a hypersurface of $\mathbb{H}^2\times \mathbb{H}^2$ with constant sectional curvature $-\frac{1}{2}$.

\begin{example}\label{exam:4.1ss1ss11ab1}
In \eqref{eqn:3.1ss1}, for some constant $0\neq a\in \mathbb{R}$, we take
$$
\begin{aligned}
&f_1(t)=\ln\Big(\sqrt{2}\cosh(\tfrac{t}{\sqrt{2}})+\sqrt{\cosh(\sqrt{2}t)}\Big),\ \
f_2(t)={\rm arcsinh}\Big(-\sqrt{2}\sinh(\tfrac{t}{\sqrt{2}})\Big)-a,\ \ t<0,\\
&\gamma_1(r)=(\cosh(r),\sinh(r),0), \ \
\gamma_2(s)=(\cosh(a),\sinh(a)\cos(s),\sinh(a)\sin(s)),\\
&N_1(r)=(0,0,1),\ \ N_2(s)=(\sinh(a),\cosh(a)\cos(s),\cosh(a)\sin(s)).
\end{aligned}	
$$
Then we have the immersion:
$\Phi_1:\Omega\subset\mathbb{R}^3\rightarrow\mathbb{H}^2\times\mathbb{H}^2$:
$(t,r,s)\rightarrow(p(t,r),q(t,s))$, where %$t<0$ and
%
\begin{equation}\label{eqn:jr}
\begin{aligned}
%&p(t,r)=\sqrt{\cosh(\sqrt{2}t)+1}(\cosh(r),\sinh(r),0)+\sqrt{\cosh(\sqrt{2}t)}(0,0,1),\\
%&q(t,s)=\sqrt{\cosh(\sqrt{2}t)}(1,0,0)+\sqrt{\cosh(\sqrt{2}t)-1}(0,\cos(s),\sin(s)),
&p(t,r)=\Big(\sqrt{\cosh(\sqrt{2}t)+1}\cosh(r),\sqrt{\cosh(\sqrt{2}t)+1}\sinh(r),\sqrt{\cosh(\sqrt{2}t)}\Big),\\
&q(t,s)=\Big(\sqrt{\cosh(\sqrt{2}t)},\sqrt{\cosh(\sqrt{2}t)-1}\cos(s),\sqrt{\cosh(\sqrt{2}t)-1}\sin(s)\Big).
\end{aligned}	
\end{equation}
By \eqref{eqn:2.2}, \eqref{eqn:3.3ss1} and \eqref{eqn:3.4ss1}, one can be checked that such hypersurface constructed by \eqref{eqn:jr} has constant sectional curvature $-\frac{1}{2}$, and
$C={\rm sech}(\sqrt{2}t)$.
\end{example}








%\begin{lemma}\label{lemma:3.5fg1}
%Let $M$ be an open part of hypersurface which constructed by
%\eqref{eqn:3.1ss1}.
%
%\begin{enumerate}
%\item[(1)]
%$M$ has constant $C$ if and only if $f_1(t)$ and $f_2(t)$ are two linear functions on $t$.

%\item[(2)]
%If $f_1(t)$ and $f_2(t)$ are two linear functions on $t$, then
%$M$ has constant principal curvatures if and only if $\kappa(r)=\tilde{\kappa}(s)=1$ or
%$\kappa(r)=-\tilde{\kappa}(s)=1$.

%\end{enumerate}
%\end{lemma}

As another special case of Example \ref{exam:4.1} by choosing $f_1(t)=\sqrt{c}t$ and $f_2(t)=\sqrt{1-c}t$
for some constant $0<c<1$, we have the following example.
\begin{example}\label{exam:3.4fg1}
For any constant $0<c<1$, and any smooth curves $\gamma_1(r)$ and $\gamma_2(s)$
of $\mathbb{H}^2$ with $r,s$ being their arc length parameters, we define a hypersurface $M_{k_1,k_2}^c$ in
$\mathbb{H}^2\times\mathbb{H}^2$ by the following immersion
$\Phi:\Omega\subset\mathbb{R}^3\rightarrow\mathbb{H}^2\times\mathbb{H}^2$:
$(t,r,s)\rightarrow(p(t,r),q(t,s))$,
%
\begin{equation}\label{eqn:3.1fg12}
\begin{aligned}
&p(t,r)=\cosh(\sqrt{c}t)\gamma_1(r)+\sinh(\sqrt{c}t)N_1(r),\\
&q(t,s)=\cosh(\sqrt{1-c}t)\gamma_2(s)+\sinh(\sqrt{1-c}t)N_2(s),
\end{aligned}	
\end{equation}
%
where $N_1(r)$ and $N_2(s)$ are unit normal vector fields of $\gamma_1(r)$ and
$\gamma_2(s)$ in $\mathbb{H}^2$, respectively.
Let $k_1(r)$ and $k_2(s)$ be the curvatures of
curves $\gamma_1(r)$ and $\gamma_2(s)$ in $\mathbb{H}^2$, respectively.
%We call such hypersurface constructed by \eqref{eqn:3.1fg12} as $M_{\kappa,\tilde{\kappa}}^c$.

It follows from \eqref{eqn:3.4ss1} that $M_{k_1,k_2}^c$ has constant
$C=g(PN,N)=1-2c$,
and its principal curvatures are
$$
\begin{aligned}
0,\ \ -\sqrt{1-c}\frac{\sinh(\sqrt{c}t)-\cosh(\sqrt{c}t)k_1(r)}
{\cosh(\sqrt{c}t)-\sinh(\sqrt{c}t)k_1(r)},\ \
\sqrt{c}\frac{\sinh(\sqrt{1-c}t)-\cosh(\sqrt{1-c}t)k_2(s)}
{\cosh(\sqrt{1-c}t)-\sinh(\sqrt{1-c}t)k_2(s)}.
\end{aligned}
$$

\begin{lemma}[cf. Lemma 3.5 and Remark 3.1 of \cite{GMY}]\label{lemma:3.5fg12}
Let $M$ be an open part of hypersurface $M_{k_1,k_2}^c$ which constructed by
\eqref{eqn:3.1fg12}, then
%
\begin{enumerate}
\item[(1)]
$M$ is minimal if and only if $c=\frac{1}{2}$, $k_1(r)$ and $k_2(s)$
are the same constant functions.

\item[(2)]
$M$ has constant sectional curvature if and only if $c=\frac{1}{2}$, $k_1(r)$ and
$k_2(s)$ are constant functions which satisfy $k_1(r) k_2(s)=1$.
Moreover, the sectional curvature of $M$ is $-\frac{1}{2}$.

\item[(3)]
$M$ has constant principal curvatures if and only if $k_1(r)=k_2(s)=1$ or
$k_1(r)=-k_2(s)=1$. We denote such hypersurfaces as $M_{1,1}^c$ and $M_{1,-1}^c$.
\end{enumerate}
\end{lemma}
\end{example}



















%\begin{example}\label{exam:4.1ss1}
%For any two functions $f_1(t),f_2(t)$ any smooth curve $\gamma(r)$
%of $\mathbb{H}^2$ with $r$ being its arc length parameter, a hypersurface in
%$\mathbb{H}^2\times\mathbb{H}^2$ can be defined by the immersion
%$\Psi_1:\Omega\subset\mathbb{R}^3\rightarrow\mathbb{H}^2\times\mathbb{H}^2$:
%$(t,r,s)\rightarrow(p(t,r),q(t,s))$
%
%\begin{equation}\label{eqn:3.2ss1ss1}
%\begin{aligned}
%&p(t,r)=\cosh(f_1(t))\gamma(r)+\sinh(f_1(t))N(r),\\
%&q(t,s)=\cosh(f_2(t))(1,0,0)+\sinh(f_2(t))(0,\cos(s),\sin(s)),
%\end{aligned}	
%\end{equation}
%
%where $N(r)$ is an unit normal vector field of $\gamma(r)\hookrightarrow\mathbb{H}^2$.
%Let $\kappa(r)$ be the curvature of curve $\gamma(r)$ in $\mathbb{H}^2$.

%Now, we have an orthonormal frame field $\left\{E_{1}, E_{2}, E_{3}\right\}$
%of hypersurface $\Psi(t,r,s)$ defined as
%\begin{equation*}
%\begin{aligned}
%E_1=&\frac{1}{\cosh(f_1(t))-\sinh(f_1(t))\kappa(r)}
%\frac{\partial}{\partial r}(p(t,r),q(t,s))
%=(\frac{d\gamma(r)}{dr},0),\\
%E_2=&\frac{1}{\sinh(f_2(t))}
%\frac{\partial }{\partial s}(p(t,r),q(t,s))
%=(0,(0,-\sin(s),\cos(s))),\\
%E_3=&\frac{1}{\sqrt{(f_1')^2+(f_2')^2}}\frac{\partial}{\partial t}(p(t,r),q(t,s))\\
%=&\frac{1}{\sqrt{(f_1')^2+(f_2')^2}}(f_1'(t)\sinh(f_1(t))\gamma(r)
%+f_1'(t)\cosh(f_1(t))N(r),\\
%&\qquad\qquad  f_2'(t)\sinh(f_2(t))(1,0,0)
%+f_2'(t)\cosh(f_2(t))(0,\cos(s),\sin(s))),\\
%\end{aligned}	
%\end{equation*}
%The unit normal vector field $N$ of $\Psi(t,r,s)$ is given by
%\begin{equation}\label{eqn:3.3ss1ss1}
%\begin{aligned}
%N&=(N_1,N_2)=\frac{1}{\sqrt{(f_1')^2+(f_2')^2}}(f_2'(t)\sinh(f_1(t))\gamma(r)
%+f_2'(t)\cosh(f_1(t))N(r),\\
%&-f_1'(t)\sinh(f_2(t))(1,0,0)-f_1'(t)\cosh(f_2(t))(0,\cos(s),\sin(s))).
%\end{aligned}
%\end{equation}
%It follows that $\Psi(t,r,s)$ has $C=g(PN,N)=\frac{(f_2')^2-(f_1')^2}{(f_1')^2+(f_2')^2}$.
%By direct calculations, we get
%\begin{equation}\label{eqn:3.4ss1ss1}
%\begin{aligned}
%&AE_1=-\frac{f_2'}{\sqrt{(f_1')^2+(f_2')^2}}\frac{\sinh(f_1(t))-\kappa(r)\cosh(f_1(t))}
%{\cosh(f_1(t))-\kappa(r)\sinh(f_1(t))}E_1,\\
%&AE_2=\frac{f_1'}{\sqrt{(f_1')^2+(f_2')^2}}\coth(f_2(t))E_2,\\
%&AE_3=-\frac{f_1'f_2''-f_1''f_2'}{((f_1')^2+(f_2')^2)^{\frac{3}{2}}}E_3.
%\end{aligned}
%\end{equation}
%That is to say:
%\begin{equation}\label{eqn:3.5ss1ss1}
%\begin{aligned}
%&b_1=-\frac{f_2'}{\sqrt{(f_1')^2+(f_2')^2}}\frac{\sinh(f_1(t))-\kappa(r)\cosh(f_1(t))}
%{\cosh(f_1(t))-\kappa(r)\sinh(f_1(t))},\ b_2=b_3=b_5=0,\\
%&b_4=\frac{f_1'}{\sqrt{(f_1')^2+(f_2')^2}}\coth(f_2(t)),\ %b_6=-\frac{f_1'f_2''-f_1''f_2'}{((f_1')^2+(f_2')^2)^{\frac{3}{2}}}.
%\end{aligned}
%\end{equation}
%\end{example}


%\begin{example}\label{exam:4.1ss1ss11ab1}
%When we take
%$$
%\begin{aligned}
%&f_1(t)=\ln\Big(\sqrt{2}\cosh(\frac{t}{\sqrt{2}})+\sqrt{\cosh(\sqrt{2}t)}\Big),\ \
%f_2(t)={\rm arcsinh}\Big(-\sqrt{2}\sinh(\frac{t}{\sqrt{2}})\Big),\ \ t<0,\\
%&\gamma(r)=(\cosh(r),\sinh(r),0), \ \ N(r)=(0,0,1),
%\end{aligned}	
%$$
%in \eqref{eqn:3.2ss1ss1}, then we have the immersion:
%$\Phi_1:\Omega\subset\mathbb{R}^3\rightarrow\mathbb{H}^2\times\mathbb{H}^2$:
%$(t,r,s)\rightarrow(p(t,r),q(t,s))$, where $t<0$ and
%
%\begin{equation}\label{eqn:jr}
%\begin{aligned}
%&p(t,r)=\sqrt{\cosh(\sqrt{2}t)+1}(\cosh(r),\sinh(r),0)+\sqrt{\cosh(\sqrt{2}t)}(0,0,1),\\
%&q(t,s)=\sqrt{\cosh(\sqrt{2}t)}(1,0,0)+\sqrt{\cosh(\sqrt{2}t)-1}(0,\cos(s),\sin(s)),
%&p(t,r)=\Big(\sqrt{\cosh(\sqrt{2}t)+1}\cosh(r),\sqrt{\cosh(\sqrt{2}t)+1}\sinh(r),\sqrt{\cosh(\sqrt{2}t)}\Big),\\
%&q(t,s)=\Big(\sqrt{\cosh(\sqrt{2}t)},\sqrt{\cosh(\sqrt{2}t)-1}\cos(s),\sqrt{\cosh(\sqrt{2}t)-1}\sin(s)\Big),
%\end{aligned}	
%\end{equation}
%It can be checked that such hypersurface constructed by \eqref{eqn:jr} has constant sectional curvature $-\frac{1}{2}$, and we call such hypersurface as $M_{-1/2}$.
%\end{example}



%%%%%%%%%%%%%%%%%%%%%%%%%%%%%%%%%%%%%%%%%%%%%%%%%%%%%%%%%%%%%%%%%
%\begin{example}\label{exam:4.1ss1ss11}
%For any two functions $f_1(t),f_2(t)$, a hypersurface in
%$\mathbb{H}^2\times\mathbb{H}^2$ can be defined by the immersion
%$\Psi_2:\Omega\subset\mathbb{R}^3\rightarrow\mathbb{H}^2\times\mathbb{H}^2$:
%$(t,r,s)\rightarrow(p(t,r),q(t,s))$
%
%\begin{equation}\label{eqn:3.2ss1ss11}
%\begin{aligned}
%&p(t,r)=\cosh(f_1(t))(1,0,0)+\sinh(f_1(t))(0,\cos(r),\sin(r)),\\
%&q(t,s)=\cosh(f_2(t))(1,0,0)+\sinh(f_2(t))(0,\cos(s),\sin(s)).
%\end{aligned}	
%\end{equation}
%


%Now, we have an orthonormal frame field $\left\{E_{1}, E_{2}, E_{3}\right\}$
%of hypersurface $\Psi(t,r,s)$ defined as
%\begin{equation*}
%\begin{aligned}
%E_1=&\frac{1}{\sinh(f_1(t))}
%\frac{\partial}{\partial r}(p(t,r),q(t,s))
%=((0,-\sin(r),\cos(r)),\vec{0}),\\
%E_2=&\frac{1}{\sinh(f_2(t))}
%\frac{\partial }{\partial s}(p(t,r),q(t,s))
%=(\vec{0},(0,-\sin(s),\cos(s))),\\
%E_3=&\frac{1}{\sqrt{(f_1')^2+(f_2')^2}}\frac{\partial}{\partial t}(p(t,r),q(t,s))\\
%=&\frac{1}{\sqrt{(f_1')^2+(f_2')^2}}(f_1'(t)\sinh(f_1(t))(1,0,0)
%+f_1'(t)\cosh(f_1(t))(0,\cos(r),\sin(r)),\\
%&\qquad\qquad  f_2'(t)\sinh(f_2(t))(1,0,0)
%+f_2'(t)\cosh(f_2(t))(0,\cos(s),\sin(s))),\\
%\end{aligned}	
%\end{equation*}
%The unit normal vector field $N$ of $\Psi(t,r,s)$ is given by
%\begin{equation}\label{eqn:3.3ss1ss11}
%\begin{aligned}
%N&=(N_1,N_2)=\frac{1}{\sqrt{(f_1')^2+(f_2')^2}}(f_2'(t)\sinh(f_1(t))(1,0,0)
%+f_2'(t)\cosh(f_1(t))(0,\cos(r),\sin(r)),\\
%&-f_1'(t)\sinh(f_2(t))(1,0,0)-f_1'(t)\cosh(f_2(t))(0,\cos(s),\sin(s))).
%\end{aligned}
%\end{equation}
%It follows that $\Psi(t,r,s)$ has $C=g(PN,N)=\frac{(f_2')^2-(f_1')^2}{(f_1')^2+(f_2')^2}$.
%By direct calculations, we get
%\begin{equation}\label{eqn:3.4ss1ss11}
%\begin{aligned}
%&AE_1=-\frac{f_2'}{\sqrt{(f_1')^2+(f_2')^2}}\coth(f_1(t))E_1,\\
%&AE_2=\frac{f_1'}{\sqrt{(f_1')^2+(f_2')^2}}\coth(f_2(t))E_2,\\
%&AE_3=-\frac{f_1'f_2''-f_1''f_2'}{((f_1')^2+(f_2')^2)^{\frac{3}{2}}}E_3.
%\end{aligned}
%\end{equation}
%That is to say:
%\begin{equation}\label{eqn:3.5ss1ss11}
%\begin{aligned}
%&b_1=-\frac{f_2'}{\sqrt{(f_1')^2+(f_2')^2}}\coth(f_1(t)),\ b_2=b_3=b_5=0,\\
%&b_4=\frac{f_1'}{\sqrt{(f_1')^2+(f_2')^2}}\coth(f_2(t)),\ %b_6=-\frac{f_1'f_2''-f_1''f_2'}{((f_1')^2+(f_2')^2)^{\frac{3}{2}}}.
%\end{aligned}
%\end{equation}
%\end{example}

%\begin{remark}\label{rem:3.5fg12}
%(1) When taking $f_1(t)=f_2(t)$,
%hypersurface constructed by \eqref{eqn:3.2ss1ss11} is minimal.
%\end{remark}


%==================================================
\section{Minimal hypersurfaces with $C=0$}\label{sect:4}
%
In this section, we give the classification of the minimal hypersurfaces with $C=0$, which
plays an important role in the proof of Theorem \ref{thm:1.2}. 

\begin{lemma}\label{lemma:2.1aaa1}
Let $M$ be a minimal hypersurface of $\mathbb{H}^2\times \mathbb{H}^2$ with $C=0$.
Then, under the frame $\{E_1,E_2,E_3\}$ defined by \eqref{eqn:2.5}, we have two cases
for $b_1,b_2$ and $b_4$ on $M$ as follows:

\noindent{\bf Case 1}:
$$
b_1=-b_4=-\tfrac{1}{\sqrt{2}}\tanh(\tfrac{t-k}{\sqrt{2}}),\ \ b_2=0,
$$

\noindent{\bf Case 2}:
\begin{equation}\label{eq:j}
\begin{aligned}
b_1&=-b_4=-\tfrac{1}{\sqrt{2}}\frac{\sinh(\frac{t-l}{\sqrt{2}})\cosh(\frac{t-l}{\sqrt{2}})}
{\sinh^2(\frac{t-l}{\sqrt{2}})+\cos^2(\frac{h}{\sqrt{2}})},\\
b_2&=-\tfrac{1}{\sqrt{2}}\frac{\sin(\frac{h}{\sqrt{2}})\cos(\frac{h}{\sqrt{2}})}
{\sinh^2(\frac{t-l}{\sqrt{2}})+\cos^2(\frac{h}{\sqrt{2}})},
\end{aligned}	
\end{equation}
where function $t$ satisfies $E_3t=1$ and functions $h, k, l:M\rightarrow \mathbb{R}$
satisfy $E_3h=E_3k=E_3l=0$.

\end{lemma}
%
\begin{proof}
By $H=\frac{1}{3}(b_1+b_4)=0$ and $C=0$, equations \eqref{eqn:2.26} and \eqref{eqn:2.27} become
\begin{equation*}
\begin{split}
E_3b_1&=-\tfrac{1}{2}+b_1^2-b_2^2,\\
E_3b_2&=2b_1b_2.
\end{split}	
\end{equation*}


When $b_2$ is constant on $M$, we have $b_1b_2=0$. Furthermore,
if $b_1=0$, then we get a contradiction from the first equation.
If $b_2=0$, then $b_1$ satisfies $E_3b_1=-\tfrac{1}{2}+b_1^2$. By directly integrating
this equation along the integral curves of $E_3$, we can obtain the {\bf Case 1}.

When $b_2$ is not constant on $M$, let $F=b_1+\mathbf{i}b_2$, we get
$$
E_3F=E_3(b_1+\mathbf{i}b_2)=-\frac{1}{2}+(b_1+\mathbf{i}b_2)^2=-\frac{1}{2}+F^2.
$$
This equation can be directly integrated along the integral curves of $E_3$
to obtain $F=-\frac{1}{\sqrt{2}}\tanh(\frac{t+z_0}{\sqrt{2}})$,
where $z_0=-l+\mathbf{i}h$, $h$ and $l$ are functions on $M$ which satisfy $E_3h=E_3l=0$.
By $F=b_1+\mathbf{i}b_2$, it follows that $b_1$ and $b_2$ are its real and imaginary parts,
which corresponds to the {\bf Case 2}.
\end{proof}



In the following, we deal with the {\bf Case 2} of Lemma \ref{lemma:2.1aaa1} 
by finding appropriate local coordinates. 
%The proof follows the similar idea of \cite{K-V} to find a appropriate local coordinates.
We assume that
\begin{equation}\label{UVT-H2}
\left(
  \begin{array}{c}
    U \\
    V \\
    T \\
  \end{array}
\right)
=
\left(
  \begin{array}{ccc}
    a_1 & a_2 & a_1d_1+a_2d_2 \\
    -a_2 & a_1 & -a_2d_1+a_1d_2 \\
    0 & 0 & 1 \\
  \end{array}
\right)
\left(
  \begin{array}{c}
    E_1 \\
    E_2 \\
    E_3 \\
  \end{array}
\right),
\end{equation}
where $a_1,a_2,d_1,d_2$ are smooth functions of $M$.

\begin{lemma}\label{lemma:2.1a1}
$\{U,V,T\}$ is a local coordinate frame if and only if
the following system of equations is satisfied:
%
\begin{equation}\label{eq:2.ab1}
E_3a_1=-a_1b_1-a_2b_2,\ \ E_3a_2=a_1b_2-a_2b_1,
\end{equation}
%
\begin{equation}\label{eq:2.ab2}
E_1a_2+E_2a_1=0,\ \ E_1a_1-E_2a_2=0,
\end{equation}
%
\begin{equation}\label{eq:2.ab3}
E_3d_1=b_1d_1-b_2d_2,\ \ E_3d_2=d_1b_2+d_2b_1,
\end{equation}
%
\begin{equation}\label{eq:2.ab4}
E_1d_2-E_2d_1=2b_2.
\end{equation}

\end{lemma}
%
\begin{proof}
We just need to show that our system of equations is equivalent to the vanishing of all the commutators $[T,U]$, $[T,V]$ and $[U,V]$.
Using $[E_i,E_j]=\nabla_{E_i} {E_j}-\nabla_{E_j} {E_i}$, $C=0$
and \eqref{eqn:2.25}, we get
$$
[E_1,E_2]=-2b_2E_3,\ [E_1,E_3]=-b_1E_1+b_2E_2,\ [E_2,E_3]=-b_2E_1+b_4E_2=-b_2E_1-b_1E_2.
$$
Due that \eqref{UVT-H2} is the same as (3.2) of \cite{LVWY}, then the proof of Lemma \ref{lemma:2.1a1} is exactly the same as the proof of Lemma 3.2 of \cite{LVWY}. 
\end{proof}

In the following, we will give a special solution to the system of equations \eqref{eq:2.ab1}--\eqref{eq:2.ab4}, which provides us adapted coordinates. 

%Any solution to the system of equations \eqref{eq:2.ab1}--\eqref{eq:2.ab4}
%provides us adapted coordinates. Now, we give a special solution as follows.

\begin{lemma}\label{lemma:2.1a2}
Assume that $b_1,b_2$ and $b_4$ are given by \eqref{eq:j} on $M$,
and let $f$ be any function satisfying $E_3f=-1$. Then the functions $a_1,a_2,d_1,d_2$ defined by
\begin{equation}\label{eq:2.ab7}
\begin{aligned}
(a_1+\mathbf{i}a_2)^2=\frac{1}{1-2(b_1-\mathbf{i}b_2)^2}, \ \ d_1=E_1f,\ d_2=E_2f,
\end{aligned}
\end{equation}
are a solution to the system of differential equations \eqref{eq:2.ab1}--\eqref{eq:2.ab4}.
\end{lemma}
%
\begin{proof}
Observe that the system of equations splits into two independent
sub-systems, \eqref{eq:2.ab1}--\eqref{eq:2.ab2} and \eqref{eq:2.ab3}--\eqref{eq:2.ab4}.

Assume now that $a_1,a_2$ are functions about $b_1,b_2$. Using the equations
\eqref{eqn:2.26}--\eqref{eqn:2.30} and $C=0$, we re-express the first sub-system with respect to
the $b_1,b_2$ to obtain that
%
\begin{equation}\label{eq:2.ab8}
(b_1^2-b_2^2-\frac{1}{2})\frac{\partial a_1}{\partial b_1}+2b_1b_2\frac{\partial a_1}{\partial b_2}
=-a_1b_1-a_2b_2,
\end{equation}
%
\begin{equation}\label{eq:2.ab9}
(b_1^2-b_2^2-\frac{1}{2})\frac{\partial a_2}{\partial b_1}+2b_1b_2\frac{\partial a_2}{\partial b_2}
=a_1b_2-a_2b_1,
\end{equation}
%
\begin{equation}\label{eq:2.ab10}
\frac{\partial a_2}{\partial b_1}E_1b_1+\frac{\partial a_2}{\partial b_2}E_1b_2
+\frac{\partial a_1}{\partial b_1}E_2b_1+\frac{\partial a_1}{\partial b_2}E_2b_2=0,
\end{equation}
%
\begin{equation}\label{eq:2.ab11}
\frac{\partial a_1}{\partial b_1}E_1b_1+\frac{\partial a_1}{\partial b_2}E_1b_2
-\frac{\partial a_2}{\partial b_1}E_2b_1-\frac{\partial a_2}{\partial b_2}E_2b_2=0.
\end{equation}
Equations \eqref{eq:2.ab10} and \eqref{eq:2.ab11} are satisfied if $(a_1,a_2)$ satisfy the Cauchy-Riemann equations with respect to $b_1,b_2$, i.e.,
$$
\frac{\partial a_2}{\partial b_1}=\frac{\partial a_1}{\partial b_2},\ \
\frac{\partial a_2}{\partial b_2}=-\frac{\partial a_1}{\partial b_1}.
$$

Therefore, we assume that $a_1(z)+\mathbf{i}a_2(z)$ is a holomorphic function with respect
to $z=b_1-\mathbf{i}b_2$. Then we have
$$
(z^2-\tfrac{1}{2})\frac{\partial}{\partial z}(a_1+\mathbf{i}a_2)=
(b_1^2-b_2^2-\tfrac{1}{2})\frac{\partial a_1}{\partial b_1}+2b_1b_2\frac{\partial a_1}{\partial b_2}
+\mathbf{i}\Big((b_1^2-b_2^2-\tfrac{1}{2})\frac{\partial a_2}{\partial b_1}+2b_1b_2\frac{\partial a_2}{\partial b_2}\Big),
$$
and
$$
-z(a_1+\mathbf{i}a_2)=(-a_1b_1-a_2b_2)+\mathbf{i}(-b_1a_2+b_2a_1).
$$
Then, \eqref{eq:2.ab8} and \eqref{eq:2.ab9} are equivalent to
$$
(z^2-\tfrac{1}{2})\frac{\partial}{\partial z}(a_1+\mathbf{i}a_2)=-z(a_1+\mathbf{i}a_2).
$$
Thus, we get a special solution to the equation system \eqref{eq:2.ab8}--\eqref{eq:2.ab11}: $(a_1+\mathbf{i}a_2)^2=\frac{1}{1-2(b_1-\mathbf{i}b_2)^2}$.
%The second sub-system can be re-written as the integrability conditions for a
%function $f$ with $E_3f=-1$. In fact, if we set $d_1=E_1f, d_2=E_2f, d_3=E_3f=-1$, we get
%$$
%\begin{aligned}
%&[E_1,E_3]f=E_1d_3-E_3d_1=-b_1d_1+b_2d_2,\\
%&[E_2,E_3]f=E_2d_3-E_3d_2=-b_2d_1-b_1d_2,\\
%&[E_1,E_2]f=E_1d_2-E_2d_1=2b_2.
%\end{aligned}
%$$
%Hence $\{d_1, d_2, -1\}$ satisfy the integrability conditions for arbitrary functions $d_1, d_2$.

Let $f$ be any function satisfying that $E_3f=-1$. We denote $d_1=E_1f, d_2=E_2f$, by using \eqref{eqn:2.25}, we obtain the following integrability conditions for $f$.
$$
\begin{aligned}
	&0=[E_1,E_3]f+b_1E_1(f)-b_2E_2(f)=-E_3d_1+b_1d_1-b_2d_2,\\
	&0=[E_2,E_3]f+b_2E_1(f)+b_1E_2(f)=-E_3d_2+b_2d_1+b_1d_2,\\
	&0=[E_1,E_2]f+2b_2E_3(f)=E_1d_2-E_2d_1-2b_2,
\end{aligned}
$$
which implies that  $d_1=E_1f, d_2=E_2f$ is a solution to the
second subsystem \eqref{eq:2.ab3}--\eqref{eq:2.ab4}.
\end{proof}

We denote the coordinates obtained from the solution given in Lemma \ref{lemma:2.1a2} by $(u,v,t)$, i.e.
$\partial_u=U,\ \partial_v=V,\ \partial_t=T=E_3$.
We will further make a special choice for $f$ to simplify \eqref{eq:j}.

\begin{lemma}\label{lemma:2.1a3}
There exists a function $f$ with $E_3f=-1$ and a constant $\tilde{c}$ such that
replacing $t$ by $t-\tilde{c}$ one has $l=0$. Moreover, $f$ is given by
$$
f=-\tfrac{1}{\sqrt{2}}{\rm arcsinh}\Big(\tfrac{-2\sqrt{2}b_1}
{\sqrt{(2b_1^2-2b_2^2-1)^2+16b_1^2b_2^2}}\Big).
$$
The variable $t$ is given by $t=-f+$const.

\end{lemma}
%
\begin{proof}
A short calculation implies that our choice of $f$ satisfies $E_3f=-1$.
By \eqref{UVT-H2}, we get
%
\begin{equation}\label{eq:2.ab12}
E_1=\frac{a_1}{a_1^2+a_2^2}\partial_u-\frac{a_2}{a_1^2+a_2^2}\partial_v-d_1\partial_t,
\end{equation}
%
\begin{equation}\label{eq:2.ab13}
E_2=\frac{a_2}{a_1^2+a_2^2}\partial_u+\frac{a_1}{a_1^2+a_2^2}\partial_v-d_2\partial_t,
\end{equation}
%
\begin{equation}\label{eq:2.ab14}
E_3=\partial_t.
\end{equation}
From $E_1f=d_1, E_2f=d_2, E_3f=-1$, we have
$$
\frac{\partial}{\partial u}f=\frac{\partial}{\partial v}f=0.
$$
We insert \eqref{eq:j} into the definition of $f$ and obtain $\sinh(-\sqrt{2}f)=\sinh(\sqrt{2}(t-l))$,
hence $\frac{\partial}{\partial u}l=\frac{\partial}{\partial v}l=0$.
Recall that $E_3l=0$ (see Lemma \ref{lemma:2.1aaa1}), we obtain that $l$ must be a constant.
\end{proof}



\begin{theorem}\label{thm:3.1}
Let $M$ be a minimal hypersurface of $\mathbb{H}^2\times \mathbb{H}^2$ with product angle function $C=0$.
Assume that with respect to frame field $\{E_1,E_2,E_3\}$ (see \eqref{eqn:2.5}), the function $b_2$ defined in \eqref{eqn:2.7} is not constant on $M$.
Then $M$ admits local coordinates $(u,v,t)$ and there is a non-zero function $h:(u,v)\mapsto \mathbb{R}$
such that the functions $b_1,b_2,b_4$ defined in \eqref{eqn:2.7} are given by
\eqref{eq:2.ab15}, where $h$ satisfies the following ``sin-Gordon equation''
%
\begin{equation}\label{eq:2.ab16}
(\frac{\partial^2}{\partial u^2}+\frac{\partial^2}{\partial v^2})h
=\tfrac{1}{\sqrt{2}}\sin(\sqrt{2}h).
\end{equation}

Conversely, let $\mathcal{D}$ be an open subset in $\mathbb{R}^2$ and  $h:\mathcal{D}\to\mathbb{R}, (u,v)\mapsto h(u,v)\in\mathbb{R}$
be a non-zero function which satisfies the differential equation \eqref{eq:2.ab16}.
Let $\Omega_0=\{(u,v,t)\in \mathcal{D}\times \mathbb{R}|\sqrt{2}h(u,v)=(2p+1)\pi,t=0, p\in\mathbb{Z}\}$, $\Omega=\{(u,v,t)\in \mathcal{D}\times \mathbb{R}\}-\Omega_0\subseteq \mathbb{R}^3$, we define a metric $g$ on $\Omega$ by \eqref{eq:2.ab17} and a $(1,1)$-tensor field $A$ on $T\Omega$ by
\eqref{eqn:A}. Let $\nu$ be a vector bundle over  $\Omega$ of rank $1$ with metric $\tilde{g}$, $\nabla^{\perp}$ a connection on $\nu$ compatible with the metric $\tilde{g}$, and $N$ a unit vector field in $\nu$.   We define a $(1,1)$-tensor field $\tilde{P}$ on $T\Omega\otimes \nu$ by
\eqref{defP}.
Then the integrability conditions are satisfied on $\Omega$ and hence on any simply
connected subset $\Omega_1$ of $\Omega$, up to an isometry of  $\mathbb{H}^2\times \mathbb{H}^2$, there is a unique isometric immersion $\Psi$ from $(\Omega_1,~g)$ into $\mathbb{H}^2\times \mathbb{H}^2$ such that $A$ is the shape operator of  $\Omega_1$ into $\mathbb{H}^2\times \mathbb{H}^2$, $\nu$ is isometric to the normal bundle of $\Psi(\Omega_1)$ in
$\mathbb{H}^2\times \mathbb{H}^2$ by an isomorphism $\tilde{\Psi}:\nu\to T^{\perp}\Psi(\Omega_1)$,
and for any $Y\in T\Omega_1$, we have
\begin{equation}
	P(\Psi_{*}Y)=\Psi_{*}((\tilde{P}Y)^{T})+\tilde{\Psi}((\tilde{P}Y)^{\perp}),~
	P(\tilde{\Psi}N)=\Psi_*(\partial_t),
\end{equation}
where $P$ is the product structure of $\mathbb{H}^2\times \mathbb{H}^2$, $(\tilde{P}Y)^{T}$ and $(\tilde{P}Y)^{\perp}$ denote the projections of  $\tilde{P}Y$ onto $T\Omega_1$ and $\nu$, respectively.
Moreover, $\Psi(\Omega_1)$ is a minimal hypersurface with $C=0$.
\end{theorem}
%
\begin{proof}
Let $M$ be a minimal hypersurface of $\mathbb{H}^2\times \mathbb{H}^2$ with $C=0$.
By Lemmas \ref{lemma:2.1aaa1}--\ref{lemma:2.1a3}, $M$ admits local coordinates $(u,v,t)$ and
there is a function $h:(u,v)\mapsto \mathbb{R}$
such that $b_1,b_2,b_4$ are given by
%
\begin{equation}\label{eq:2.ab15}
b_1=-b_4=-\tfrac{1}{\sqrt{2}}\frac{\sinh(\frac{t}{\sqrt{2}})\cosh(\frac{t}{\sqrt{2}})}
{\sinh^2(\frac{t}{\sqrt{2}})+\cos^2(\frac{h}{\sqrt{2}})},\ \
b_2=-\tfrac{1}{\sqrt{2}}\frac{\sin(\frac{h}{\sqrt{2}})\cos(\frac{h}{\sqrt{2}})}
{\sinh^2(\frac{t}{\sqrt{2}})+\cos^2(\frac{h}{\sqrt{2}})},
\end{equation}
%By Lemma \ref{lemma:2.1a3} we can assume without loss of generality that the
%functions $b_1,b_2,b_4$ are given by \eqref{eq:2.ab15}.
Now, by Lemma \ref{lemma:2.1a2}, we can obtain that
%
\begin{equation}\label{eq:2.ab18}
a_1=\cosh(\tfrac{t}{\sqrt{2}})\cos(\tfrac{h}{\sqrt{2}}),\
a_2=-\sinh(\tfrac{t}{\sqrt{2}})\sin(\tfrac{h}{\sqrt{2}}).
\end{equation}

Using \eqref{eq:2.ab12}--\eqref{eq:2.ab14} and \eqref{eq:2.ab15},
then \eqref{eqn:2.29} and \eqref{eqn:2.30} are equivalent to
$$
d_1=\tfrac{2\big(\cosh(\frac{t}{\sqrt{2}})\cos(\frac{h}{\sqrt{2}})h_v
-\sinh(\frac{t}{\sqrt{2}})\sin(\frac{h}{\sqrt{2}})h_u\big)}{\cosh(\sqrt{2}t)+\cos(\sqrt{2}h)},\ \
d_2=-\tfrac{2\big(\sinh(\frac{t}{\sqrt{2}})\sin(\frac{h}{\sqrt{2}})h_v
+\cosh(\frac{t}{\sqrt{2}})\cos(\frac{h}{\sqrt{2}})h_u\big)}{\cosh(\sqrt{2}t)+\cos(\sqrt{2}h)}.
$$

Now, we can check the compatibility of functions $b_1,b_2,b_4$ with respect to the Lie bracket $[E_i,E_j]$,
$1\leq i<j\leq3$, and obtain that
$$
(\frac{\partial^2}{\partial u^2}+\frac{\partial^2}{\partial v^2})h
=\tfrac{1}{\sqrt{2}}\sin(\sqrt{2}h).
$$
In fact, the  compatibility conditions of  $b_1,b_2,b_4$  with respect to the Lie brackets $[E_1,E_3]$ and $[E_2,E_3]$ are automatically satisfied and the compatibility conditions of  $b_1,b_2,b_4$  with respect to the Lie bracket $[E_1,E_2]$  imply the above differential equation.


Conversely, let $\mathcal{D}$ be an open subset in $\mathbb{R}^2$ and  $h:\mathcal{D}\to\mathbb{R}, (u,v)\mapsto h(u,v)\in\mathbb{R}$
be a non-zero function which satisfies the differential equation \eqref{eq:2.ab16}.
We assume that the metric $g$ on $\Omega$ is defined by
%
\begin{equation}\label{eq:2.ab17}
\begin{aligned}
g&=\Big(\tfrac{1}{2}(\cosh(\sqrt{2}t)+\cos(\sqrt{2}h))+h_v^2\Big)du^2
+\Big(\tfrac{1}{2}(\cosh(\sqrt{2}t)+\cos(\sqrt{2}h))+h_u^2\Big)dv^2+dt^2\\
&\ \ -2h_uh_vdudv+2h_vdudt-2h_udvdt,
\end{aligned}
\end{equation}
the $(1,1)$-tensor field $A$ on $T\Omega$ is defined by
\begin{equation}\label{eqn:A}
\begin{aligned}
A\partial_u&=-\tfrac{\cos(\sqrt{2}h)\sinh(\sqrt{2}t)}{\sqrt{2}(\cosh(\sqrt{2}t)+\cos(\sqrt{2}h))}
\partial_u
-\tfrac{\cosh(\sqrt{2}t)\sin(\sqrt{2}h)}{\sqrt{2}(\cosh(\sqrt{2}t)+\cos(\sqrt{2}h))}\partial_v\\
&\ \ +\tfrac{\cos(\sqrt{2}h)\sinh(\sqrt{2}t)h_v-\cosh(\sqrt{2}t)\sin(\sqrt{2}h)h_u}
{\sqrt{2}(\cosh(\sqrt{2}t)+\cos(\sqrt{2}h))}\partial_t,\\
A\partial_v&=-\tfrac{\cosh(\sqrt{2}t)\sin(\sqrt{2}h)}{\sqrt{2}(\cosh(\sqrt{2}t)+\cos(\sqrt{2}h))}
\partial_u
+\tfrac{\cos(\sqrt{2}h)\sinh(\sqrt{2}t)}{\sqrt{2}(\cosh(\sqrt{2}t)+\cos(\sqrt{2}h))}\partial_v\\
&\ \ +\tfrac{\sin(\sqrt{2}h)\cosh(\sqrt{2}t)h_v+\sinh(\sqrt{2}t)\cos(\sqrt{2}h)h_u}
{\sqrt{2}(\cosh(\sqrt{2}t)+\cos(\sqrt{2}h))}\partial_t,\\
A\partial_t&=0,
\end{aligned}
\end{equation}
and the $(1,1)$-tensor field $\tilde{P}$ on $T\Omega\otimes \nu$ is defined by 
\begin{equation}\label{defP}
\begin{aligned}
	\tilde{P}\partial_u&=\tfrac{1+\cos(\sqrt{2}h)\cosh(\sqrt{2}t)}{\cos(\sqrt{2}h)+\cosh(\sqrt{2}t)}
	\partial_u
	+\tfrac{\sin(\sqrt{2}h)\sinh(\sqrt{2}t)}{\cos(\sqrt{2}h)+\cosh(\sqrt{2}t)}\partial_v\\
	&\ \ -
\tfrac{h_v+h_v\cos(\sqrt{2}h)\cosh(\sqrt{2}t)-h_u\sin(\sqrt{2}h)\sinh(\sqrt{2}t)}
{{\cos(\sqrt{2}h)+\cosh(\sqrt{2}t)}}\partial_t+h_vN,\\
	\tilde{P}\partial_v&=\tfrac{\sin(\sqrt{2}h)\sinh(\sqrt{2}t)}{\cos(\sqrt{2}h)+\cosh(\sqrt{2}t)}
	\partial_u-
	\tfrac{1+\cos(\sqrt{2}h)\cosh(\sqrt{2}t)}{\cos(\sqrt{2}h)+\cosh(\sqrt{2}t)}\partial_v\\
	&\ \  -\tfrac{h_u+h_u\cos(\sqrt{2}h)\cosh(\sqrt{2}t)+h_v\sin(\sqrt{2}h)\sinh(\sqrt{2}t)}
{{\cos(\sqrt{2}h)+\cosh(\sqrt{2}t)}}\partial_t-h_uN,\\
	\tilde{P}\partial_t&=N,\\
	\tilde{P}N&=\partial_t.\\	
\end{aligned}
\end{equation}
%assume that the function $h$ satisfies \eqref{eq:2.ab16}, the metric $g$ is defined according
%to \eqref{eq:2.ab17}, the $(1,1)$-tensor $A$ is defined by \eqref{eqn:A}, and we have a $(1,1)$-tensor $\tilde{P}$ defined by \eqref{defP},
It can be checked  directly that all the integrability conditions are satisfied.
Therefore, the conclusions can be obtained by applying  the existence and uniqueness theorems in  \cite{K,L-T-V} directly.
\end{proof}

%\begin{remark}\label{rem:3.1aa}
%If $M$ is a minimal hypersurface described by Theorem \ref{thm:3.1},
%then $M$ does not have constant scalar curvature. In fact, according to Corollary 1 (3)
%of \cite{Ur}, $M_0$ is the only hypersurface of $\mathbb{S}^2\times\mathbb{S}^2$
%which has the mean curvature, the scalar curvature and the function $C$ constants.
%But, under the frame $\{E_1,E_2,E_3\}$, $b_2=\frac{1}{\sqrt{2}}$ holds on $M_0$,
%which contradicts to the statement of Theorem \ref{thm:3.1}.
%\end{remark}

\begin{remark}\label{rem:3.1aaa1}
{\bf Case 1} of Lemma \ref{lemma:2.1aaa1}
corresponds to a trivial solution to the differential equation \eqref{eq:2.ab16}: $h\equiv0$,
and the corresponding example is the hypersurface
$M_{k_1,k_1}^{1/2}$ given in Example \ref{exam:3.4fg1} for some constant $k_1$,
see Theorem  \ref{thm:3.2aa} for more details.
\end{remark}

Before we deal with the {\bf Case 1} of Lemma \ref{lemma:2.1aaa1},
we shall prove the following proposition.

\begin{proposition}\label{prop:3.61}
Let $M$ be a hypersurface of $\mathbb{H}^2\times\mathbb{H}^2$ with constant $C\neq\pm1$.
Then, under the frame $\{E_1,E_2,E_3\}$ defined by \eqref{eqn:2.5}, $b_2$ is constant on $M$ if and only if either
%
\begin{enumerate}
\item[(1)]
$b_2\neq 0$ and $M$ is an open part of $M_\tau$ for some $\tau<-1$; or

\item[(2)]
$b_2=0$ and $M$ is an open part of $M_{k_1,k_2}^c$ given in Example \ref{exam:3.4fg1}. %where $C=1-2c$.

%locally given by map
%$\Phi: \Omega\subset\mathbb{R}^{3}\longrightarrow \mathbb{H}^2\times\mathbb{H}^2$:
%\ $(t,r,s)\longrightarrow(p,q)$
%
%\begin{equation}\label{eqn:6.1}
%\begin{aligned}
%&p(t,r)=\cosh(\sqrt{\tfrac{1-C}{2}}t)\gamma(r)+\sinh(\sqrt{\tfrac{1-C}{2}}t)N(r),\\
%&q(t,s)=\cosh(\sqrt{\tfrac{1+C}{2}}t)\tilde{\gamma}(s)+\sinh(\sqrt{\tfrac{1+C}{2}}t)\tilde{N}(s),
%\end{aligned}	
%\end{equation}
\end{enumerate}
%where $\gamma(r)$ and $\tilde{\gamma}(s)$ are any two curves in $\mathbb{H}^2$,
%$r$ and $s$ are their arc length parameters,
%$N(r)$ and $\tilde{N}(s)$ are unit normal vector fields of $\gamma(r)$ and $\tilde{\gamma}(s)$
%in $\mathbb{H}^2$, respectively.
\end{proposition}

\begin{proof}
To verify the ``only if" part, we assume that $b_2$ is constant on $M$.

If $b_2$ is a non-zero constant on $M$, then by \eqref{eqn:2.27}, it holds
$$
b_1\sqrt{\tfrac{1-C}{1+C}}-b_4\sqrt{\tfrac{1+C}{1-C}}=0,
$$
which implies that $b_4=\frac{1-C}{1+C}b_1$.
Substituting $b_4=\frac{1-C}{1+C}b_1$ into \eqref{eqn:2.26} and \eqref{eqn:2.28}, we have
$$
\begin{aligned}
E_3b_1&=\tfrac{1+C}{2}\sqrt{\tfrac{1+C}{1-C}}+b_2^2\sqrt{\tfrac{1+C}{1-C}}
-b_1^2\sqrt{\tfrac{1-C}{1+C}},\\
E_3b_1&=-\tfrac{\sqrt{1-C^2}}{2}+b_1^2\sqrt{\tfrac{1-C}{1+C}}-b_2^2\sqrt{\tfrac{1+C}{1-C}}.
\end{aligned}
$$
From above two equations, it follows that $b_1$ and $b_4$ are also constant on $M$.
Now, we know that $M$ has constant principal curvatures and constant $C$.
According to Theorem 1.1 of \cite{GMY}, it can be directly checked that
$M$ is an open part of $M_\tau$ for some $\tau<-1$.

If $b_2=0$ on $M$, then by Theorem 3.11 of \cite{GMY}, we know that $M$ is an open part of $M_{k_1,k_2}^c$ given in Example \ref{exam:3.4fg1}.

%$M$ is locally given by map $\Phi: \Omega\subset\mathbb{R}^{3}\longrightarrow \mathbb{H}^2\times\mathbb{H}^2$:
%\ $(t,r,s)\longrightarrow(p,q)$ defined by \eqref{eqn:6.1}.

For the ``if" part, it is a directly consequence from verifying that hypersurfaces $M_\tau$ and
$M_{k_1,k_2}^c$ have constant $C$ and constant $b_2$.
\end{proof}


%\begin{remark}\label{rem:3.1aaaa}
%For any given $C\neq\pm1$, Proposition \ref{prop:3.61} (2) constructs abundant hypersurfaces
%with such product angle function $C$.
%\end{remark}


Now, combining Lemma \ref{lemma:2.1aaa1}, Theorem \ref{thm:3.1} and Proposition \ref{prop:3.61},
we can obtain the following classification result.
\begin{theorem}\label{thm:3.2aa}
Let $M$ be a minimal hypersurface of $\mathbb{H}^2\times \mathbb{H}^2$ with constant $C=0$.
Then, either
%
\begin{enumerate}
\item[(1)]
$M$ is an open part of $M_{k_1,k_1}^{1/2}$ given in Example \ref{exam:3.4fg1} for some constant $k_1$; or

\item[(2)]
$M$ is an open part of a minimal hypersurface described by Theorem \ref{thm:3.1}.
\end{enumerate}
\end{theorem}
%
\begin{proof}
By Lemma \ref{lemma:2.1aaa1}, under the frame $\{E_1,E_2,E_3\}$ defined by \eqref{eqn:2.5}, there are two cases
for $b_1,b_2,b_4$ on $M$, and the {\bf Case 2} has been classified by Theorem \ref{thm:3.1}.

For {\bf Case 1}, by Proposition \ref{prop:3.61}, $M$ is an open part of $M_{k_1,k_2}^c$ given in Example \ref{exam:3.4fg1}.
Then, by Lemma \ref{lemma:3.5fg12}, we know that $c=\frac{1}{2}$, $k_1(r)$ and $k_2(s)$
are equal to the same constant functions.
%In the following, we need to
%determine the specific curves $\gamma(r)$ and $\tilde{\gamma}(s)$.
%We assume that $\kappa(r)$ and $\tilde{\kappa}(s)$ are curvatures of $\gamma(r)$ and
%$\tilde{\gamma}(s)$, respectively. Then, for $\gamma(r)$ and $\tilde{\gamma}(s)$, it holds
%\begin{equation*}
%\begin{aligned}
%\frac{d^2\gamma(r)}{dr^2}&=\gamma(r)+\kappa(r)N(r), \ \
%\frac{dN(r)}{dr}=-\kappa(r)\frac{d\gamma(r)}{dr},\\
%\frac{d^2\tilde{\gamma}(s)}{ds^2}&=\tilde{\gamma}(s)+\tilde{\kappa}(s)\tilde{N}(s), \ \
%\frac{d\tilde{N}(s)}{ds}=-\tilde{\kappa}(s)\frac{d\tilde{\gamma}(s)}{ds}.
%\end{aligned}
%\end{equation*}
%It follows that
%$E_1=(\frac{d\gamma(r)}{d r},0),\ E_2=(0,\frac{d\tilde{\gamma}(s)}{d s}),
%\ E_3=(\frac{\partial p}{\partial t},\frac{\partial q}{\partial t})$,
%$N=(\frac{\partial p}{\partial t},-\frac{\partial q}{\partial t})$
%and
%$$
%\begin{aligned}
%b_1&=g(AE_1,E_1)=-\tfrac{1}{\sqrt{2}}\frac{\sinh(\frac{t}{\sqrt{2}})
%-\cosh(\frac{t}{\sqrt{2}})\kappa(r)}{\cosh(\frac{t}{\sqrt{2}})
%-\sinh(\frac{t}{\sqrt{2}})\kappa(r)}, \ \ b_2=g(AE_1,E_2)=0,\\
%b_4&=g(AE_2,E_2)=\tfrac{1}{\sqrt{2}}\frac{\sinh(\frac{t}{\sqrt{2}})
%-\cosh(\frac{t}{\sqrt{2}})\tilde{\kappa}(s)}{\cosh(\frac{t}{\sqrt{2}})
%-\sinh(\frac{t}{\sqrt{2}})\tilde{\kappa}(s)}.
%\end{aligned}
%$$
%Then, by $H=\frac{1}{3}(b_1+b_4)=0$, we have that $\kappa(r)$ and $\tilde{\kappa}(s)$
%are equal to the same constant. %Up to an isometry of $\mathbb{H}^2\times \mathbb{H}^2$,
%...
\end{proof}







\section{Proof of Theorem \ref{thm:1.1}}\label{sect:5}
%
In this section, we prove Theorem \ref{thm:1.1}.
Firstly, by using the so-called Tsinghua principle, we establish a very useful identity
on any hypersurface $M$ of $\mathbb{H}^2\times\mathbb{H}^2$ with constant sectional curvature.


\begin{lemma}\label{lemma:3.2aa}
Let $M$ be a hypersurface of $\mathbb{H}^2\times \mathbb{H}^2$ with constant sectional curvature $\kappa$.
Then, for any tangent vector fields $W,U,Y,Z\in TM$, the following equation holds:
%
\begin{equation}\label{eqn:3.1}
\mathop\mathfrak{S}\limits_{W,U,Y}\mathbf{I}(W,U,Y,Z)=0,
\end{equation}
where the symbol $\mathfrak{S}$ stands for the cyclic summation,
and $\mathbf{I}(W,U,Y,Z)$ is defined by
\begin{equation*}
\mathbf{I}(W,U,Y,Z):=-\tfrac{1}{2}\{-g(TY,Z)g(AW,TU)+g(TU,Z)g(AW,TY)\}.
\end{equation*}
\end{lemma}
%
\begin{proof}

In order to prove Lemma \ref{lemma:3.2aa}, we calculate the expression of the cyclic
summation
%
\begin{equation}\label{eqn:3.2}
\mathfrak{A}:=\mathop\mathfrak{S}\limits_{W,U,Y} \big\{g((\nabla^2 A)(W,U,Y),Z)-
g((\nabla^2 A)(W,Y,U),Z)\big\}
\end{equation}
in two different ways. On the one hand, we take the covariant derivative
of the Codazzi equation \eqref{eqn:2.3}, we have
\begin{equation*}
\begin{split}
&g((\nabla^2 A)(W,U,Y),Z)-g((\nabla^2 A)(W,Y,U),Z)\\
&=-\tfrac{1}{2}\{g((\nabla_W T)Y,Z)\mu(U)+g(TY,Z)(\nabla_W \mu)U\\
&\ \ -g((\nabla_W T)U,Z)\mu(Y)-g(TU,Z)(\nabla_W \mu)Y\}.
\end{split}
\end{equation*}
By direct calculations, with the use of Lemma \ref{lemma:2.1},
we can have
\begin{equation*}
\begin{split}
g&((\nabla^2 A)(W,U,Y),Z)-g((\nabla^2 A)(W,Y,U),Z)\\
&=-\tfrac{1}{2}\{g(AW,Y)\mu(Z)\mu(U)-g(TY,Z)g(AW,TU)+Cg(TY,Z)g(AW,U)\\
&\ \ -g(AW,U)\mu(Z)\mu(Y)+g(TU,Z)g(AW,TY)-Cg(TU,Z)g(AW,Y)\}.
\end{split}
\end{equation*}
Thus, it holds that
%
\begin{equation}\label{eqn:3.3aa}
\mathfrak{A}=\mathop\mathfrak{S}\limits_{W,U,Y}\mathbf{I}(W,U,Y,Z).
\end{equation}

On the other hand, the cyclic sum $\mathfrak{A}$ can be rewritten as
%
\begin{equation}\label{eqn:3.3}
\mathfrak{A}=\mathop\mathfrak{S}\limits_{W,U,Y}\big\{g((\nabla^{2}A)(W,U,Y),Z)
-g((\nabla^{2}A)(U,W,Y),Z)\big\}.
\end{equation}
%
Then by applying the Ricci identity \eqref{eqn:ric}, it follows from \eqref{eqn:3.3} that
%
\begin{equation}\label{eqn:3.4}
\mathfrak{A}=-\mathop\mathfrak{S}\limits_{W,U,Y} \big\{g(R(W,U)Y,AZ)
+g(R(W,U)Z,AY)\big\}.
\end{equation}


Since $M$ has constant sectional curvature $\kappa$, we have
%
\begin{equation}\label{eqn:3.6aa}
g(R(W,U)Y,Z)=\kappa\{g(U,Y)g(W,Z)-g(W,Y)g(U,Z)\},
\end{equation}
%
it follows directly that
$$
\mathfrak{A}=-\mathop\mathfrak{S}\limits_{W,U,Y} \big\{g(R(W,U)Y,AZ)
+g(R(W,U)Z,AY)\big\}=0.
$$
Thus, from \eqref{eqn:3.3aa}, we have completed the proof of Lemma \ref{lemma:3.2aa}.
\end{proof}

%\begin{remark}\label{rem:4.1aaaa}

%\end{remark}

Next, we prove that the sectional curvature of $M$ must be
$\kappa=-\tfrac{1}{2}$.
\begin{proposition}\label{prop:4.5}
Let $M$ be a hypersurface of $\mathbb{H}^2\times\mathbb{H}^2$ with constant sectional curvature $\kappa$.
Then, $\kappa=-\tfrac{1}{2}$.
\end{proposition}
%
\begin{proof}
Assume that the hypersurface $M$ has constant sectional curvature $\kappa$.
We will discuss two cases depending on the value of
the product angle function $C$ on $M$.

{\bf Case I}. $C\neq\pm1$.

In this case, we can take the local orthonormal frame fields $\{E_1,E_2,E_3\}$
defined by \eqref{eqn:2.5}.
Then, the shape operator $A$, $(1,1)$-tensor $T$ and $1$-form $\mu$
are given by \eqref{eqn:2.6}--\eqref{eqn:2.7}.
By using $\bar{\nabla}P=0$, \eqref{eqn:2.6} and Lemma \ref{lemma:2.1},
we can obtain the connections with respect to $\{E_1,E_2,E_3\}$ as follows:
%
\begin{equation}\label{eqn:LL}
\left\{
\begin{aligned}
&E_1C=-2b_3\sqrt{1-C^2},\ E_2C=-2b_5\sqrt{1-C^2},\ E_3C=-2b_6\sqrt{1-C^2},\\ 		
&\nabla_{E_1}E_1=b_1\sqrt{\tfrac{1-C}{1+C}}E_3,\
\nabla_{E_1}E_2=-b_2\sqrt{\tfrac{1+C}{1-C}}E_3,\ \nabla_{E_1}E_3=-b_1\sqrt{\tfrac{1-C}{1+C}}E_1
+b_2\sqrt{\tfrac{1+C}{1-C}}E_2,\\
&\nabla_{E_2} E_1=b_2\sqrt{\tfrac{1-C}{1+C}}E_3,\
\nabla_{E_2} E_2=-b_4\sqrt{\tfrac{1+C}{1-C}}E_3,\ \nabla_{E_2}E_3=-b_2\sqrt{\tfrac{1-C}{1+C}}E_1
+b_4\sqrt{\tfrac{1+C}{1-C}}E_2,\\
&\nabla_{E_3}E_1=b_3\sqrt{\tfrac{1-C}{1+C}}E_3,\
\nabla_{E_3}E_2=-b_5\sqrt{\tfrac{1+C}{1-C}}E_3,\ \nabla_{E_3}E_3=-b_3\sqrt{\tfrac{1-C}{1+C}}E_1
+b_5\sqrt{\tfrac{1+C}{1-C}}E_2.
\end{aligned}\right.
\end{equation}


Taking in \eqref{eqn:3.1}, respectively, %$1\leq i,j,k,l\leq 3$,
%
$$
(W,U,Y,Z)=(E_1,E_2,E_3,E_1), \ (E_1,E_2,E_3,E_2), \ (E_1,E_2,E_3,E_3),
$$
we obtain $b_3(1+C)=b_5(1-C)=b_2C=0$.
Since $C\neq\pm1$, we have $b_3=b_5=0$. In the following, we further divide {\bf Case I} into two subcases depending on whether $C$ is $0$ or not.

{\bf Case I-i}: $C\neq0$.

In this subcase, we have $b_3=b_5=b_2=0$. Then, by direct calculations, we get
$$
g(R(E_1,E_2)E_2,E_1)=b_1b_4,\
g(R(E_1,E_3)E_3,E_1)=\tfrac{C}{2}-\tfrac{1}{2}+b_1b_6,\
g(R(E_2,E_3)E_3,E_2)=b_4b_6-\tfrac{1}{2}-\tfrac{C}{2}.
$$
By $C\neq\pm1$, it follows $\kappa\neq0$, $b_6=\tfrac{2\kappa+1-C}{2b_1}=\tfrac{2\kappa+1+C}{2b_4}$ 
and $b_6^2=\tfrac{(2\kappa+1)^2-C^2}{4\kappa}$. If $\kappa>0$, without loss of generality, we have
\begin{equation}\label{eqn:6.2}
b_6=-\tfrac{\sqrt{(2\kappa+1)^2-C^2}}{2\sqrt{\kappa}},\
b_1=-\tfrac{\sqrt{\kappa}(2\kappa+1-C)}{\sqrt{(2\kappa+1)^2-C^2}},
\ b_4=-\tfrac{\sqrt{\kappa}(2\kappa+1+C)}{\sqrt{(2\kappa+1)^2-C^2}}.
\end{equation}

Next, taking $(Y,Z)=(E_3,E_1)$ and $(Y,Z)=(E_3,E_2)$ into Codazzi equation \eqref{eqn:2.3},
with the use of $b_2=b_3=b_5=0$, \eqref{eqn:2.6} and \eqref{eqn:LL},
we compare the components of $E_1$ and $E_2$, respectively, then we can have
%Then, we check the Codazzi equations \eqref{eqn:2.17} and \eqref{eqn:2.21}, we have
\begin{equation}\label{eqn:6.3}
E_3b_1=-\tfrac{\sqrt{1-C}(1+C-2b_1(b_1-b_6))}{2\sqrt{1+C}},\
E_3b_4=\tfrac{\sqrt{1+C}(1-C-2b_4(b_4-b_6))}{2\sqrt{1-C}}.
\end{equation}
In the following, we take the derivatives of $b_1$ and $b_4$ with respect to $E_3$,
with the use of \eqref{eqn:6.2} and $E_3C=-2b_6\sqrt{1-C^2}$, we obtain
\begin{equation}\label{eqn:6.4}
E_3b_1=\tfrac{\sqrt{1-C^2}(2\kappa+1)}{1+C+2\kappa},\
E_3b_4=\tfrac{\sqrt{1-C^2}(2\kappa+1)}{-1+C-2\kappa}.
\end{equation}
By \eqref{eqn:6.2}, \eqref{eqn:6.3} and \eqref{eqn:6.4}, we have $\sqrt{1-C^2}(2\kappa+1)=0$,
which implies that $\kappa=-\tfrac{1}{2}$.
But, this contradicts with $\kappa>0$. Thus, it holds $\kappa<0$.

If $\kappa<0$, without loss of generality, we have
\begin{equation}\label{eqn:6.5}
b_6=-\tfrac{\sqrt{-(2\kappa+1)^2+C^2}}{2\sqrt{-\kappa}},\
b_1=-\tfrac{\sqrt{-\kappa}(2\kappa+1-C)}{\sqrt{-(2\kappa+1)^2+C^2}},
\ b_4=-\tfrac{\sqrt{-\kappa}(2\kappa+1+C)}{\sqrt{-(2\kappa+1)^2+C^2}}.
\end{equation}

By taking $(Y,Z)=(E_3,E_1)$ and $(Y,Z)=(E_3,E_2)$ into Codazzi equation \eqref{eqn:2.3},
we compare the components of $E_1$ and $E_2$ respectively, then we still have \eqref{eqn:6.3}.
%Then, we check the Codazzi equations \eqref{eqn:2.17} and \eqref{eqn:2.21}, we have
%\begin{equation}\label{eqn:6.6}
%E_3b_1=-\tfrac{\sqrt{1-C}(1+C-2b_1(b_1-b_6))}{2\sqrt{1+C}},\
%E_3b_4=\tfrac{\sqrt{1+C}(1-C-2b_4(b_4-b_6))}{2\sqrt{1-C}}.
%\end{equation}
Taking the derivatives of $b_1$ and $b_4$ in \eqref{eqn:6.5} with respect to $E_3$,
with the use of \eqref{eqn:6.5} and $E_3C=-2b_6\sqrt{1-C^2}$, then comparing with \eqref{eqn:6.3}, we can obtain
%by similar discussion, we can have
$$
\kappa=-\frac{1}{2},\ b_1=\frac{C}{\sqrt{2C^2}},\ b_4=-\frac{C}{\sqrt{2C^2}},\ b_6=-\frac{\sqrt{C^2}}{\sqrt{2}}.
$$

Notice that, the isometry map
$$
\mathbb{F}:\mathbb{H}^2\times \mathbb{H}^2\rightarrow\mathbb{H}^2\times \mathbb{H}^2,\ \
\mathbb{F}(p,q)=(q,p),
$$
satisfies $d \mathbb{F}\circ P=-P \circ d \mathbb{F}$. If $C<0$ on an open subset $\Omega$ of $M$, then using the isometry $\mathbb{F}$, we can obtain another hypersurface $\mathbb{F}(M)$ with constant sectional curvature and $C>0$ on $\mathbb{F}(\Omega)$. 
Thus, we can assume that $C>0$ on an open subset $\Omega$ of $M$, and have
$$
b_1=\frac{1}{\sqrt{2}},\ b_4=-\frac{1}{\sqrt{2}},\ b_6=-\frac{C}{\sqrt{2}},\ b_2=b_3=b_5=0.
$$
%When $C<0$ in local, we have
%$$
%b_1=-\frac{1}{\sqrt{2}},\ b_4=\frac{1}{\sqrt{2}},\ b_6=\frac{C}{\sqrt{2}},\ b_2=b_3=b_5=0.
%$$
%Note that, any isometry $\mathbb{F}$ of $\mathbb{H}^2\times \mathbb{H}^2$ satisfies $d \mathbb{F}\circ P=\pm P \circ d \mathbb{F}$. Then we can use an isometry $\mathbb{F}$ with $d \mathbb{F}\circ P=-P \circ d \mathbb{F}$
%to obtain $\mathbb{F}(M)$ with constant section curvature and $C>0$ on $\mathbb{F}(M)$.
By directly integrating $E_3C=C\sqrt{2(1-C^2)}$ along the integral curves of $E_3$ and $C>0$ on $\Omega$, we have
$$
C=\frac{1}{\cosh(\sqrt{2}\tilde{t}+\tilde{f})},
$$
where $\tilde{t},\tilde{f}$ are functions on $\Omega$ satisfying $E_3\tilde{t}=1$ and $E_3\tilde{f}=0$.
According to the fact that $C$ is a continuous function on $M$, we know that $C>0$ on $M$. 


{\bf Case I-ii}: $C=0$.

In this subcase, by Lemma \ref{lemma:2.1}, we have $b_3=b_5=b_6=0$.
By Gauss equation \eqref{eqn:2.2}, we have
$$
g(R(E_1,E_2)E_2,E_1)=b_1b_4-b_2^2, \ \ g(R(E_1,E_3)E_3,E_1)=g(R(E_2,E_3)E_3,E_2)=-\tfrac{1}{2}.
$$
It implies that the sectional curvature of $M$ is $\kappa=-\tfrac{1}{2}$
and it holds $b_1b_4-b_2^2=-\tfrac{1}{2}$.

{\bf Case II}. $C^2=1$.

In this case, by Lemma \ref{lem:2.2}, we know that $M$ is an open part of
$\Gamma\times \mathbb{H}^2$, where $\Gamma$
is a curve of $\mathbb{H}^2$. It follows that the tangent bundle of $M$ has
a decomposition $TM=T\Gamma\oplus T\mathbb{H}^2$,
and we have $A|_{T\mathbb{H}^2}=0$.

Now, we take an orthonormal basic $\{v_1,v_2,v_3\}$ with $v_1\in T\Gamma$ and $v_2,v_3\in T\mathbb{H}^2$.
By direct calculations, we have $g(R(v_1,v_2)v_2,v_1)=0$ and $g(R(v_2,v_3)v_3,v_2)=-1$.
Therefore, in this case, there exists no hypersurface with constant sectional curvature.
\end{proof}

%\begin{remark}\label{rem:4.1aa}
%Notice that, if hypersurface $M$ has constant sectional curvature, then $M$ does not have
%constant mean curvature. In fact, by Proposition \ref{prop:4.5}, $M$ has constant $C=0$
%and $\kappa=\frac{1}{2}$. According to Corollary 1 (3) of \cite{Ur},
%%the scalar curvature and the function $C$ constants. But $M_t$ does not have
%constant sectional curvature.
%\end{remark}

In the following, we discuss the {\bf Case I-i} of Proposition \ref{prop:4.5}.

\begin{theorem}\label{thm:4.1abc}
Let $M$ be a hypersurface of $\mathbb{H}^2\times \mathbb{H}^2$ with constant sectional curvature $-\frac{1}{2}$.
If under the frame field $\{E_1,E_2,E_3\}$ defined by \eqref{eqn:2.5}, it holds
$$
C>0,\ b_1=\frac{1}{\sqrt{2}},\ b_4=-\frac{1}{\sqrt{2}},\ b_6=-\frac{C}{\sqrt{2}},\ b_2=b_3=b_5=0.
$$
%Then, $M$ is locally given by the immersion
%$\Phi:\Omega\subset\mathbb{R}^3\rightarrow\mathbb{H}^2\times\mathbb{H}^2$:
%$(t,r,s)\rightarrow(p(t,r),q(t,s))$, where $t<0$ and
%
%\begin{equation}\label{eqn:jr}
%\begin{aligned}
%&p(t,r)=\sqrt{\cosh(\sqrt{2}t)+1}(\cosh(r),\sinh(r),0)+\sqrt{\cosh(\sqrt{2}t)}(0,0,1),\\
%&q(t,s)=\sqrt{\cosh(\sqrt{2}t)}(1,0,0)+\sqrt{\cosh(\sqrt{2}t)-1}(0,\cos(s),\sin(s)),
%&p(t,r)=\Big(\sqrt{\cosh(\sqrt{2}t)+1}\cosh(r),\sqrt{\cosh(\sqrt{2}t)+1}\sinh(r),\sqrt{\cosh(\sqrt{2}t)}\Big),\\
%&q(t,s)=\Big(\sqrt{\cosh(\sqrt{2}t)},\sqrt{\cosh(\sqrt{2}t)-1}\cos(s),\sqrt{\cosh(\sqrt{2}t)-1}\sin(s)\Big),
%\end{aligned}	
%\end{equation}
Then, $M$ is an open part of Example \ref{exam:4.1ss1ss11ab1}.
\end{theorem}
%
\begin{proof}
By $b_3=b_5=0$, $b_6=-\frac{C}{\sqrt{2}}$ and $\nabla C=-2AX$, we have that product angle function satisfies
$$
E_1C=E_2C=0,\ E_3C=C\sqrt{2(1-C^2)}.
$$
Now, we choose two appropriate non-zero functions $\rho_1$ and $\rho_2$, such that
the new frame
$$
\{X_1=\rho_1E_1,\ X_2=\rho_2E_2,\ X_3=E_3\}
$$
satisfies that $[X_1,X_2]=[X_1,X_3]=[X_2,X_3]=0$. Specifically, by using
$[E_i,E_j]=\nabla_{E_i} {E_j}-\nabla_{E_j} {E_i}$, \eqref{eqn:LL} and $b_2=b_3=b_5=0$, we get
$$
[E_1,E_2]=0,\ [E_1,E_3]=-\sqrt{\tfrac{1-C}{2(1+C)}}E_1,\ [E_2,E_3]=-\sqrt{\tfrac{1+C}{2(1-C)}}E_2.
$$
Therefore
$$
\begin{aligned}
{\text [X_1,X_3]}&=-(\sqrt{\tfrac{1-C}{2(1+C)}}\rho_1+E_3\rho_1)E_1.
\end{aligned}
$$
Then $[X_1,X_3]=0$ is equivalent to
%
\begin{equation}\label{eqn:m1}
E_3\rho_1=-\sqrt{\tfrac{1-C}{2(1+C)}}\rho_1.
\end{equation}
Similarly, we get
$$
{\text [X_2,X_3]}=(-\sqrt{\tfrac{1+C}{2(1-C)}}\rho_2-E_3\rho_2)E_2,\ \
{\text [X_1,X_2]}=\rho_1(E_1\rho_2)E_2-\rho_2(E_2\rho_1)E_1.
$$
Then $[X_2,X_3]=[X_1,X_2]=0$ are equivalent to
%
\begin{equation}\label{eqn:m2}
E_3\rho_2=-\sqrt{\tfrac{1+C}{2(1-C)}}\rho_2,
\end{equation}
%
\begin{equation}\label{eqn:dm1dm4}
\rho_1 E_1\rho_2=\rho_2E_2\rho_1=0.
\end{equation}

Let $\rho_1=\sqrt{\frac{1+C}{2C}}$, $\rho_2=\sqrt{\frac{1-C}{2C}}$, then by $E_1C=E_2C=0$ and $E_3C=C\sqrt{2(1-C^2)}$,
we can check that \eqref{eqn:m1}--\eqref{eqn:dm1dm4} hold, which implies that $[X_1,X_2]=[X_1,X_3]=[X_2,X_3]=0$.

By the definition of the frame $\{X_1,X_2,X_3\}$ and using $[X_1,X_2]=[X_1,X_3]=[X_2,X_3]=0$,
we can identify $M$ with an open subset $\Omega_1$ of $\mathbb{R}^3$ locally and express the hypersurface $M$ by an immersion
$$
\begin{aligned}
\Phi_1:\Omega_1
\longrightarrow\mathbb{H}^2\times\mathbb{H}^2,\ \ \
(t_1,r,s)\mapsto (p(t_1,r,s),q(t_1,r,s)),
\end{aligned}
$$
such that $(\frac{\partial p}{\partial t_1},\frac{\partial q}{\partial t_1})=E_3$,
$(\frac{\partial p}{\partial r},\frac{\partial q}{\partial r})=\rho_1E_1$ and
$(\frac{\partial p}{\partial s},\frac{\partial q}{\partial s})=\rho_2E_2$.


By $C\neq0$, we integrate $E_3C=C\sqrt{2(1-C^2)}$ along the integral curves of $E_3$, and obtain
$$
C=\frac{1}{\cosh(\sqrt{2}t_1+\tilde{h})},
$$
where $\tilde{h}$ is a function on $M$ satisfying $E_3\tilde{h}=0$ and
$\sqrt{2}t_1+\tilde{h}<0$.


By $E_1C=E_2C=0$, $(\frac{\partial p}{\partial r},\frac{\partial q}{\partial r})=\rho_1E_1$ and
$(\frac{\partial p}{\partial s},\frac{\partial q}{\partial s})=\rho_2E_2$, we have that
$$
\frac{\partial C}{\partial r}=\frac{\partial C}{\partial s}=0\ \  {\rm and}\ \
\frac{\partial \tilde{h}}{\partial r}=\frac{\partial \tilde{h}}{\partial s}=0.
$$
It follows from $E_3\tilde{h}=0$ that $\tilde{h}$ is constant function. Thus, without loss of generality,
we can identify $M$ with an open subset $\Omega$ of $\mathbb{R}^3$ locally and express the hypersurface $M$ by an immersion
$$
\begin{aligned}
\Phi:\Omega
\longrightarrow\mathbb{H}^2\times\mathbb{H}^2,\ \ \
(t,r,s)\mapsto (p(t,r,s),q(t,r,s)),
\end{aligned}
$$
such that $(\frac{\partial p}{\partial t},\frac{\partial q}{\partial t})=E_3$,
$(\frac{\partial p}{\partial r},\frac{\partial q}{\partial r})=\rho_1E_1$,
$(\frac{\partial p}{\partial s},\frac{\partial q}{\partial s})=\rho_2E_2$ and
$C=\frac{1}{\cosh(\sqrt{2}t)}$, where $t=t_1+\frac{\tilde{h}}{\sqrt{2}}<0$. By $E_3=(\frac{\partial p}{\partial t},\frac{\partial q}{\partial t})$ and
$PN=CN+\sqrt{1-C^2}E_3$, we have $N=(\sqrt{\frac{1+C}{1-C}}\frac{\partial p}{\partial t},-\sqrt{\frac{1-C}{1+C}}\frac{\partial q}{\partial t})$.

From the definition of $P$  and using
$$
PE_1=E_1,\ \ PE_2=-E_2,
$$
we obtain that $dp,dq:T(\Omega)\rightarrow
T\mathbb{H}^2$ satisfy
\begin{equation}\label{eqn:5.49}
\left\{
\begin{aligned}
(dp(\tfrac{\partial}{\partial r}),0)&=\tfrac{1}{2}(d\Phi(\tfrac{\partial}{\partial r})+Pd\Phi(\tfrac{\partial}{\partial r}))=d\Phi(\tfrac{\partial}{\partial r}),\\
(0,dq(\tfrac{\partial}{\partial r}))&=\tfrac{1}{2}(d\Phi(\tfrac{\partial}{\partial r})-Pd\Phi(\tfrac{\partial}{\partial r}))=0,
\end{aligned}
\right.
\end{equation}
%
%\vskip-3mm
%
and
\begin{equation}\label{eqn:5.50}
\left\{
\begin{aligned}
(dp(\tfrac{\partial}{\partial s}),0)&=\tfrac{1}{2}(d\Phi(\tfrac{\partial}{\partial s})+Pd\Phi(\tfrac{\partial}{\partial s}))=0,\\
(0,dq(\tfrac{\partial}{\partial s}))&=\tfrac{1}{2}(d\Phi(\tfrac{\partial}{\partial s})-Pd\Phi(\tfrac{\partial}{\partial s}))=d\Phi(\tfrac{\partial}{\partial s}).
\end{aligned}
\right.
\end{equation}

The first equation of \eqref{eqn:5.50} shows that $p$ depends only on the $(t,r)$.
Similarly, from the second equation in \eqref{eqn:5.49} we derive that $q$ depends only on
the $(t,s)$. By $g(PE_3,E_3)=-C$, $AE_3=-\frac{C}{\sqrt{2}}E_3$ and $\nabla_{E_3}{E_3}=0$,
we have that
\begin{equation}\label{eqn:4.mm1}
\begin{aligned}
(\tfrac{\partial^2 p}{\partial t^2},\tfrac{\partial^2 q}{\partial t^2})&%=D_{E_3}{E_3}
=\bar{\nabla}_{E_3}{E_3}-g( D_{E_3}{E_3},\tfrac{(p,q)}{\sqrt{2}})\tfrac{(p,q)}{\sqrt{2}}
-g( D_{E_3}{E_3},\tfrac{(p,-q)}{\sqrt{2}})\tfrac{(p,-q)}{\sqrt{2}}\\
&=-\tfrac{C}{\sqrt{2}}N+\tfrac{1}{2}(p,q)-\tfrac{C}{2}(p,-q)\\
&=(-C\sqrt{\tfrac{1+C}{2(1-C)}}\tfrac{\partial p}{\partial t},C\sqrt{\tfrac{1-C}{2(1+C)}}\tfrac{\partial q}{\partial t})+(\tfrac{1-C}{2}p,\tfrac{1+C}{2}q),
\end{aligned}
\end{equation}
where $D$ is the canonical Euclidean connection on $\mathbb{R}_2^6$.
From $\frac{dC}{dt}=E_3C=C\sqrt{2(1-C^2)}$, we know that the solution of
\eqref{eqn:4.mm1} is given by
\begin{equation}\label{eqn:pq}
\begin{aligned}
p(t,r)=\sqrt{\tfrac{1+C}{C}}V_1(r)+\tfrac{1}{\sqrt{C}}V_2(r),\ \
q(t,s)=\tfrac{1}{\sqrt{C}}W_1(s)+\sqrt{\tfrac{1-C}{C}}W_2(s),
\end{aligned}
\end{equation}
where $V_1(r),V_2(r),W_1(s),W_2(s)\in \mathbb{R}_1^3$ are vector-value functions depending on variables $r$ and $s$. Thus, we have
$$
\begin{aligned}
E_1&=\tfrac{1}{\rho_1}(\tfrac{\partial p}{\partial r},\tfrac{\partial q}{\partial r})
%=\sqrt{\frac{2C}{1+C}}(\sqrt{\frac{1+C}{C}}\frac{d V_1(r)}{dr}+\frac{1}{\sqrt{C}}\frac{d V_2(r)}{dr},\vec{0})
=(\sqrt{2}\tfrac{d V_1(r)}{dr}+\sqrt{\tfrac{2}{1+C}}\tfrac{d V_2(r)}{dr},0),\\
E_2&=\tfrac{1}{\rho_2}(\tfrac{\partial p}{\partial s},\tfrac{\partial q}{\partial s})
%=\sqrt{\frac{2C}{1-C}}(\vec{0},\frac{1}{\sqrt{C}}\frac{d W_1(s)}{ds}+\sqrt{\frac{1-C}{C}}\frac{d W_2(s)}{ds})
=(0,\sqrt{\tfrac{2}{1-C}}\tfrac{d W_1(s)}{ds}+\sqrt{2}\tfrac{d W_2(s)}{ds}),\\
E_3&=(\tfrac{\partial p}{\partial t},\tfrac{\partial q}{\partial t})
=(-\sqrt{\tfrac{1-C}{2C}}V_1(r)-\sqrt{\tfrac{1-C^2}{2C}}V_2(r),
-\sqrt{\tfrac{1-C^2}{2C}}W_1(s)-\sqrt{\tfrac{1+C}{2C}}W_2(s)),
\end{aligned}
$$
and
$$
N=(\sqrt{\tfrac{1+C}{1-C}}\tfrac{\partial p}{\partial t},-\sqrt{\tfrac{1-C}{1+C}}\tfrac{\partial q}{\partial t})
=(-\sqrt{\tfrac{1+C}{2C}}V_1(r)-\tfrac{1+C}{\sqrt{2C}}V_2(r),
\tfrac{1-C}{\sqrt{2C}}W_1(s)+\sqrt{\tfrac{1-C}{2C}}W_2(s)).
$$



From the fact that $\langle p(t,r),p(t,r)\rangle=\langle q(t,s),q(t,s)\rangle=-1$, we have
$$
\begin{aligned}
&1+\langle V_1(r),V_1(r)\rangle+\tfrac{1}{C}(\langle V_1(r),V_1(r)\rangle+\langle V_2(r),V_2(r)\rangle)
+\tfrac{2\sqrt{1+C}}{C}\langle V_1(r),V_2(r)\rangle=0,\\
&1-\langle W_2(s),W_2(s)\rangle+\tfrac{1}{C}(\langle W_1(s),W_1(s)\rangle+\langle W_2(s),W_2(s)\rangle)
+\tfrac{2\sqrt{1-C}}{C}\langle W_1(s),W_2(s)\rangle=0.
\end{aligned}
$$
According to the independence of parameters $t,r$ and $s$, it follows from above two equations that
$$
\begin{aligned}
&\langle V_1(r),V_1(r)\rangle=-1,\ \langle V_1(r),V_2(r)\rangle=0,\ \langle V_2(r),V_2(r)\rangle=1.\\
&\langle W_1(s),W_1(s)\rangle=-1,\ \langle W_1(s),W_2(s)\rangle=0,\ \langle W_2(s),W_2(s)\rangle=1.
\end{aligned}
$$
Then, by $g(E_1,E_3)=g(E_2,E_3)=0$,
with the use of $\langle \frac{dV_1(r)}{dr},V_2(r)\rangle+\langle V_1(r),\frac{dV_2(r)}{dr}\rangle=0$
and $\langle \frac{dW_1(s)}{ds},W_2(s)\rangle+\langle W_1(s),\frac{dW_2(s)}{ds}\rangle=0$,
we can have
$$
\begin{aligned}
&\langle \frac{d V_1(r)}{dr},V_2(r)\rangle=\langle \frac{d W_1(s)}{ds},W_2(s)\rangle=0.
\end{aligned}
$$

Now, on one hand, for $V_1(r)\in \mathbb{H}^2$, if $V_1(r)$ is a constant vector, we have
$\frac{d V_1(r)}{dr}=0$ and $\frac{d V_2(r)}{dr}\neq0$. Then, it holds that
$$
\begin{aligned}
E_1&=\sqrt{\tfrac{2}{1+C}}(\tfrac{d V_2(r)}{d r},0),\\
D_{E_1}N&=\sqrt{\tfrac{2C}{1+C}}\tfrac{\partial N}{\partial r}=
\sqrt{\tfrac{2C}{1+C}}(-\tfrac{1+C}{\sqrt{2C}}\tfrac{d V_2(r)}{d r},0)=-\tfrac{1+C}{\sqrt{2}}E_1,
\end{aligned}
$$
which implies that $b_1=-\langle D_{E_1}N,E_1\rangle=\frac{1+C}{\sqrt{2}}$.
By the assumption that $b_1=\frac{1}{\sqrt{2}}$, we have $C=0$. It is a contradiction.
Thus, $V_1(r)$ is a curve of $\mathbb{H}^2$,
and $V_2(r)$ is a unit normal vector field of $V_1(r)\hookrightarrow\mathbb{H}^2$.
By taking re-parameterization to make $r$ being arc length parameter of $V_1(r)$, and assume that
$k(r)$ is the curvature of $V_1(r)$, we have
$$
\frac{d^2V_1(r)}{dr^2}=V_1(r)+k(r)V_2(r), \ \
\frac{dV_2(r)}{dr}=-k(r)\frac{dV_1(r)}{dr}.
$$
Now, we can have
$$
\begin{aligned}
E_1&=\Big((\sqrt{2}-k(r)\sqrt{\tfrac{2}{1+C}})\tfrac{d V_1(r)}{dr},0\Big),\\
D_{E_1}N&=\sqrt{\tfrac{2C}{1+C}}\tfrac{\partial N}{\partial r}=
((-1+k(r)\sqrt{1+C})\tfrac{d V_1(r)}{d r},0)=\tfrac{-1+k(r)\sqrt{1+C}}{\sqrt{2}-k(r)\sqrt{\tfrac{2}{1+C}}}E_1,
\end{aligned}
$$
which implies that $b_1=-\langle D_{E_1}N,E_1\rangle=\frac{1-k(r)\sqrt{1+C}}{\sqrt{2}-k(r)\sqrt{\frac{2}{1+C}}}$.
By the assumption that $b_1=\frac{1}{\sqrt{2}}$ and $C\neq0$, we have $k(r)=0$,
i.e., $V_1(r)$ is a geodesic of $\mathbb{H}^2$.

On the other hand, for $W_1(s)\in \mathbb{H}^2$, if $W_1(s)$ is not constant vector,
then $W_1(s)$ is a curve of $\mathbb{H}^2$,
and $W_2(s)$ is a unit normal vector field of $W_1(s)\hookrightarrow\mathbb{H}^2$.
By taking re-parameterization to make $s$ being arc length parameter of $W_1(s)$, and assume that
$\tilde{k}(s)$ is the curvature of $W_1(s)$, we have
$$
\frac{d^2W_1(s)}{ds^2}=W_1(s)+\tilde{k}(s)W_2(s), \ \
\frac{dW_2(s)}{ds}=-\tilde{k}(s)\frac{dW_1(s)}{ds}.
$$
So, it follows that
$$
\begin{aligned}
E_2&=\Big(0,(\sqrt{\tfrac{2}{1-C}}-\sqrt{2}\tilde{k}(s))\tfrac{d W_1(s)}{ds}\Big),\\
D_{E_2}N&=\sqrt{\tfrac{2C}{1-C}}\tfrac{\partial N}{\partial s}=
(0,(\sqrt{1-C}-\tilde{k}(s))\tfrac{d W_1(s)}{d s})=\tfrac{\sqrt{1-C}-\tilde{k}(s)}{\sqrt{\tfrac{2}{1-C}}-\sqrt{2}\tilde{k}(s)}E_2,
\end{aligned}
$$
which implies that $b_4=-\langle D_{E_2}N,E_2\rangle=\frac{-\sqrt{1-C}+\tilde{k}(s)}{\sqrt{\frac{2}{1-C}}-\sqrt{2}\tilde{k}(s)}$.
By the assumption that $b_4=-\frac{1}{\sqrt{2}}$, we have $\sqrt{1-C}-\sqrt{\frac{1}{1-C}}=0$.
It is a contradiction. Thus, $W_1(s)$ is a constant vector, we have
$\frac{d W_1(s)}{ds}=0$ and $\frac{d W_2(s)}{ds}\neq0$.
Then, we get
$$
\begin{aligned}
E_2&=(0,\sqrt{2}\tfrac{d W_2(s)}{d s}),\\
D_{E_2}N&=\sqrt{\tfrac{2C}{1-C}}\tfrac{\partial N}{\partial s}=
(0,\tfrac{d W_2(s)}{d s})=\tfrac{1}{\sqrt{2}}E_2,
\end{aligned}
$$
which implies that $b_4=-\langle D_{E_2}N,E_2\rangle=-\frac{1}{\sqrt{2}}$.

In conclusion, we know that $V_1(r)$ is a geodesic of $\mathbb{H}^2$ and $V_2(r)$ is a unit normal
vector field of $V_1(r)\hookrightarrow\mathbb{H}^2$;
$W_1(s)$ is a constant vector in $\mathbb{H}^2$ and $W_2(s)\in T^1_{W_1(s)}\mathbb{H}^2$.
By taking re-parameterizations of $r$ and $s$, and up to an isometry of $\mathbb{H}^2\times \mathbb{H}^2$, we can assume that
$$
\begin{aligned}
V_1(r)&=(\cosh(r),\sinh(r),0),\ \ V_2(r)=(0,0,1),\\
W_1(s)&=(1,0,0), \ \ W_2(s)=(0,\cos(s),\sin(s)).
\end{aligned}
$$
Thus, combining \eqref{eqn:pq}, $C=\frac{1}{\cosh(\sqrt{2}t)}$ and $t<0$, we exactly get the immersion \eqref{eqn:jr}.
%Conversely, it can be directly checked that hypersurface constructed by \eqref{eqn:jr}
%has constant sectional curvature.
\end{proof}


%\begin{remark}\label{rem:5.1}
%The hypersurface constructed by \eqref{eqn:jr} can be regarded as taking
%$$
%f_1(t)=\ln\Big(\sqrt{2}\cosh(\frac{t}{\sqrt{2}})+\sqrt{\cosh(\sqrt{2}t)}\Big),\ \
%f_2(t)={\rm arcsinh}\Big(-\sqrt{2}\sinh(\frac{t}{\sqrt{2}})\Big),\ \ t<0,
%$$
%and
%$$
%\gamma(r)=(\cosh(r),\sinh(r),0), \ \ N(r)=(0,0,1)
%$$
%in Example \ref{exam:4.1ss1}.
%\end{remark}


%For the Case of ``$C<0$, $b_1=-\frac{1}{\sqrt{2}},\ b_4=\frac{1}{\sqrt{2}},\ b_6=\frac{C}{\sqrt{2}}$''.
%It can obtain a similar theorem as Theorem \ref{thm:4.1abc}.

%I think that hypersurface of $\mathbb{H}^2\times \mathbb{H}^2$ with constant sectional curvature $-\frac{1}{2}$
%and ``$C>0$, $b_1=\frac{1}{\sqrt{2}},\ b_4=-\frac{1}{\sqrt{2}},\ b_6=-\frac{C}{\sqrt{2}}$'' is
%a hypersurface constructed in Example \ref{exam:4.1} with specific $f_1(t),f_2(t)$ and curves $\gamma(r)$ and $\tilde{\gamma}(s)$.



In the following, we discuss the hypersurface of $\mathbb{H}^2\times \mathbb{H}^2$ with constant sectional curvature $-\frac{1}{2}$ and $C=0$.
By Lemma \ref{lemma:2.3}, {\bf Case I-ii} of Proposition \ref{prop:4.5} and $C=0$, we have the following lemma.

\begin{lemma}\label{lemma:4.7}
Let $M$ be a hypersurface of $\mathbb{H}^2\times \mathbb{H}^2$ with constant sectional curvature $\kappa=-\frac{1}{2}$ and $C=0$.
Then, under the frame $\{E_1,E_2,E_3\}$ defined by \eqref{eqn:2.5},
the connections satisfy
\begin{equation}\label{eqn:2.25a11}
\begin{aligned}
&A E_1=b_1 E_1+b_2 E_2,\ A E_2=b_2 E_1+b_4 E_2,\ A E_3=0,\\		
&\nabla_{E_1}E_1=b_1E_3,\
\nabla_{E_1}E_2=-b_2E_3,\ \nabla_{E_1}E_3=-b_1E_1
+b_2E_2,\\
&\nabla_{E_2} E_1=b_2E_3,\
\nabla_{E_2} E_2=-b_4E_3,\ \nabla_{E_2}E_3=-b_2E_1
+b_4E_2,\\
&\nabla_{E_3}E_1=\nabla_{E_3}E_2=\nabla_{E_3}E_3=0.
\end{aligned}
\end{equation}
%Moreover, all the Gauss and Codazzi equations are equivalent to the following
%five independent equations:
Moreover, the following six equations hold:
\begin{equation}\label{eqn:4.26}
E_3b_1=-\tfrac{1}{2}+b_1^{2}-b_2^2,
\end{equation}
%
\begin{equation}\label{eqn:4.27}
E_3b_2=b_2\big(b_1-b_4\big),
\end{equation}
%
\begin{equation}\label{eqn:4.28}
E_3b_4=\tfrac{1}{2}+b_2^{2}-b_4^2,
\end{equation}
%
\begin{equation}\label{eqn:4.29}
E_1b_2-E_2b_1=0,	
\end{equation}
%
\begin{equation}\label{eqn:4.30}
E_1b_4-E_2b_2=0,
\end{equation}
%
\begin{equation}\label{eqn:4.31asas}
b_1b_4-b_2^2=-\tfrac{1}{2}.
\end{equation}
%Constant sectional curvature means that $b_1b_4-b_2^2=-\frac{1}{2}$.
\end{lemma}

Then, by using \eqref{eqn:4.26}--\eqref{eqn:4.28} and \eqref{eqn:4.31asas},
we have the following lemma.

\begin{lemma}\label{lemma:4.1aaa1}
Let $M$ be a hypersurface of $\mathbb{H}^2\times \mathbb{H}^2$ with constant sectional curvature $\kappa=-\frac{1}{2}$ and $C=0$.
Then, under the frame $\{E_1,E_2,E_3\}$ defined by \eqref{eqn:2.5}, we have the following six cases
of $b_1,b_2$ and $b_4$ on $M$:

\noindent{\bf Case 1}:
$$
b_1=-\tfrac{1}{\sqrt{2}},\ \ b_2=0,\ \
b_4=\tfrac{1}{\sqrt{2}}.
$$


\noindent{\bf Case 2}:
$$
b_1=-\tfrac{1}{\sqrt{2}}\tanh(\tfrac{t+k}{\sqrt{2}}),\ \ b_2=0,\ \
b_4=\tfrac{1}{\sqrt{2}}\coth(\tfrac{t+k}{\sqrt{2}}).
$$

\noindent{\bf Case 3}:
\begin{equation}\label{eq:jj1}
\begin{aligned}
&b_1=-\tfrac{1}{\sqrt{2}}+b_2,\ b_4=\tfrac{1}{\sqrt{2}}+b_2,\ b_2=e^{-\sqrt{2}(t+h_1)}.\\
\end{aligned}	
\end{equation}


\noindent{\bf Case 4}:
\begin{equation}\label{eq:jj2}
\begin{aligned}
&b_1=-\tfrac{1}{\sqrt{2}}-b_2,\ b_4=\tfrac{1}{\sqrt{2}}-b_2,\ b_2=e^{-\sqrt{2}(t+h_2)}.\\
\end{aligned}	
\end{equation}

\noindent{\bf Case 5}:
\begin{equation}\label{eq:jj3}
\begin{aligned}
b_1&=\tfrac{1}{\sqrt{2}}{\rm sech}(\sqrt{2}(t+l_1))\Big(\sinh (h_3)-\sinh(\sqrt{2}(t+l_1))\Big),\\
b_2&=\tfrac{1}{\sqrt{2}}{\rm sech}(\sqrt{2}(t+l_1))\cosh (h_3) ,\\
b_4&=\tfrac{1}{\sqrt{2}}{\rm sech}(\sqrt{2}(t+l_1))\Big(\sinh (h_3)+\sinh(\sqrt{2}(t+l_1))\Big).
\end{aligned}	
\end{equation}

%\noindent{\bf Case 6 $\Leftrightarrow$ Case 5}:
%\begin{equation}\label{eq:jj3}
%\begin{aligned}
%b_1&=\frac{1}{\sqrt{2}}{\rm sech}(\sqrt{2}(t+h))\Big(\sinh k-\sinh(\sqrt{2}(t+h))\Big),\\
%b_2&=-\frac{1}{\sqrt{2}}{\rm sech}(\sqrt{2}(t+h))\cosh k ,\\
%b_4&=\frac{1}{\sqrt{2}}{\rm sech}(\sqrt{2}(t+h))\Big(\sinh k+\sinh(\sqrt{2}(t+h))\Big).
%\end{aligned}	
%\end{equation}

%b_1=\tfrac{1}{2}{\rm csch}(\sqrt{2}(t+h))\Big(c_3-\sqrt{2}\cosh(\sqrt{2}(t+h))\Big),\
%b_4=\tfrac{1}{2}{\rm csch}(\sqrt{2}(t+h))\Big(c_3+\sqrt{2}\cosh(\sqrt{2}(t+h))\Big)

\noindent{\bf Case 6}: %\noindent{\bf Case 7}:
\begin{equation}\label{eq:jj4}
\begin{aligned}
b_1&=\tfrac{1}{\sqrt{2}}{\rm csch}(\sqrt{2}(t+l_2))\Big(\cosh(h_4)-\cosh(\sqrt{2}(t+l_2))\Big),\\
b_2&=\tfrac{1}{\sqrt{2}}{\rm csch}(\sqrt{2}(t+l_2))\sinh(h_4),\\
b_4&=\tfrac{1}{\sqrt{2}}{\rm csch}(\sqrt{2}(t+l_2))\Big(\cosh(h_4)+\cosh(\sqrt{2}(t+l_2))\Big).
\end{aligned}	
\end{equation}
Here, function $t$ satisfies $E_3t=1$ and functions $h_1,h_2,h_3,h_4, k, l_1, l_2:M\rightarrow \mathbb{R}$
satisfy $E_3h_1=E_3h_2=E_3h_3=E_3h_4=E_3k=E_3l_1=E_3l_2=0$.
\end{lemma}
%
\begin{proof}
Assume that the hypersurface $M$ has constant sectional curvature $-\frac{1}{2}$ and $C=0$.
We will discuss two cases depending on the function $b_2$ on $M$.

{\bf Case I}: $b_2$ is constant.

When $b_2$ is constant on $M$, by \eqref{eqn:4.27}, we have $b_2(b_1-b_4)=0$.
If $b_2$ is non-zero constant, then it holds $b_1=b_4$ and $b_1^2-b_2^2=-\frac{1}{2}$. Now,
$$
E_3b_1=-\tfrac{1}{2}+b_1^{2}-b_2^2=-1,\ \ E_3b_4=\tfrac{1}{2}+b_2^{2}-b_4^2=1,
$$
which contradicts with $b_1=b_4$.

If $b_2=0$, then $b_1$ and $b_4$ satisfy $E_3b_1=-\tfrac{1}{2}+b_1^2$ and $E_3b_4=\tfrac{1}{2}-b_4^2$.
By directly integrating these equations along the integral curves of $E_3$,
and combining $b_1b_4=-\tfrac{1}{2}$, up to a sign of normal vector field,
we have
$$
\begin{aligned}
(b_1,b_4)=(-\tfrac{1}{\sqrt{2}},\tfrac{1}{\sqrt{2}}),\ &{\rm or}\
(b_1,b_4)=(-\tfrac{1}{\sqrt{2}}\tanh(\tfrac{t+k}{\sqrt{2}}),
\tfrac{1}{\sqrt{2}}\coth(\tfrac{t+k}{\sqrt{2}})),\\
&  {\rm or}\ (b_1,b_4)=(-\tfrac{1}{\sqrt{2}}\coth(\tfrac{t+k}{\sqrt{2}}),
\tfrac{1}{\sqrt{2}}\tanh(\tfrac{t+k}{\sqrt{2}})),
\end{aligned}
$$
where $k$ is a function on $M$ satisfying $E_3k=0$.
In fact, if $(b_1,b_4)=(-\tfrac{1}{\sqrt{2}}\coth(\tfrac{t+k}{\sqrt{2}}),
\tfrac{1}{\sqrt{2}}\tanh(\tfrac{t+k}{\sqrt{2}}))$, then by using the isometry
$$
\mathbb{F}:\mathbb{H}^2\times \mathbb{H}^2\rightarrow\mathbb{H}^2\times \mathbb{H}^2,\ \
\mathbb{F}(p,q)=(q,p),
$$
we can have another hypersurface $\mathbb{F}(M)$ of $\mathbb{H}^2\times \mathbb{H}^2$
with constant section curvature $-\frac{1}{2}$ and $C=0$. The unit normal vector field
of $\mathbb{F}(M)$ is $\tilde{N}=d\mathbb{F}(N)$.
Now, define the frame $\{\tilde{E}_1,\tilde{E}_2,\tilde{E}_3\}$ as \eqref{eqn:2.5} on $\mathbb{F}(M)$,
then the corresponding components of the shape operator on $\mathbb{F}(M)$ are given by:
$$
\tilde{b}_1=b_4,\ \tilde{b}_2=b_2=0,\ \tilde{b}_4=b_1.
$$
Thus, if $b_2=0$, up to an isometry of $\mathbb{H}^2\times \mathbb{H}^2$
or up to a sign of normal vector field, we obtain {\bf Case 1} and {\bf Case 2}.

\vskip 2mm

{\bf Case II}: $b_2$ is not constant.

%Up to a sign of $N$, we always assume that $b_2>0$ in local.
When $b_2$ is not constant on $M$, let $F=b_1-b_4$, we get
$$
E_3F=E_3(b_1-b_4)=-2+(b_1-b_4)^2=-2+F^2.
$$
This equation can be directly integrated along the integral curves of $E_3$
to obtain
\begin{equation}\label{eqn:5.1a}
b_1-b_4=\pm\sqrt{2},\ {\rm or}\ b_1-b_4=-\sqrt{2}\tanh(\sqrt{2}(t+l)),
\ {\rm or}\ b_1-b_4=-\sqrt{2}\coth(\sqrt{2}(t+l)),
\end{equation}
where $l$ is a function on $M$ satisfying $E_3l=0$.
Note that, for the case of $b_1-b_4=\pm\sqrt{2}$, up to the isometry $\mathbb{F}$ of $\mathbb{H}^2\times \mathbb{H}^2$, and up to a sign of $N$, %we always assume that $b_2>0$ in local.
we can always assume that $b_1-b_4=-\sqrt{2}$ and $b_2>0$.

Then by $E_3b_2=b_2\big(b_1-b_4\big)=b_2F$, with the use of \eqref{eqn:5.1a}, we integrate along the integral curves of $E_3$
to obtain
\begin{equation}\label{eqn:5.2a}
b_2=e^{-\sqrt{2}(t+h)}, \ {\rm or}\ b_2=h\, {\rm sech}(\sqrt{2}(t+l)),
\ {\rm or}\ b_2=h\, {\rm csch}(\sqrt{2}(t+l)),
\end{equation}
where $h$ is a function on $M$ satisfying $E_3h=0$.

Now, let $\tilde{F}=b_1+b_4$, by \eqref{eqn:4.26} and \eqref{eqn:4.28}, we get
$$
E_3\tilde{F}=E_3(b_1+b_4)=(b_1+b_4)(b_1-b_4)=F\tilde{F}.
$$
Using \eqref{eqn:5.1a}, we integrate along the integral curves of $E_3$
to obtain
$$
b_1+b_4=c_3e^{-\sqrt{2}t}, \ {\rm or}\ \ b_1+b_4=c_3{\rm sech}(\sqrt{2}(t+l)),
\ {\rm or}\ \ b_1+b_4=c_3{\rm csch}(\sqrt{2}(t+l)),
$$
where $c_3$ is a function on $M$ satisfying $E_3c_3=0$. Thus, we have the following three situations:
$$
\begin{aligned}
&b_1=-\tfrac{1}{\sqrt{2}}+\tfrac{c_3}{2}e^{-\sqrt{2}t},\ b_4=\tfrac{1}{\sqrt{2}}+\tfrac{c_3}{2}e^{-\sqrt{2}t};\\
&b_1=\tfrac{1}{2}{\rm sech}(\sqrt{2}(t+l))\Big(c_3-\sqrt{2}\sinh(\sqrt{2}(t+l))\Big),\
b_4=\tfrac{1}{2}{\rm sech}(\sqrt{2}(t+l))\Big(c_3+\sqrt{2}\sinh(\sqrt{2}(t+l))\Big);\\
&b_1=\tfrac{1}{2}{\rm csch}(\sqrt{2}(t+l))\Big(c_3-\sqrt{2}\cosh(\sqrt{2}(t+l))\Big),\
b_4=\tfrac{1}{2}{\rm csch}(\sqrt{2}(t+l))\Big(c_3+\sqrt{2}\cosh(\sqrt{2}(t+l))\Big).
\end{aligned}
$$
From the fact that $b_1b_4-b_2^2=-\frac{1}{2}$ and using \eqref{eqn:5.2a},
corresponding to the above three situations, we have
$$
\begin{aligned}
&c_3^2=4e^{-2\sqrt{2}h};\ \ \ {\rm or}\ \ 4h^2-c_3^2=2;\ \ \ {\rm or}\ \ 4h^2-c_3^2=-2.
\end{aligned}
$$

For the first situation that $c_3^2=4e^{-2\sqrt{2}h}$, we have $c_3=\pm2e^{-\sqrt{2}h}$.
Now, we have two subcases as follows:
\begin{equation*}
\left\{
\begin{aligned}
%&(1)\ b_1=\tfrac{1}{\sqrt{2}}+b_2,\ b_2=e^{\sqrt{2}(t+h)},\ b_4=-\tfrac{1}{\sqrt{2}}+b_2.\\
%&(2)\ b_1=\tfrac{1}{\sqrt{2}}-b_2,\ b_2=e^{\sqrt{2}(t+h)},\ b_4=-\tfrac{1}{\sqrt{2}}-b_2.\\
&(1)\ b_1=-\tfrac{1}{\sqrt{2}}+b_2,\ b_2=e^{-\sqrt{2}(t+h)},\ b_4=\tfrac{1}{\sqrt{2}}+b_2.\\
&(2)\ b_1=-\tfrac{1}{\sqrt{2}}-b_2,\ b_2=e^{-\sqrt{2}(t+h)},\ b_4=\tfrac{1}{\sqrt{2}}-b_2.
\end{aligned}\right.
\end{equation*}
%Note that, up to the isometry $\mathbb{F}$ of $\mathbb{H}^2\times \mathbb{H}^2$, subcase (1) is
%isometric to subcase (3), and subcase (2) is isometric to subcase (4).
Thus, we have {\bf Case 3} and {\bf Case 4}.

For the second situation that $4h^2-c_3^2=2$, up to an isometry of $\mathbb{H}^2\times \mathbb{H}^2$,
we assume $h=\frac{1}{\sqrt{2}}\cosh h_3$ and $c_3=\sqrt{2}\sinh h_3$, where $h_3$ is a function on $M$ satisfying $E_3h_3=0$. Thus we obtain {\bf Case 5}.

For the third situation that $4h^2-c_3^2=-2$, up to an isometry of $\mathbb{H}^2\times \mathbb{H}^2$,
we assume $h=\frac{1}{\sqrt{2}}\sinh h_4$ and $c_3=\sqrt{2}\cosh h_4$, where $h_4$ is a function on $M$ satisfying $E_3h_4=0$. Thus we obtain {\bf Case 6}.
\end{proof}

For {\bf Case 1} and {\bf Case 2} of Lemma \ref{lemma:4.1aaa1}, we have the following proposition.

\begin{proposition}\label{prop:4.1a1}
Let $M$ be a hypersurface of $\mathbb{H}^2\times \mathbb{H}^2$ with constant sectional curvature $-\frac{1}{2}$ and $C=0$.
Under the frame $\{E_1,E_2,E_3\}$ defined by \eqref{eqn:2.5}, if $b_1,b_2$ and $b_4$ are
given by {\bf Case 1} or {\bf Case 2} of Lemma \ref{lemma:4.1aaa1}, then
$M$ is an open part of $M_{k_1,\frac{1}{k_1}}^{1/2}$ given in Example \ref{exam:3.4fg1} for some constant $0<k_1\leq1$.
More specifically, {\bf Case 1} corresponds to $M_{1,1}^{1/2}$, and
{\bf Case 2} corresponds to $M_{k_1,\frac{1}{k_1}}^{1/2}$ for some constant $0<k_1<1$.
\end{proposition}
%
\begin{proof}
Recall that the functions $b_2=0$ in {\bf Case 1} and {\bf Case 2} of Lemma \ref{lemma:4.1aaa1}, and we have classified the hypersurfaces of $\mathbb{H}^2\times \mathbb{H}^2$ with constant $C$
and constant $b_2$ in Proposition \ref{prop:3.61}. Then, Proposition \ref{prop:4.1a1} can be obtained directly by using (2) of Lemma \ref{lemma:3.5fg12}.
%checking that hypersurface $M_{\kappa,\tilde{\kappa}}^c$
%has constant sectional curvature if and only if $c=\frac{1}{2}$, and $\kappa(r)$ and
%$\tilde{\kappa}(s)$ are constant which satisfies $\kappa(r) \tilde{\kappa}(s)=1$.
\end{proof}



Now, we begin to deal with {\bf Case 3} and {\bf Case 4} of Lemma \ref{lemma:4.1aaa1}.


\begin{theorem}\label{thm:5.10a1}
Let $M$ be a hypersurface of $\mathbb{H}^2\times \mathbb{H}^2$ with constant sectional curvature $-\frac{1}{2}$ and $C=0$.
Under the frame $\{E_1,E_2,E_3\}$ defined by \eqref{eqn:2.5}, if $b_1,b_2$ and $b_4$ are
given by {\bf Case 3} of Lemma \ref{lemma:4.1aaa1}, then $M$ admits local coordinates $(u,v,t)$
such that the functions $b_1,b_2,b_4$ are given by \eqref{eqn:3.1ff1}.



Conversely, let $\mathcal{D}$ be an open subset in $\mathbb{R}^2$, $I$ be an open subset in $\mathbb{R}$ and $\beta_1:I\to\mathbb{R},\  v\mapsto \beta_1(v)\in\mathbb{R}$ be any function. We define a metric $g$ on $\Omega=\{(u,v,t)\in \mathcal{D}\times I\}\subseteq \mathbb{R}^3$ by
\eqref{eqn:3.1ff3}, and a $(1,1)$-tensor field $A$ on $T\Omega$ by
\eqref{eqn:3.1ff4}.
Let $\nu$ be a vector bundle over  $\Omega$ of rank $1$ with metric $\tilde{g}$, $\nabla ^{\perp} $ a connection on $\nu$ compatible with the metric $\tilde{g}$, and $N$ a unit vector field in $\nu$.   We define a $(1,1)$-tensor field $\tilde{P}$ on $T\Omega\otimes \nu$ by
\eqref{eqn:3.1ff5}.
Then the integrability conditions are satisfied on $\Omega$ and hence on any simply
connected subset $\Omega_1$ of $\Omega$, up to an isometry of  $\mathbb{H}^2\times \mathbb{H}^2$, there is a unique isometric immersion $\Psi$ from $(\Omega_1,~g)$ into $\mathbb{H}^2\times \mathbb{H}^2$ such that $A$ is the shape operator of  $\Omega_1$ into $\mathbb{H}^2\times \mathbb{H}^2$, $\nu$ is isometric to the normal bundle of $\Psi(\Omega_1)$ in
$\mathbb{H}^2\times \mathbb{H}^2$ by an isomorphism $\tilde{\Psi}:\nu\to T^{\perp}\Psi(\Omega_1)$,
and for any $Y\in T\Omega_1$, we have
\begin{equation}
	P(\Psi_{*}Y)=\Psi_{*}((\tilde{P}Y)^{T})+\tilde{\Psi}((\tilde{P}Y)^{\perp}),~
	P(\tilde{\Psi}N)=\Psi_*(\partial_t),
\end{equation}
where $P$ is the product structure of $\mathbb{H}^2\times \mathbb{H}^2$, $(\tilde{P}Y)^{T}$ and $(\tilde{P}Y)^{\perp}$ denote the projections of  $\tilde{P}Y$ onto $T\Omega_1$ and $\nu$, respectively.
Moreover, $\Psi(\Omega_1)$ has constant sectional curvature $-\frac{1}{2}$ and product angle function $C=0$.
\end{theorem}
%
\begin{proof}
For {\bf Case 3} of Lemma \ref{lemma:4.1aaa1}, we have $b_1=-\frac{1}{\sqrt{2}}+b_2,\ b_4=\frac{1}{\sqrt{2}}+b_2$. By \eqref{eqn:4.29}, we have $E_1b_2=E_2b_2$. Let $f=-\frac{1}{\sqrt{2}}\ln b_2$ and $d_1=E_1f$,
then we have $E_2b_2=-\sqrt{2}b_2d_1$.

Now, we calculate $0=([E_i,E_j]-(\nabla_{E_i}{E_j}-\nabla_{E_j}{E_i}))b_2$, $1\leq i,j\leq3$, and assume that $E_1d_1+E_2d_1=d_2$, then we have
\begin{equation}\label{eqn:Eid1}
E_1d_1=-b_2+\frac{d_2}{2},\ \ E_2d_1=b_2+\frac{d_2}{2},\ \ E_3d_1=-\frac{d_1}{\sqrt{2}}.
\end{equation}

Next, we calculate $0=([E_i,E_j]-(\nabla_{E_i}{E_j}-\nabla_{E_j}{E_i}))d_1$, $1\leq i,j\leq3$, and have
\begin{equation}\label{eqn:Eid2}
E_2d_2=-6\sqrt{2}b_2d_1+E_1d_2,\ \ E_3d_2=-4b_2^2-\sqrt{2}d_2.
\end{equation}

Then, we define the new frame $\{X_1,X_2,X_3\}$ on $M$ as follows:
$$
\begin{aligned}
X_1&=\frac{-1}{2b_2}(E_1-E_2),\ X_3=\frac{-d_1}{2\sqrt{2}b_2}(E_1-E_2)+E_3,\\
X_2&=\frac{-2b_2-\sqrt{2}d_1^2-d_2}{2b_2^{\frac{3}{2}}}E_1+\frac{-2b_2+\sqrt{2}d_1^2+d_2}
{2b_2^{\frac{3}{2}}}E_2+\frac{2d_1}{\sqrt{b_2}}E_3.
\end{aligned}
$$
One can check that frame $\{X_1,X_2,X_3\}$ satisfies $[X_i,X_j]=0$. By using $E_1f=E_2f=d_1,E_3f=1$ and
\eqref{eqn:Eid1}, we can have
$$
X_1d_1=1,\ X_2d_1=X_3d_1=0,\ X_1f=X_2f=0,\ X_3f=1.
$$
Due that $[X_i,X_j]=0$, then there exists a local coordinates $\{u,v,t\}$, such that
$$
\partial_u=X_1,\ \partial_v=X_2,\ \partial_t=X_3,\ f=t,\ b_2=e^{-\sqrt{2}t},\ d_1=u.
$$
So the functions $b_1,b_2,b_4$ are given by
%
\begin{equation}\label{eqn:3.1ff1}
b_1=-\frac{1}{\sqrt{2}}+e^{-\sqrt{2}t},\ b_2=e^{-\sqrt{2}t},\ b_4=\frac{1}{\sqrt{2}}+e^{-\sqrt{2}t}.
\end{equation}
It follows that
\begin{equation}\label{eqn:d1d2}
\left(
  \begin{array}{c}
    E_1 \\
    E_2 \\
    E_3 \\
  \end{array}
\right)
=\left(\begin{array}{cccc}
-b_2+\frac{d_2}{2} & \frac{-\sqrt{b_2}}{2} & d_1 \\
b_2+\frac{d_2}{2} & \frac{-\sqrt{b_2}}{2} & d_1 \\
\frac{-d_1}{\sqrt{2}} & 0 & 1
\end{array}\right)
\left(
  \begin{array}{c}
    \partial_u \\
    \partial_v \\
    \partial_t \\
  \end{array}
\right).
\end{equation}

Now, by using \eqref{eqn:Eid2} and \eqref{eqn:d1d2}, we have
\begin{equation}\label{eqn:dud2}
\frac{\partial d_2}{\partial u}=-3\sqrt{2}u,
\end{equation}
\begin{equation}\label{eqn:dtd2}
-\frac{u}{\sqrt{2}}\frac{\partial d_2}{\partial u}+\frac{\partial d_2}{\partial t}=
-4e^{-2\sqrt{2}t}-\sqrt{2}d_2.
\end{equation}
By \eqref{eqn:dud2}, we have $d_2=\frac{-3}{\sqrt{2}}u^2+\alpha(v,t)$,
where $\alpha(v,t)$ is a function depending on variables $v$ and $t$.
By \eqref{eqn:dtd2}, we have $\frac{\partial \alpha(v,t)}{\partial t}=-\sqrt{2}\alpha(v,t)-4e^{-2\sqrt{2}t}$, which deduces that
$\alpha(v,t)=e^{-2\sqrt{2}t}(2\sqrt{2}+e^{\sqrt{2}t}\beta_1(v))$.
Here, $\beta_1(v)$ is a function depending on variable $v$.
Thus, we have $d_2=\frac{-3}{\sqrt{2}}u^2+e^{-2\sqrt{2}t}(2\sqrt{2}+e^{\sqrt{2}t}\beta_1(v))$.
Finally, we can check that the compatibility conditions of  $b_1,b_2,b_4,d_1,d_2,\beta_1(v)$  with respect to the Lie brackets $[E_i,E_j]$ are automatically satisfied.


Conversely, we assume that the metric $g$ on $\Omega$ is defined by
%
\begin{equation}\label{eqn:3.1ff3}
\begin{aligned}
g=&\tfrac{1}{4}(16 e^{-\sqrt{2}t}+e^{3 \sqrt{2}t}u^4+8\sqrt{2}\beta_1(v)
-2\sqrt{2}e^{2\sqrt{2}t}u^2\beta_1(v)+2e^{\sqrt{2}t}(4+4u^2+\beta_1^2(v)))dv^2\\
&+\tfrac{1}{2}e^{2\sqrt{2}t}du^2+(1+\tfrac{1}{4}e^{2\sqrt{2}t}u^2)dt^2+\tfrac{1}{2}e^{\tfrac{t}{\sqrt{2}}}
(\sqrt{2}(4-e^{2\sqrt{2}t}u^2)+2e^{\sqrt{2}t}\beta_1(v))dudv\\
&+\tfrac{e^{2\sqrt{2}t}u}{\sqrt{2}}dudt
+\tfrac{1}{2}e^{\tfrac{t}{\sqrt{2}}}u(12-e^{2\sqrt{2}t}u^2+\sqrt{2}e^{\sqrt{2}t}\beta_1(v))dvdt,
\end{aligned}
\end{equation}
the $(1,1)$-tensor $A$ on $T\Omega$ is defined by
\begin{equation}\label{eqn:3.1ff4}
\begin{aligned}
A\partial_u=&(e^{-\sqrt{2}t}-\tfrac{3}{4}e^{\sqrt{2}t}u^2+\tfrac{\beta_1(v)}{2\sqrt{2}})\partial_u
-\tfrac{e^{\tfrac{t}{\sqrt{2}}}}{2\sqrt{2}}\partial_v+\tfrac{e^{\sqrt{2}t}u}{\sqrt{2}}\partial_t,\\
A\partial_v=&\tfrac{1}{8}e^{-\tfrac{5t}{\sqrt{2}}}(-16\sqrt{2}+3\sqrt{2}e^{4\sqrt{2}t}u^4-8e^{3\sqrt{2}t} u^2\beta_1(v)+2\sqrt{2}e^{2\sqrt{2}t}(-4+4u^2+\beta_1^2(v)))\partial_u\\
&+(e^{-\sqrt{2}t}+\tfrac{1}{4}e^{\sqrt{2}t}u^2-\tfrac{\beta_1(v)}{2\sqrt{2}})\partial_v
-\tfrac{1}{2}e^{-\tfrac{t}{\sqrt{2}}}u(4+e^{2\sqrt{2}t}u^2-\sqrt{2}e^{\sqrt{2}t}\beta_1(v))\partial_t,\\
A\partial_t=&\tfrac{1}{8}u(\sqrt{2}e^{-\sqrt{2}t}(4-3e^{2\sqrt{2}t}u^2)+2\beta_1(v))\partial_u
-\tfrac{e^{\tfrac{t}{\sqrt{2}}}u}{4}\partial_v+\tfrac{1}{2}e^{\sqrt{2}t}u^2\partial_t,
\end{aligned}
\end{equation}
and the $(1,1)$-tensor $\tilde{P}$ on $T\Omega\otimes \nu$ is defined by
\begin{equation}\label{eqn:3.1ff5}
\begin{aligned}
\tilde{P}\partial_u=&\tfrac{1}{4}(\sqrt{2}e^{-\sqrt{2}t}(-4+3e^{2\sqrt{2}t}u^2)-2 \beta_1(v))\partial_u+\tfrac{e^{\tfrac{t}{\sqrt{2}}}}{2}\partial_v
-e^{\sqrt{2}t}u\partial_t,\\
\tilde{P}\partial_v=&\tfrac{1}{4}e^{-\tfrac{5t}{\sqrt{2}}}(-16-3e^{4\sqrt{2}t}u^4-8\sqrt{2}e^{\sqrt{2}t} \beta(v)+4\sqrt{2}e^{3\sqrt{2}t}u^2\beta_1(v)+2e^{2\sqrt{2}t}(4+8u^2-\beta_1^2(v)))\partial_u\\
&+\tfrac{1}{4}(\sqrt{2}e^{-\sqrt{2}t}(4-e^{2\sqrt{2}t}u^2)+2\beta_1(v))\partial_v
+\tfrac{1}{2}e^{-\tfrac{t}{\sqrt{2}}}u(\sqrt{2}(-4+e^{2\sqrt{2}t}u^2)-2e^{\sqrt{2}t} \beta_1(v))\partial_t\\
&+2ue^{\tfrac{t}{\sqrt{2}}}N,\\
\tilde{P}\partial_t=&\tfrac{1}{4}u(-4e^{-\sqrt{2}t}+3e^{\sqrt{2}t}u^2-\sqrt{2} \beta_1(v))\partial_u+\tfrac{ue^{\frac{t}{\sqrt{2}}}}{2\sqrt{2}}\partial_v
-\tfrac{e^{\sqrt{2}t}u^2}{\sqrt{2}}\partial_t+N,\\
\tilde{P}N=&-\tfrac{u}{\sqrt{2}}\partial_u+\partial_t.\\	
\end{aligned}
\end{equation}
Then it can be  checked  that all the integrability conditions are satisfied.
Therefore, the conclusions can be obtained by applying  the existence and uniqueness theorems in  \cite{K,L-T-V} directly.
\end{proof}


\begin{theorem}\label{thm:5.11a1}
Let $M$ be a hypersurface of $\mathbb{H}^2\times \mathbb{H}^2$ with constant sectional curvature $-\frac{1}{2}$ and $C=0$.
Under the frame $\{E_1,E_2,E_3\}$ defined by \eqref{eqn:2.5}, if $b_1,b_2$ and $b_4$ are
given by {\bf Case 4} of Lemma \ref{lemma:4.1aaa1}, then $M$ admits local coordinates $(u,v,t)$
such that the functions $b_1,b_2,b_4$ are given by \eqref{eqn:3.1ffff1}.


Conversely, let $\mathcal{D}$ be an open subset in $\mathbb{R}^2$, $I$ be an open subset in $\mathbb{R}$ and $\beta_2:I\to\mathbb{R},\ v\mapsto \beta_2(v)\in\mathbb{R}$ be any function. We define a metric $g$ on  $\Omega=\{(u,v,t)\in \mathcal{D}\times I\}\subseteq  \mathbb{R}^3$ by
\eqref{eqn:3.1ffff3}
and a $(1,1)$-tensor field $A$ on $T\Omega$ by
\eqref{eqn:3.1ffff4}.
Let $\nu$ be a vector bundle over  $\Omega$ of rank $1$ with metric $\tilde{g}$, $\nabla ^{\perp} $ a connection on $\nu$ compatible with the metric $\tilde{g}$, and $N$ a unit vector field in $\nu$.   We define a $(1,1)$-tensor field $\tilde{P}$ on $T\Omega\otimes \nu$ by
\eqref{eqn:3.1ffff5}.

Then the integrability conditions are satisfied on $\Omega$ and hence on any simply
connected subset $\Omega_1$ of $\Omega$, up to an isometry of  $\mathbb{H}^2\times \mathbb{H}^2$, there is a unique isometric immersion $\Psi$ from $(\Omega_1,~g)$ into $\mathbb{H}^2\times \mathbb{H}^2$ such that $A$ is the shape operator of  $\Omega_1$ into $\mathbb{H}^2\times \mathbb{H}^2$, $\nu$ is isometric to the normal bundle of $\Psi(\Omega_1)$ in
$\mathbb{H}^2\times \mathbb{H}^2$ by an isomorphism $\tilde{\Psi}:\nu\to T^{\perp}\Psi(\Omega_1)$,
and for any $Y\in T\Omega_1$, we have
\begin{equation}
	P(\Psi_{*}Y)=\Psi_{*}((\tilde{P}Y)^{T})+\tilde{\Psi}((\tilde{P}Y)^{\perp}),~
	P(\tilde{\Psi}N)=\Psi_*(\partial_t),
\end{equation}
where $P$ is the product structure of $\mathbb{H}^2\times \mathbb{H}^2$, $(\tilde{P}Y)^{T}$ and $(\tilde{P}Y)^{\perp}$ denote the projections of  $\tilde{P}Y$ onto $T\Omega_1$ and $\nu$, respectively.
Moreover, $\Psi(\Omega_1)$ has constant sectional curvature $-\frac{1}{2}$ and product angle function $C=0$.
\end{theorem}
%
\begin{proof}
For {\bf Case 4}, we have $b_1=-\frac{1}{\sqrt{2}}-b_2,\ b_4=\frac{1}{\sqrt{2}}-b_2$. By \eqref{eqn:4.29}, we get $E_1b_2=-E_2b_2$. Let $f=-\frac{1}{\sqrt{2}}\ln b_2$ and $d_1=E_1f$,
then it holds $E_2b_2=\sqrt{2}b_2d_1$.

Now, we calculate $0=([E_i,E_j]-(\nabla_{E_i}{E_j}-\nabla_{E_j}{E_i}))b_2$, $1\leq i,j\leq3$, and assume that
$E_1d_1-E_2d_1=d_2$, it follows that
\begin{equation}\label{eqn:Eid1ff}
E_1d_1=b_2+\frac{d_2}{2},\ \ E_2d_1=b_2-\frac{d_2}{2},\ \ E_3d_1=-\frac{d_1}{\sqrt{2}}.
\end{equation}

Next, we calculate $0=([E_i,E_j]-(\nabla_{E_i}{E_j}-\nabla_{E_j}{E_i}))d_1$, $1\leq i,j\leq3$, and have
\begin{equation}\label{eqn:Eid2ff}
E_2d_2=-6\sqrt{2}b_2d_1-E_1d_2,\ \ E_3d_2=-4b_2^2-\sqrt{2}d_2.
\end{equation}

Then, we define the new frame $\{X_1,X_2,X_3\}$ on $M$ as follows:
$$
\begin{aligned}
X_1&=\frac{1}{2b_2}(E_1+E_2),\ X_3=\frac{d_1}{2\sqrt{2}b_2}(E_1+E_2)+E_3,\\
X_2&=\frac{2b_2-\sqrt{2}d_1^2-d_2}{2b_2^{\frac{3}{2}}}E_1+\frac{-2b_2-\sqrt{2}d_1^2-d_2}{2b_2^{\frac{3}{2}}}E_2
-\frac{2d_1}{\sqrt{b_2}}E_3.\\
\end{aligned}
$$
One can check that frame $\{X_1,X_2,X_3\}$ satisfies $[X_i,X_j]=0$. By using $E_1f=-E_2f=d_1,E_3f=1$ and
\eqref{eqn:Eid1ff}, we can get
$$
X_1d_1=1,\ X_2d_1=X_3d_1=0,\ X_1f=X_2f=0,\ X_3f=1.
$$
Due that $[X_i,X_j]=0$, there exists a local coordinates $\{u,v,t\}$, such that
$$
\partial_u=X_1,\ \partial_v=X_2,\ \partial_t=X_3,\ f=t,\ b_2=e^{-\sqrt{2}t},\ d_1=u.
$$
So the functions $b_1,b_2,b_4$ are given by
%
\begin{equation}\label{eqn:3.1ffff1}
b_1=-\frac{1}{\sqrt{2}}-e^{-\sqrt{2}t},\ b_2=e^{-\sqrt{2}t},\ b_4=\frac{1}{\sqrt{2}}-e^{-\sqrt{2}t}.
\end{equation}
It follows that
\begin{equation}\label{eqn:d1d2ff}
\left(
  \begin{array}{c}
    E_1 \\
    E_2 \\
    E_3 \\
  \end{array}
\right)
=\left(\begin{array}{cccc}
b_2+\frac{d_2}{2} & \frac{\sqrt{b_2}}{2} & d_1 \\
b_2-\frac{d_2}{2} & \frac{-\sqrt{b_2}}{2} & -d_1 \\
\frac{-d_1}{\sqrt{2}} & 0 & 1
\end{array}\right)
\left(
  \begin{array}{c}
    \partial_u \\
    \partial_v \\
    \partial_t \\
  \end{array}
\right).
\end{equation}

Next, by using \eqref{eqn:Eid2ff} and \eqref{eqn:d1d2ff}, we obtain
\begin{equation}\label{eqn:dud2ff}
\frac{\partial d_2}{\partial u}=-3\sqrt{2}u,
\end{equation}
\begin{equation}\label{eqn:dtd2ff}
-\frac{u}{\sqrt{2}}\frac{\partial d_2}{\partial u}+\frac{\partial d_2}{\partial t}=
-4e^{-2\sqrt{2}t}-\sqrt{2}d_2.
\end{equation}
Now, by \eqref{eqn:dud2ff}, we have $d_2=\frac{-3}{\sqrt{2}}u^2+\alpha(v,t)$,
where $\alpha(v,t)$ is a function depending on variables $v$ and $t$.
By \eqref{eqn:dtd2ff}, we have $\frac{\partial \alpha(v,t)}{\partial t}=-\sqrt{2}\alpha(v,t)-4e^{-2\sqrt{2}t}$, which deduces that
$\alpha(v,t)=e^{-2\sqrt{2}t}(2\sqrt{2}+e^{\sqrt{2}t}\beta_2(v))$.
Here, $\beta_2(v)$ is a function depending on variable $v$.
Thus, we have
$$
d_2=\frac{-3}{\sqrt{2}}u^2+e^{-2\sqrt{2}t}(2\sqrt{2}+e^{\sqrt{2}t}\beta_2(v)).
$$
Finally, we can check that the compatibility conditions of  $b_1,b_2,b_4,d_1,d_2,\beta_2(v)$  with respect to the Lie brackets $[E_i,E_j]$ are automatically satisfied.


Conversely, we assume that the metric $g$ on $\Omega$ is defined by
%
\begin{equation}\label{eqn:3.1ffff3}
\begin{aligned}
g=&\frac{1}{4}(16 e^{-\sqrt{2}t}+e^{3 \sqrt{2}t}u^4+8\sqrt{2}\beta_2(v)
-2\sqrt{2}e^{2\sqrt{2}t}u^2\beta_2(v)+2e^{\sqrt{2}t}(4+4u^2+\beta_2^2(v)))dv^2\\
&+\frac{1}{2}e^{2\sqrt{2}t}du^2+(1+\frac{1}{4}e^{2\sqrt{2}t}u^2)dt^2+\frac{1}{2}e^{\frac{t}{\sqrt{2}}}(\sqrt{2}(-4+e^{2\sqrt{2}t}u^2)-2e^{\sqrt{2}t}\beta_2(v))dudv\\
&+\frac{e^{2\sqrt{2}t}u}{\sqrt{2}}dudt
+\frac{1}{2}e^{\frac{t}{\sqrt{2}}}u(-12+e^{2\sqrt{2}t}u^2-\sqrt{2}e^{\sqrt{2}t}\beta_2(v))dvdt,
\end{aligned}
\end{equation}
the $(1,1)$-tensor $A$ on $T\Omega$ is defined by
\begin{equation}\label{eqn:3.1ffff4}
\begin{aligned}
A\partial_u=&(-e^{-\sqrt{2}t}+\frac{3}{4}e^{\sqrt{2}t}u^2-\frac{\beta_2(v)}{2\sqrt{2}})\partial_u
-\frac{e^{\frac{t}{\sqrt{2}}}}{2\sqrt{2}}\partial_v-\frac{e^{\sqrt{2}t}u}{\sqrt{2}}\partial_t,\\
A\partial_v=&\frac{1}{8}e^{-\frac{5t}{\sqrt{2}}}(-16\sqrt{2}+3\sqrt{2}e^{4\sqrt{2}t}u^4-8e^{3\sqrt{2}t} u^2\beta_2(v)+2\sqrt{2}e^{2\sqrt{2}t}(-4+4u^2+\beta_2^2(v)))\partial_u\\
&+(-e^{-\sqrt{2}t}-\frac{1}{4}e^{\sqrt{2}t}u^2+\frac{\beta_2(v)}{2\sqrt{2}})\partial_v
-\frac{1}{2}e^{-\frac{t}{\sqrt{2}}}u(4+e^{2\sqrt{2}t}u^2-\sqrt{2}e^{\sqrt{2}t}\beta_2(v))\partial_t,\\
A\partial_t=&\frac{1}{8}u(\sqrt{2}e^{-\sqrt{2}t}(-4+3e^{2\sqrt{2}t}u^2)-2\beta_2(v))\partial_u
-\frac{e^{\frac{t}{\sqrt{2}}}u}{4}\partial_v-\frac{1}{2}e^{\sqrt{2}t}u^2\partial_t,
\end{aligned}
\end{equation}
and the $(1,1)$-tensor $\tilde{P}$ on $T\Omega\otimes \nu$ is defined by
\begin{equation}\label{eqn:3.1ffff5}
\begin{aligned}
\tilde{P}\partial_u=&\frac{1}{4}(\sqrt{2}e^{-\sqrt{2}t}(4-3e^{2\sqrt{2}t}u^2)+2 \beta_2(v))\partial_u+\frac{e^{\frac{t}{\sqrt{2}}}}{2}\partial_v
+e^{\sqrt{2}t}u\partial_t,\\
\tilde{P}\partial_v=&\frac{1}{4}e^{-\frac{5t}{\sqrt{2}}}(-16-3e^{4\sqrt{2}t}u^4-8\sqrt{2}e^{\sqrt{2}t} \beta_2(v)+4\sqrt{2}e^{3\sqrt{2}t}u^2\beta_2(v)+2e^{2\sqrt{2}t}(4+8u^2-\beta_2^2(v)))\partial_u\\
&+\frac{1}{4}(\sqrt{2}e^{-\sqrt{2}t}(-4+e^{2\sqrt{2}t}u^2)-2\beta_2(v))\partial_v
+\frac{1}{2}e^{-\frac{t}{\sqrt{2}}}u(\sqrt{2}(-4+e^{2\sqrt{2}t}u^2)-2e^{\sqrt{2}t} \beta_2(v))\partial_t\\
&-2ue^{\frac{t}{\sqrt{2}}}N,\\
\tilde{P}\partial_t=&\frac{1}{4}u(4e^{-\sqrt{2}t}-3e^{\sqrt{2}t}u^2+\sqrt{2} \beta_2(v))\partial_u+\frac{ue^{\frac{t}{\sqrt{2}}}}{2\sqrt{2}}\partial_v+\frac{e^{\sqrt{2}t}u^2}{\sqrt{2}}\partial_t+N,\\
\tilde{P}N=&-\frac{u}{\sqrt{2}}\partial_u+\partial_t.\\	
\end{aligned}
\end{equation}
Then it can be  checked  that all the integrability conditions are satisfied.
Therefore, the conclusions can be obtained by applying  the existence and uniqueness theorems in  \cite{K,L-T-V} directly.
\end{proof}


In the following, we deal with {\bf Case 5} and {\bf Case 6} of Lemma \ref{lemma:4.1aaa1}. By $b_1-b_4\neq\sqrt{2}$, we define
\begin{equation}\label{eqn:bibar}
b_1^*=\frac{b_1+b_4}{2-\sqrt{2}(b_1-b_4)},\ \
b_2^*=\frac{2b_2}{2-\sqrt{2}(b_1-b_4)}.
\end{equation}
Then, combining with $b_1b_4-b_2^2=-\tfrac{1}{2}$, we have
\begin{equation}\label{eqn:bi}
b_1=\frac{1}{\sqrt{2}}-\frac{\sqrt{2}-2b_1^*}{1-2b_1^{*2}+2b_2^{*2}},\
b_2=\frac{2b_2^*}{1-2b_1^{*2}+2b_2^{*2}},\
b_4=-\frac{1}{\sqrt{2}}+\frac{\sqrt{2}+2b_1^*}{1-2b_1^{*2}+2b_2^{*2}}.
\end{equation}

Now, by \eqref{eqn:4.26}--\eqref{eqn:4.30}, we have
\begin{equation}\label{eqn:Eibibar}
\begin{aligned}
&E_3b_1^*=-\sqrt{2}b_1^*,\ \
E_1b_1^*=-E_2b_2^*+\frac{2(\sqrt{2}b_2^*E_1b_2^*+(1-\sqrt{2}b_1^*)E_2b_2^*)}{1-2b_1^{*2}+2b_2^{*2}},\\
&E_3b_2^*=-\sqrt{2}b_2^*,\ \ E_2b_1^*=-E_1b_2^*+\frac{2((1+\sqrt{2}b_1^*)E_1b_2^*-\sqrt{2}b_2^*E_2b_2^*)}{1-2b_1^{*2}+2b_2^{*2}}.
\end{aligned}
\end{equation}
Next, we define
\begin{equation}\label{eqn:Eibar}
\begin{aligned}
&E_1^*=(1+\sqrt{2}b_1^{*}-\sqrt{2}b_2^{*})E_1+(-1+\sqrt{2}b_1^{*}-\sqrt{2}b_2^{*})E_2,\\
&E_2^*=(1+\sqrt{2}b_1^{*}+\sqrt{2}b_2^{*})E_1+(1-\sqrt{2}b_1^{*}-\sqrt{2}b_2^{*})E_2,\ \
E_3^*=E_3.
\end{aligned}
\end{equation}
Then, by \eqref{eqn:2.25a11}, \eqref{eqn:Eibibar} and \eqref{eqn:Eibar}, one can check that
\begin{equation}\label{eqn:LieEibar}
[E_1^*,E_3^*]=\frac{1}{\sqrt{2}}E_1^*,\ \ [E_2^*,E_3^*]=\frac{1}{\sqrt{2}}E_2^*,\ \
[E_1^*,E_2^*]=-8b_2^*E_3^*.
\end{equation}







\begin{theorem}\label{thm:5.12a1}
Let $M$ be a hypersurface of $\mathbb{H}^2\times \mathbb{H}^2$ with constant sectional curvature $-\frac{1}{2}$ and $C=0$.
Under the frame $\{E_1,E_2,E_3\}$ defined by \eqref{eqn:2.5}, if functions $b_1,b_2$ and $b_4$ defined in \eqref{eqn:2.7} are given by {\bf Case 5} of Lemma \ref{lemma:4.1aaa1},
then $M$ admits local coordinates $(u,v,t)$
and there is a function $h_3:(u,v)\mapsto \mathbb{R}$
such that the functions $b_1,b_2,b_4$ are given by
\eqref{eqn:bijth5},
where $h_3$ satisfies the following ``cosh-Gordon equation'':
%
\begin{equation}\label{eqn:3.1fffff2}
\frac{\partial^2 h_3}{\partial u \partial v}
=4\sqrt{2}\cosh(h_3).
\end{equation}

Conversely, let $\mathcal{D}$ be an open subset in $\mathbb{R}^2$ and  $h_3:\mathcal{D}\to\mathbb{R}, (u,v)\mapsto h_3(u,v)\in\mathbb{R}$
be a function which satisfies the differential equation \eqref{eqn:3.1fffff2}.
We define a metric $g$ on $\Omega=\{(u,v,t)\in \mathcal{D}\times \mathbb{R}\}$ by
\eqref{eqn:3.1fffff3},
and a $(1,1)$-tensor field $A$ on $T\Omega$ by
\eqref{eqn:3.1fffff4}.
Let $\nu$ be a vector bundle over  $\Omega$ of rank $1$ with metric $\tilde{g}$, $\nabla ^{\perp} $ a connection on $\nu$ compatible with the metric $\tilde{g}$, and $N$ a unit vector field in $\nu$.   We define a $(1,1)$-tensor field $\tilde{P}$ on $T\Omega\otimes \nu$ by
\eqref{eqn:3.1fffff5}.
Then the integrability conditions are satisfied on $\Omega$ and hence on any simply
connected subset $\Omega_1$ of $\Omega$, up to an isometry of  $\mathbb{H}^2\times \mathbb{H}^2$, there is a unique isometric immersion $\Psi$ from $(\Omega_1,~g)$ into $\mathbb{H}^2\times \mathbb{H}^2$ such that $A$ is the shape operator of  $\Omega_1$ into $\mathbb{H}^2\times \mathbb{H}^2$, $\nu$ is isometric to the normal bundle of $\Psi(\Omega_1)$ in
$\mathbb{H}^2\times \mathbb{H}^2$ by an isomorphism $\tilde{\Psi}:\nu\to T^{\perp}\Psi(\Omega_1)$,
and for any $Y\in T\Omega_1$, we have
\begin{equation}
	P(\Psi_{*}Y)=\Psi_{*}((\tilde{P}Y)^{T})+\tilde{\Psi}((\tilde{P}Y)^{\perp}),~
	P(\tilde{\Psi}N)=\Psi_*(\partial_t),
\end{equation}
where $P$ is the product structure of $\mathbb{H}^2\times \mathbb{H}^2$, $(\tilde{P}Y)^{T}$ and $(\tilde{P}Y)^{\perp}$ denote the projections of  $\tilde{P}Y$ onto $T\Omega_1$ and $\nu$, respectively.
Moreover, $\Psi(\Omega_1)$ has constant sectional curvature $-\frac{1}{2}$ and product angle function $C=0$.
\end{theorem}
%
\begin{proof}
For {\bf Case 5} of Lemma \ref{lemma:4.1aaa1}, we have
$$
\begin{aligned}
&b_1^*=\frac{1}{\sqrt{2}}e^{-\sqrt{2}(t+l_1)}\sinh(h_3),\ \
b_2^*=\frac{1}{\sqrt{2}}e^{-\sqrt{2}(t+l_1)}\cosh(h_3),\\
&b_2^*-b_1^*>0,\ \ b_2^*+b_1^*>0,\ \ 1-2b_1^{*2}+2b_2^{*2}\neq0,
\end{aligned}
$$
where function $t$ satisfies $E_3t=1$ and functions $l_1, h_3:M\rightarrow \mathbb{R}$
satisfy $E_3l_1=E_3h_3=0$.


%By \eqref{eqn:4.29}, \eqref{eqn:4.29} and $b_1=-\frac{1}{\sqrt{2}}-b_2$, $b_4=\frac{1}{\sqrt{2}}-b_2$,
%we have $E_1b_2+E_2b_2=0$, which implies that $E_2^*b_2=0$.

Now, we define $\rho_1=\frac{1}{\sqrt{b_2^*-b_1^*}}$ and $\rho_2=\frac{1}{\sqrt{b_2^*+b_1^*}}$.
Let $\tilde{E}_1^*=\rho_1E_1^*$, $\tilde{E}_2^*=\rho_2E_2^*$ and $\tilde{E}_3^*=E_3^*$.
Then, it can be checked that
\begin{equation}\label{eqn:Lietilde5}
[\tilde{E}_1^*,\tilde{E}_3^*]=0,\ \ [\tilde{E}_2^*,\tilde{E}_3^*]=0,\ \
[\tilde{E}_1^*,\tilde{E}_2^*]=-8\frac{b_2^*}{\sqrt{b_2^{*2}-b_1^{*2}}}\tilde{E}_3^*.
\end{equation}

Next, we define a differential system:
\begin{equation}\label{eqn:diff}
\left\{
\begin{aligned}
&\tilde{E}_1^*(u)=1,\ \ \tilde{E}_2^*(u)=0,\ \ \tilde{E}_3^*(u)=0,\\
&\tilde{E}_1^*(v)=0,\ \ \tilde{E}_2^*(v)=1,\ \ \tilde{E}_3^*(v)=0.
\end{aligned}\right.
\end{equation}
Then by \eqref{eqn:Lietilde5}, we know that there are functions $u$ and $v$ satisfying \eqref{eqn:diff} on $M$. Now, we choose %$f=\frac{1}{\sqrt{2}}{\rm arcsech}(\sqrt{2(b_2^2-(b_1+b_4)^2)})=t+l_1$.
$f=-\frac{1}{2\sqrt{2}}\ln\Big(2(b_2^{*2}-b_1^{*2})\Big)=t+l_1$.
In the following, we use $t$ to denote $f$.
%where function $h_3:M\rightarrow \mathbb{R}$ satisfies $E_3h_3=0$.
We define the map
$$
\Phi: M \longrightarrow \mathbb{R}^{3},
\quad p \mapsto(u(p),v(p),t(p)).
$$
Then $\Phi_*$ is given by
\begin{equation}\label{eqn:dPhi5}
\Phi_*\left(
        \begin{array}{c}
          \tilde{E}_1^* \\
          \tilde{E}_2^* \\
          \tilde{E}_3^* \\
        \end{array}
      \right)
=\left(\begin{array}{cccc}
1 & 0 & \tilde{d}_1:=\tilde{E}_1^*(t) \\
0 & 1 & \tilde{d}_2:=\tilde{E}_2^*(t) \\
0 & 0 & 1
\end{array}\right),
\end{equation}
which implies that $\{u,v,t\}$ are
local coordinates on $M$. It follows from $E_3h_3=0$ that function $h_3$ depends only on variables $\{u,v\}$. Then for {\bf Case 5} of Lemma \ref{lemma:4.1aaa1},
we have
%
\begin{equation}\label{eqn:bijth5}
\begin{aligned}
b_1&=\tfrac{1}{\sqrt{2}}{\rm sech}(\sqrt{2}(t))\Big(\sinh (h_3)-\sinh(\sqrt{2}(t))\Big),\\
b_2&=\tfrac{1}{\sqrt{2}}{\rm sech}(\sqrt{2}(t))\cosh (h_3) ,\\
b_4&=\tfrac{1}{\sqrt{2}}{\rm sech}(\sqrt{2}(t))\Big(\sinh (h_3)+\sinh(\sqrt{2}(t))\Big),\\
b_1^*&=\frac{1}{\sqrt{2}}e^{-\sqrt{2}t}\sinh(h_3),\ \
b_2^*=\frac{1}{\sqrt{2}}e^{-\sqrt{2}t}\cosh(h_3).
\end{aligned}
\end{equation}

By using \eqref{eqn:Eibar}, \eqref{eqn:dPhi5}, $\tilde{E}_1^*=\rho_1E_1^*$
and $\tilde{E}_2^*=\rho_2E_2^*$, we have
\begin{equation}\label{eqn:transframe5}
\left(
  \begin{array}{c}
    E_1 \\
    E_2 \\
    E_3 \\
  \end{array}
\right)
=\left(\begin{array}{cccc}
\frac{-1+\sqrt{2}b_1^*+\sqrt{2}b_2^*}{\rho_1(-2+4b_1^{*2}-4b_2^{*2})} & \frac{-1+\sqrt{2}b_1^*-\sqrt{2}b_2^*}{\rho_2(-2+4b_1^{*2}-4b_2^{*2})} & %\frac{-\tilde{d}_1\rho_1-\tilde{d}_2\rho_2+\sqrt{2}(b_2^*(\tilde{d}_1\rho_1-\tilde{d}_2\rho_2)
%+b_1^*(\tilde{d}_1\rho_1+\tilde{d}_2\rho_2))}{-2+4b_1^{*2}-4b_2^{*2}}
\frac{\tilde{d}_1\rho_2(-1+\sqrt{2}b_1^*+\sqrt{2}b_2^*)+\tilde{d}_2\rho_1(-1+\sqrt{2}b_1^*-\sqrt{2}b_2^*)}
{\rho_1\rho_2(-2+4b_1^{*2}-4b_2^{*2})}\\
\frac{1+\sqrt{2}b_1^*+\sqrt{2}b_2^*}{\rho_1(-2+4b_1^{*2}-4b_2^{*2})} & \frac{-1-\sqrt{2}b_1^*+\sqrt{2}b_2^*}{\rho_2(-2+4b_1^{*2}-4b_2^{*2})} & %\frac{\tilde{d}_1\rho_1-\tilde{d}_2\rho_2+\sqrt{2}(b_1^*(\tilde{d}_1\rho_1-\tilde{d}_2\rho_2)
%+b_2^*(\tilde{d}_1\rho_1+\tilde{d}_2\rho_2))}{-2+4b_1^{*2}-4b_2^{*2}}
\frac{\tilde{d}_1\rho_2(1+\sqrt{2}b_1^*+\sqrt{2}b_2^*)+\tilde{d}_2\rho_1(-1-\sqrt{2}b_1^*+\sqrt{2}b_2^*)}
{\rho_1\rho_2(-2+4b_1^{*2}-4b_2^{*2})}
\\
0 & 0 & 1
\end{array}\right)
\left(
  \begin{array}{c}
    \partial_u \\
    \partial_v \\
    \partial_t \\
  \end{array}
\right).
\end{equation}

Now, by using \eqref{eqn:4.29}, \eqref{eqn:4.30} and \eqref{eqn:transframe5}, we can get
$$
\tilde{d}_1=\frac{(h_3)_u}{\sqrt{2}},\ \ \tilde{d}_2=\frac{-(h_3)_v}{\sqrt{2}}.
$$
Then we can check the compatibility conditions of $b_1,b_2,b_4$ with respect to the Lie brackets $[E_i,E_j]$, $1\leq i<j\leq3$,
and obtain that
$$
\frac{\partial^2 h_3}{\partial u \partial v}
=4\sqrt{2}\cosh(h_3).
$$
In fact, the  compatibility conditions of  $b_1,b_2,b_4$  with respect to the Lie brackets $[E_1,E_3]$ and $[E_2,E_3]$ are automatically satisfied and the  compatibility conditions of  $b_1,b_2,b_4$  with respect to the Lie bracket $[E_1,E_2]$  imply the above differential equation.

Conversely, we assume that the function $h_3$ satisfies \eqref{eqn:3.1fffff2}, the metric $g$ on $\Omega$ is defined by
\begin{equation}\label{eqn:3.1fffff3}
\begin{aligned}
g=&\frac{1}{2}(8\sqrt{2}\cosh(\sqrt{2}t+h_3)+(h_3)_u^2)du^2
+\frac{1}{2}(8\sqrt{2}\cosh(\sqrt{2}t-h_3)+(h_3)_v^2)dv^2+dt^2\\
&+(8\sqrt{2}\sinh(h_3)-(h_3)_u(h_3)_v)dudv-(\sqrt{2}(h_3)_u)dudt+(\sqrt{2}(h_3)_v)dvdt,
\end{aligned}
\end{equation}
the $(1,1)$-tensor $A$ on $T\Omega$ is defined by
\begin{equation}\label{eqn:3.1fffff4}
\begin{aligned}
A\partial_u=&\frac{{\rm sech}(\sqrt{2}t)\sinh(h_3)}{\sqrt{2}}\partial_u
-\frac{\cosh(h_3)+\sinh(h_3)\tanh(\sqrt{2}t)}{\sqrt{2}}\partial_v\\
&+\frac{1}{2}{\rm sech}(\sqrt{2}t)(\cosh(\sqrt{2}t+h_3)(h_3)_v+\sinh(h_3)(h_3)_u)\partial_t,\\
A\partial_v=&\frac{-\cosh(h_3)+\sinh(h_3)\tanh(\sqrt{2}t)}{\sqrt{2}}\partial_u+
\frac{{\rm sech}(\sqrt{2}t)\sinh(h_3)}{\sqrt{2}}\partial_v\\
&-\frac{1}{2}{\rm sech}(\sqrt{2}t)(\cosh(\sqrt{2}t-h_3)(h_3)_u+\sinh(h_3)(h_3)_v)\partial_t,\\
A\partial_t=&0,
\end{aligned}
\end{equation}
and the $(1,1)$-tensor $\tilde{P}$ on $T\Omega\otimes \nu$ is defined by
\begin{equation}\label{eqn:3.1fffff5}
\begin{aligned}
\tilde{P}\partial_u=&-\cosh(h_3){\rm sech}(\sqrt{2}t)\partial_u+(\sinh(h_3)+\cosh(h_3)\tanh(\sqrt{2}t))\partial_v\\
&-\frac{{\rm sech}(\sqrt{2}t)(\sinh(\sqrt{2}t+h_3)(h_3)_v+\cosh(h_3)(h_3)_u)}{\sqrt{2}}\partial_t-\frac{(h_3)_u}{\sqrt{2}}N,\\
\tilde{P}\partial_v=&(-\sinh(h_3)+\cosh(h_3)\tanh(\sqrt{2}t))\partial_u+\cosh(h_3){\rm sech}(\sqrt{2}t)\partial_v\\
&+\frac{{\rm sech}(\sqrt{2}t)(\sinh(\sqrt{2}t-h_3)(h_3)_u-\cosh(h_3)(h_3)_v)}{\sqrt{2}}\partial_t+\frac{(h_3)_v}{\sqrt{2}}N,\\
\tilde{P}\partial_t=&N,\\
\tilde{P}N=&\partial_t.\\	
\end{aligned}
\end{equation}
Then it can be  checked  that all the integrability conditions are satisfied.
Therefore, the conclusions can be obtained by applying  the existence and uniqueness theorems in  \cite{K,L-T-V} directly.
\end{proof}







\begin{theorem}\label{thm:5.13a1}
Let $M$ be a hypersurface of $\mathbb{H}^2\times \mathbb{H}^2$ with constant sectional curvature $-\frac{1}{2}$ and $C=0$.
Under the frame $\{E_1,E_2,E_3\}$ defined by \eqref{eqn:2.5}, if $b_1,b_2$ and $b_4$ are
given by {\bf Case 6} of Lemma \ref{lemma:4.1aaa1}, then $M$ admits local coordinates $(u,v,t)$
and there is a non-zero function $h_4:(u,v)\mapsto \mathbb{R}$
such that the functions $b_1,b_2,b_4$ are given by \eqref{eqn:bijth6},
where $h_4$ satisfies the following ``sinh-Gordon equation'':
%
\begin{equation}\label{eqn:3.1ffffff2}
\frac{\partial^2 h_4}{\partial u \partial v}
=4\sqrt{2}\sinh(h_4).
\end{equation}

Conversely, let $\mathcal{D}$ be an open subset in $\mathbb{R}^2$ and $h_4:\mathcal{D}\to\mathbb{R}, (u,v)\mapsto h_4(u,v)\in\mathbb{R}$
be a non-zero function which satisfies the differential equation \eqref{eqn:3.1ffffff2}.
Let $\Omega_0=\{(u,v,t)\in \mathcal{D}\times \mathbb{R}|t=0\}$, $\Omega=\{(u,v,t)\in \mathcal{D}\times \mathbb{R}\}-\Omega_0\subseteq \mathbb{R}^3$, we define a metric $g$ on $\Omega$ by
\eqref{eqn:3.1ffffff3},
and a $(1,1)$-tensor field $A$ on $T\Omega$ by
\eqref{eqn:3.1ffffff4}.
Let $\nu$ be a vector bundle over  $\Omega$ of rank $1$ with metric $\tilde{g}$, $\nabla ^{\perp} $ a connection on $\nu$ compatible with the metric $\tilde{g}$, and $N$ a unit vector field in $\nu$.   We define a $(1,1)$-tensor field $\tilde{P}$ on $T\Omega\otimes \nu$ by
\eqref{eqn:3.1ffffff5}.
Then the integrability conditions are satisfied on $\Omega$ and hence on any simply
connected subset $\Omega_1$ of $\Omega$, up to an isometry of  $\mathbb{H}^2\times \mathbb{H}^2$, there is a unique isometric immersion $\Psi$ from $(\Omega_1,~g)$ into $\mathbb{H}^2\times \mathbb{H}^2$ such that $A$ is the shape operator of  $\Omega_1$ into $\mathbb{H}^2\times \mathbb{H}^2$, $\nu$ is isometric to the normal bundle of $\Psi(\Omega_1)$ in
$\mathbb{H}^2\times \mathbb{H}^2$ by an isomorphism $\tilde{\Psi}:\nu\to T^{\perp}\Psi(\Omega_1)$,
and for any $Y\in T\Omega_1$, we have
\begin{equation}
	P(\Psi_{*}Y)=\Psi_{*}((\tilde{P}Y)^{T})+\tilde{\Psi}((\tilde{P}Y)^{\perp}),~
	P(\tilde{\Psi}N)=\Psi_*(\partial_t),
\end{equation}
where $P$ is the product structure of $\mathbb{H}^2\times \mathbb{H}^2$, $(\tilde{P}Y)^{T}$ and $(\tilde{P}Y)^{\perp}$ denote the projections of  $\tilde{P}Y$ onto $T\Omega_1$ and $\nu$, respectively.
Moreover, $\Psi(\Omega_1)$ has constant sectional curvature $-\frac{1}{2}$ and product angle function $C=0$.
\end{theorem}
%
\begin{proof}
For {\bf Case 6} of Lemma \ref{lemma:4.1aaa1}, we have
$$
\begin{aligned}
&b_1^*=\frac{1}{\sqrt{2}}e^{-\sqrt{2}(t+l_2)}\cosh(h_4),\ \
b_2^*=\frac{1}{\sqrt{2}}e^{-\sqrt{2}(t+l_2)}\sinh(h_4),\\
&b_1^*-b_2^*>0,\ \ b_2^*+b_1^*>0,\ \ 1-2b_1^{*2}+2b_2^{*2}\neq0,
\end{aligned}
$$
where function $t$ satisfies $E_3t=1$ and functions $l_2, h_4:M\rightarrow \mathbb{R}$
satisfy $E_3l_2=E_3h_4=0$.




%we have $b_1^*-b_2^*>0$, $b_2^*+b_1^*>0$, and
%$1-2b_1^{*2}+2b_2^{*2}\neq0$.

%By \eqref{eqn:4.29}, \eqref{eqn:4.29} and $b_1=-\frac{1}{\sqrt{2}}-b_2$, $b_4=\frac{1}{\sqrt{2}}-b_2$,
%we have $E_1b_2+E_2b_2=0$, which implies that $E_2^*b_2=0$.

Now, we define $\rho_1=\frac{1}{\sqrt{b_1^*-b_2^*}}$ and $\rho_2=\frac{1}{\sqrt{b_2^*+b_1^*}}$.
Let $\tilde{E}_1^*=\rho_1E_1^*$, $\tilde{E}_2^*=\rho_2E_2^*$ and $\tilde{E}_3^*=E_3^*$.
Then, it can be checked that
\begin{equation}\label{eqn:Lietilde6}
[\tilde{E}_1^*,\tilde{E}_3^*]=0,\ \ [\tilde{E}_2^*,\tilde{E}_3^*]=0,\ \
[\tilde{E}_1^*,\tilde{E}_2^*]=-8\frac{b_2^*}{\sqrt{b_1^{*2}-b_2^{*2}}}\tilde{E}_3^*.
\end{equation}

We still define the differential system \eqref{eqn:diff}. By  \eqref{eqn:Lietilde6},
we know that there are functions $\{u,v\}$ satisfying \eqref{eqn:diff} on $M$.
Now, we choose %$f=\frac{1}{\sqrt{2}}{\rm arccsch}(\sqrt{2((b_1+b_4)^2-b_2^2)})=t+l_2$.
$f=-\frac{1}{2\sqrt{2}}\ln\Big(2(b_1^{*2}-b_2^{*2})\Big)=t+l_2$. In the following,
we still use $t$ to denote $f$.
%where function $h_4:M\rightarrow \mathbb{R}$ satisfies $E_3h_4=0$.

We define the map
$$
\Phi: M \longrightarrow \mathbb{R}^{3},
\quad p \mapsto(u(p),v(p),t(p)).
$$
Then $\Phi_*$ is given by
\begin{equation}\label{eqn:dPhi6}
\Phi_*\left(
        \begin{array}{c}
          \tilde{E}_1^* \\
          \tilde{E}_2^* \\
          \tilde{E}_3^* \\
        \end{array}
      \right)
=\left(\begin{array}{cccc}
1 & 0 & \tilde{d}_1:=\tilde{E}_1^*(t) \\
0 & 1 & \tilde{d}_2:=\tilde{E}_2^*(t) \\
0 & 0 & 1
\end{array}\right),
\end{equation}
which implies that $\{u,v,t\}$ are
local coordinates on $M$. It follows from $E_3h_4=0$ that function $h_4$ depends only on variables $\{u,v\}$. Then for {\bf Case 6} of Lemma \ref{lemma:4.1aaa1},
we have
%
\begin{equation}\label{eqn:bijth6}
\begin{aligned}
b_1&=\tfrac{1}{\sqrt{2}}{\rm csch}(\sqrt{2}t)\Big(\cosh(h_4)-\cosh(\sqrt{2}t)\Big),\\
b_2&=\tfrac{1}{\sqrt{2}}{\rm csch}(\sqrt{2}t)\sinh(h_4),\\
b_4&=\tfrac{1}{\sqrt{2}}{\rm csch}(\sqrt{2}t)\Big(\cosh(h_4)+\cosh(\sqrt{2}t)\Big),\\
b_1^*&=\frac{1}{\sqrt{2}}e^{-\sqrt{2}t}\cosh(h_4),\ \
b_2^*=\frac{1}{\sqrt{2}}e^{-\sqrt{2}t}\sinh(h_4).
\end{aligned}
\end{equation}

By using \eqref{eqn:Eibar}, \eqref{eqn:dPhi6}, $\tilde{E}_1^*=\rho_1E_1^*$
and $\tilde{E}_2^*=\rho_2E_2^*$, we have
\begin{equation}\label{eqn:transframe6}
\left(
  \begin{array}{c}
    E_1 \\
    E_2 \\
    E_3 \\
  \end{array}
\right)
=\left(\begin{array}{cccc}
\frac{-1+\sqrt{2}b_1^*+\sqrt{2}b_2^*}{\rho_1(-2+4b_1^{*2}-4b_2^{*2})} & \frac{-1+\sqrt{2}b_1^*-\sqrt{2}b_2^*}{\rho_2(-2+4b_1^{*2}-4b_2^{*2})} & %\frac{-\tilde{d}_1\rho_1-\tilde{d}_2\rho_2+\sqrt{2}(b_2^*(\tilde{d}_1\rho_1-\tilde{d}_2\rho_2)
%+b_1^*(\tilde{d}_1\rho_1+\tilde{d}_2\rho_2))}{-2+4b_1^{*2}-4b_2^{*2}}
\frac{\tilde{d}_1\rho_2(-1+\sqrt{2}b_1^*+\sqrt{2}b_2^*)+\tilde{d}_2\rho_1(-1+\sqrt{2}b_1^*-\sqrt{2}b_2^*)}
{\rho_1\rho_2(-2+4b_1^{*2}-4b_2^{*2})}\\
\frac{1+\sqrt{2}b_1^*+\sqrt{2}b_2^*}{\rho_1(-2+4b_1^{*2}-4b_2^{*2})} & \frac{-1-\sqrt{2}b_1^*+\sqrt{2}b_2^*}{\rho_2(-2+4b_1^{*2}-4b_2^{*2})} & %\frac{\tilde{d}_1\rho_1-\tilde{d}_2\rho_2+\sqrt{2}(b_1^*(\tilde{d}_1\rho_1-\tilde{d}_2\rho_2)
%+b_2^*(\tilde{d}_1\rho_1+\tilde{d}_2\rho_2))}{-2+4b_1^{*2}-4b_2^{*2}}
\frac{\tilde{d}_1\rho_2(1+\sqrt{2}b_1^*+\sqrt{2}b_2^*)+\tilde{d}_2\rho_1(-1-\sqrt{2}b_1^*+\sqrt{2}b_2^*)}
{\rho_1\rho_2(-2+4b_1^{*2}-4b_2^{*2})}
\\
0 & 0 & 1
\end{array}\right)
\left(
  \begin{array}{c}
    \partial_u \\
    \partial_v \\
    \partial_t \\
  \end{array}
\right).
\end{equation}

Now, by using \eqref{eqn:4.29}, \eqref{eqn:4.30} and \eqref{eqn:transframe6}, we can have
$$
\tilde{d}_1=\frac{(h_4)_u}{\sqrt{2}},\ \ \tilde{d}_2=\frac{-(h_4)_v}{\sqrt{2}}.
$$
Then we can check the compatibility conditions of $b_1,b_2,b_4$ with respect to the Lie brackets $[E_i,E_j]$, $1\leq i<j\leq3$,
and obtain that
$$
\frac{\partial^2 h_4}{\partial u \partial v}
=4\sqrt{2}\sinh(h_4).
$$
In fact, the  compatibility conditions of  $b_1,b_2,b_4$  with respect to the Lie brackets $[E_1,E_3]$ and $[E_2,E_3]$ are automatically satisfied and the  compatibility conditions of  $b_1,b_2,b_4$  with respect to the Lie bracket $[E_1,E_2]$  imply the above differential equation.

Conversely, we assume that the function $h_4$ satisfies \eqref{eqn:3.1ffffff2}, the metric $g$ on $\Omega$ is defined by %
%
\begin{equation}\label{eqn:3.1ffffff3}
\begin{aligned}
g=&\frac{1}{2}(8\sqrt{2}\cosh(\sqrt{2}t+h_4)+(h_4)_u^2)du^2
+\frac{1}{2}(8\sqrt{2}\cosh(\sqrt{2}t-h_4)+(h_4)_v^2)dv^2+dt^2\\
&+(8\sqrt{2}\cosh(h_4)-(h_4)_u(h_4)_v)dudv-(\sqrt{2}(h_4)_u)dudt+(\sqrt{2}(h_4)_v)dvdt,
\end{aligned}
\end{equation}
the $(1,1)$-tensor $A$ on $T\Omega$ is defined by
\begin{equation}\label{eqn:3.1ffffff4}
\begin{aligned}
A\partial_u=&\frac{{\rm csch}(\sqrt{2}t)\cosh(h_4)}{\sqrt{2}}\partial_u
-\frac{\sinh(h_4)+\cosh(h_4)\coth(\sqrt{2}t)}{\sqrt{2}}\partial_v\\
&+\frac{1}{2}{\rm csch}(\sqrt{2}t)(\cosh(\sqrt{2}t+h_4)(h_4)_v+\cosh(h_4)(h_4)_u)\partial_t,\\
A\partial_v=&\frac{\sinh(h_4)-\cosh(h_4)\coth(\sqrt{2}t)}{\sqrt{2}}\partial_u+
\frac{{\rm csch}(\sqrt{2}t)\cosh(h_4)}{\sqrt{2}}\partial_v\\
&-\frac{1}{2}{\rm csch}(\sqrt{2}t)(\cosh(\sqrt{2}t-h_4)(h_4)_u+\cosh(h_4)(h_4)_v)\partial_t,\\
A\partial_t=&0,
\end{aligned}
\end{equation}
and the $(1,1)$-tensor $\tilde{P}$ on $T\Omega\otimes \nu$ is defined by
\begin{equation}\label{eqn:3.1ffffff5}
\begin{aligned}
\tilde{P}\partial_u=&-\sinh(h_4){\rm csch}(\sqrt{2}t)\partial_u+(\cosh(h_4)+\sinh(h_4)\coth(\sqrt{2}t))\partial_v\\
&-\frac{{\rm csch}(\sqrt{2}t)(\sinh(\sqrt{2}t+h_4)(h_4)_v+\sinh(h_4)(h_4)_u)}{\sqrt{2}}\partial_t-\frac{(h_4)_u}{\sqrt{2}}N,\\
\tilde{P}\partial_v=&(\cosh(h_4)-\sinh(h_4)\coth(\sqrt{2}t))\partial_u+\sinh(h_4){\rm csch}(\sqrt{2}t)\partial_v\\
&+\frac{{\rm csch}(\sqrt{2}t)(\sinh(\sqrt{2}t-h_4)(h_4)_u-\sinh(h_4)(h_4)_v)}{\sqrt{2}}\partial_t+\frac{(h_4)_v}{\sqrt{2}}N,\\
\tilde{P}\partial_t=&N,\\
\tilde{P}N=&\partial_t.\\	
\end{aligned}
\end{equation}
Then it can be  checked  that all the integrability conditions are satisfied.
Therefore, the conclusions can be obtained by applying  the existence and uniqueness theorems in  \cite{K,L-T-V} directly.
\end{proof}

%\vskip 0.2 cm
\noindent{\bf Completion of the proof of Theorem \ref{thm:1.1}}.

Since $M$  has  constant sectional curvature $\kappa$, from Proposition \ref{prop:4.5}, we know that $M$ has constant sectional curvature $\kappa=-\tfrac{1}{2}$. According to the proof of Proposition \ref{prop:4.5}, we consider two situations depending on whether $C$ is $0$ or not on $M$.
When $C>0$ on $M$, then by Theorem \ref{thm:4.1abc}, 
$M$ is an open part of Example \ref{exam:4.1ss1ss11ab1}.
When $C\equiv0$ on $M$, then by Lemma \ref{lemma:4.1aaa1}, there are six cases.
For {\bf Case 1} and {\bf Case 2} of Lemma \ref{lemma:4.1aaa1}, by Proposition \ref{prop:4.1a1},
$M$ is an open part of $M_{k_1,\frac{1}{k_1}}^{1/2}$ given in Example \ref{exam:3.4fg1} for some constant $0<k_1\leq1$. For {\bf Case 3}, {\bf Case 4}, {\bf Case 5} and {\bf Case 6} of Lemma \ref{lemma:4.1aaa1}, we know that $M$ is an open part of one of the hypersurfaces described by Theorems \ref{thm:5.10a1}--\ref{thm:5.13a1}, respectively. We complete the proof
of Theorem \ref{thm:1.1}.




\section{Proofs of Theorem \ref{thm:1.2} and Theorem \ref{thm:1.3}}\label{sect:6}
%
\subsection{Proof of Theorem \ref{thm:1.2}}\label{sect:6.1}~

In this subsection, we assume that $M$ has constant product angle function and
constant mean curvature.
When $C=\pm1$ on $M$, by Lemma \ref{lem:2.2}, we know that $M$ is locally the
product of a curve $\Gamma$ in $\mathbb{H}^2$ and an open
subset of $\mathbb{H}^2$. Furthermore, if $M$ has constant mean curvature,
then $M$ is an open part of $M_\Gamma$, where $\Gamma$ is a curve of $\mathbb{H}^2$ with constant curvature.

When $C\neq\pm1$, we choose the frame $\{E_1,E_2,E_3\}$ defined by \eqref{eqn:2.5}.
Then, by $\nabla C=-2 AX$, we have $b_3=b_5=b_6=0$. It follows from \eqref{eqn:2.8a1212} that
$\rho=-2+2(b_1b_4-b_2^2)$. By \eqref{eqn:2.6a11} and \eqref{eqn:2.5aaa}, we have
\begin{equation}\label{eqn:3.2a1}
C\|A\|^2={\rm tr}(TA^2)=3H(b1-b4), 
\end{equation}
\begin{equation}\label{eqn:3.2a2}
9H^2-\|A\|^2=2(b_1b_4-b_2^2).
\end{equation}

By using \eqref{eqn:2.26}--\eqref{eqn:2.28}, \eqref{eqn:3.2a1}, \eqref{eqn:3.2a2},
$\|A\|^2=b_1^2+2b_2^2+b_4^2$ and $X=\sqrt{1-C^2}E_3$, we can check that
\begin{equation}\label{eqn:3.2aas1}
3H(X\|A\|^2)=C[-(1-C^2)\|A\|^2+(9H^2-\|A\|^2)^2].
\end{equation}

If $H\neq0$, by \eqref{eqn:3.2aas1}, we have
$$
X(\|A\|^2)=\frac{C[-(1-C^2)\|A\|^2+(9H^2-\|A\|^2)^2]}{3H}.
$$

Taking the derivative of $(b_1-b_4)$ with respect to $X$, with the use of \eqref{eqn:2.26}--\eqref{eqn:2.28}, \eqref{eqn:3.2a1} and \eqref{eqn:3.2a2}, we can check that
$$
X(b_1-b_4)=-1+C^2-C^2\|A\|^2+(9H^2-\|A\|^2)+\frac{C^2\|A\|^4}{9H^2}.
$$

Now, taking the derivative of \eqref{eqn:3.2a1} with respect to $X$, with the use of the expressions of $X(\|A\|^2)$ and $X(b_1-b_4)$, it follows that
$$
(C^2-9H^2)\|A\|^2=9H^2(1-9H^2).
$$
If $9H^2=C^2$, then $0=9H^2(1-9H^2)=C^2(1-C^2)$, so $C=0$ and $H=0$, which is a contradiction.
Then we have $\|A\|^2=\tfrac{9H^2(1-9H^2)}{C^2-9H^2}$,
which means that $\|A\|$ is constant. It implies that $M$ has constant principal curvatures.
Then, according to Theorem 1.1 of \cite{GMY},
$M$ is either an open part of $M_\tau$ for some $\tau<-1$, or
$M$ is an open part of $M_{1,1}^c$ for some $c\in(0,\frac{1}{2})\cup(\frac{1}{2},1)$, or
$M$ is an open part of $M_{1,-1}^c$ for some $c\in(0,1)$.

If $H=0$, then by \eqref{eqn:3.2a1}, we have $C=0$ or $\|A\|=0$. Note that,
according to Corollary 1.3 of \cite{Me-T}, we know that the hypersurfaces with $\|A\|=0$ corresponds to $C=\pm1$.
Thus, $M$ is a minimal hypersurface with $C=0$, which has been classified by Theorem \ref{thm:3.2aa}.
We complete the proof of Theorem \ref{thm:1.2}. 
%\qed

%\vskip 3mm

\subsection{Proof of Theorem \ref{thm:1.3}}\label{sect:6.2}~

In this subsection, we assume that $M$ has constant product angle function and
constant scalar curvature.
When $C=\pm1$ on $M$, by Lemma \ref{lem:2.2}, we know that $M$ is locally the
product of a curve $\Gamma$ in $\mathbb{H}^2$ and an open
subset of $\mathbb{H}^2$. Furthermore, if $M$ has constant scalar curvature,
then $M$ is an open part of $M_\Gamma$, where $\Gamma$ is a curve of $\mathbb{H}^2$ with constant curvature.

When $C\neq\pm1$, we choose the frame $\{E_1,E_2,E_3\}$ defined by \eqref{eqn:2.5}. Since the scalar curvature
$\rho=-2+2(b_1b_4-b_2^2)$ is constant, by using \eqref{eqn:2.26}--\eqref{eqn:2.28} and $X=\sqrt{1-C^2}E_3$, it follows from $X\rho=0$ that
\begin{equation}\label{eqn:4.1a1}
(1-C)b_1(\rho+C+3)=(1+C)b_4(\rho-C+3).
\end{equation}
Taking the derivative of equation \eqref{eqn:4.1a1} with respect to $X$, using \eqref{eqn:2.26}--\eqref{eqn:2.28}, \eqref{eqn:4.1a1} and $\rho=-2+2(b_1b_4-b_2^2)$, we can obtain
$$
(\rho+1)(\rho+3)+C^2=0.
$$
Hence there are two possible values of the scalar curvature: $\rho=-2\pm\sqrt{1-C^2}$.

{\bf Case i}: $\rho=-2+\sqrt{1-C^2}$.

Putting this value in \eqref{eqn:4.1a1}, it follows that
\begin{equation}\label{eqn:4.1a121212}
b_1\sqrt{1-C}=b_4\sqrt{1+C}.
\end{equation}
Then, by $b_1b_4-b_2^2=\frac{1}{2}(\rho+2)=\frac{\sqrt{1-C^2}}{2}$, we have $b_1\neq0$ and
$\sqrt{\frac{1-C}{1+C}}b_1^2-b_2^2=\frac{\sqrt{1-C^2}}{2}$.

Taking the derivative of $\sqrt{\frac{1-C}{1+C}}b_1^2-b_2^2=\frac{\sqrt{1-C^2}}{2}$ with respect to $E_1$ and $E_2$, with the use of \eqref{eqn:2.29}, \eqref{eqn:2.30} and \eqref{eqn:4.1a121212}, we get
$$
b_2E_1b_2-b_1E_2b_2=0,\ \ \sqrt{\frac{1-C}{1+C}}b_1E_1b_2-b_2E_2b_2=0.
$$
Now, by $\sqrt{\frac{1-C}{1+C}}b_1^2-b_2^2=\frac{\sqrt{1-C^2}}{2}\neq0$, the solution to above two equations is $E_1b_2=E_2b_2=0$.

From the fact that $[E_1,E_2]=\frac{-2b_2}{\sqrt{1-C^2}}E_3$, with the use of \eqref{eqn:2.27} and
\eqref{eqn:4.1a121212}, we have
$$
0=[E_1,E_2]b_2=\frac{-2b_2}{\sqrt{1-C^2}}E_3b_2
=\frac{-2b_2^2b_1}{\sqrt{1-C^2}}\frac{\sqrt{1-C}-\sqrt{1+C}}{\sqrt{1+C}}.
$$
Then we have $b_2=0$ or $C=0$.

{\bf Case i-1}: If $b_2=0$, then $b_1^2=\frac{1+C}{2}$ and $b_4^2=\frac{1-C}{2}$ and
$M$ has constant principal curvature. Thus, from Theorem 1.1 of \cite{GMY}, we know that $M$ is an open part of $M_{1,-1}^c$ for $c\in(0,1)$.

{\bf Case i-2}: If $C=0$, then $b_1=b_4$. By \eqref{eqn:2.27} and $E_1b_2=E_2b_2=0$, we have that
$b_2$ is constant. It follows that $b_1$ and $b_4$ are also constant.
So, it implies that $M$ has constant principal curvature.
Thus, from Theorem 1.1 of \cite{GMY}, we know that $M$ is either an open part of $M_{1,-1}^{1/2}$,
or an open part of $M_\tau$ for some $\tau<-1$.

{\bf Case ii}: $\rho=-2-\sqrt{1-C^2}$.

Putting this value in \eqref{eqn:4.1a1}, it follows that
$$
b_1\sqrt{1-C}(\sqrt{1+C}-\sqrt{1-C})=b_4\sqrt{1+C}(\sqrt{1-C}-\sqrt{1+C}).
$$
We divide the discussions into two subcases depending on the value of $C$.

{\bf Case ii-1}: If $C\neq0$, then we have
\begin{equation}\label{eqn:4.2a121212}
b_1\sqrt{1-C}=-b_4\sqrt{1+C}.
\end{equation}
Then by $b_1b_4-b_2^2=\frac{1}{2}(\rho+2)=-\frac{\sqrt{1-C^2}}{2}$, it holds  $\sqrt{\frac{1-C}{1+C}}b_1^2+b_2^2=\frac{\sqrt{1-C^2}}{2}$.

Taking the derivative of $\sqrt{\frac{1-C}{1+C}}b_1^2+b_2^2=\frac{\sqrt{1-C^2}}{2}$ with respect to $E_1$ and $E_2$, with the use of \eqref{eqn:2.29}, \eqref{eqn:2.30} and \eqref{eqn:4.2a121212}, we get
$$
-b_2E_1b_2+b_1E_2b_2=0,\ \ \sqrt{\frac{1-C}{1+C}}b_1E_1b_2+b_2E_2b_2=0.
$$
Now, by $\sqrt{\frac{1-C}{1+C}}b_1^2+b_2^2=\frac{\sqrt{1-C^2}}{2}\neq0$,
the solution to above two equations is $E_1b_2=E_2b_2=0$.

From the fact that $[E_1,E_2]=\frac{-2b_2}{\sqrt{1-C^2}}E_3$, we have
$0=[E_1,E_2]b_2=\frac{-2b_2}{\sqrt{1-C^2}}E_3b_2$. Then by \eqref{eqn:2.27} and
\eqref{eqn:4.2a121212}, we get
$$0=b_2^2b_1(\sqrt{1-C}+\sqrt{1+C}).$$

Then we have $b_2=0$ or $b_1=0$. If $b_1=0$, then $M$ has constant principal curvature,
by Theorem 1.1 of \cite{GMY}, we know that this subcase does not occur.
Thus, we have $b_2=0$, which also means that $M$ has constant principal curvature.
By Theorem 1.1 of \cite{GMY}, we know that $M$ is an open part of $M_{1,1}^c$ for $c\in(0,\frac{1}{2})\cup(\frac{1}{2},1)$.

{\bf Case ii-2}: If $C=0$, then $\rho=-3$ and $b_1b_4-b_2^2=-\frac{1}{2}$.

Now, by \eqref{eqn:2.8a1212}, the Ricci tensor is expressed by
$$
{\rm Ric}(E_1)=-E_1,\ \ {\rm Ric}(E_2)=-E_2,\ \ {\rm Ric}(E_3)=-E_3.
$$
Since ${\rm dim}~M=3$, $M$ has constant sectional curvature $-\tfrac{1}{2}$,
which has been classified by Theorem \ref{thm:1.1}.
We complete the proof of Theorem \ref{thm:1.3}.
%In conclusion, we have completed the proof of Theorem \ref{thm:1.3}.
%\qed



\section{Appendix}\label{sect:7}

In this section, we give a complete classification of the hypersurfaces in $\mathbb{S}^m\times \mathbb{S}^n$ ($m\geq3$, $n\geq2$), 
$\mathbb{H}^m\times \mathbb{H}^n$ ($m\geq3$, $n\geq2$), and $\mathbb{S}^m\times \mathbb{H}^n$ ($m,n\geq2$) with constant sectional curvature. 
For $\epsilon\in \{-1,1\}$, 
we denote by $\bar{M}^m(\epsilon)$, the unit $m$-sphere
$\mathbb S^{m}$ 
endowed with the canonical metric of constant sectional curvature $1$ when $\epsilon= 1$,
and the hyperbolic space 
$\mathbb{H}^{m}$ endowed with the canonical metric of constant sectional curvature $-1$ when $\epsilon= -1$. 
We endow $\bar{M}^m(\epsilon_1)\times \bar{M}^n(\epsilon_2)$ 
with the product metric $g$, where $(\epsilon_1,\epsilon_2)\in\{(1,1), (1,-1), (-1,-1)\}$.  
%Moreover, for $\mathbb{S}^m\times\mathbb{S}^n$ and $\mathbb{H}^m\times\mathbb{H}^n$, we assume that $m\geq3$ and $n\geq2$;  for $\mathbb{S}^m\times\mathbb{H}^n$, we assume that $m, n\geq2$.  
The product structure $P$ on $\bar{M}^m(\epsilon_1)\times \bar{M}^n(\epsilon_2)$
is defined by
%
$$
P(v_1,v_2)=(v_1,-v_2),~\hbox{for any}~(v_1,v_2)\in T(\bar{M}^m(\epsilon_1)\times \bar{M}^n(\epsilon_2)).
$$
Let $\bar{\nabla}$ be the Levi-Civita connection of the metric 
$g$ on $\bar{M}^m(\epsilon_1)\times \bar{M}^n(\epsilon_2)$. 
The curvature tensor $\bar{R}$ of $\bar{M}^m(\epsilon_1)\times \bar{M}^n(\epsilon_2)$ 
with the Riemannian product metric $g$ is given by
\begin{equation}\label{eqn:apxql}
\begin{aligned}
g(\bar{R}(U,Y)Z,W)=&\varepsilon_1\Big(g(U_1,W_1)g(Y_1,Z_1)-g(U_1,Z_1)g(Y_1,W_1)\Big)\\
&+\varepsilon_2\Big(g(U_2,W_2)g(Y_2,Z_2)-g(U_2,Z_2)g(Y_2,W_2)\Big), 
\end{aligned}
\end{equation}
where $U_1=\frac{PU+U}{2}$ and $U_2=\frac{U-PU}{2}$ for any $U\in T(\bar{M}^m(\epsilon_1)\times \bar{M}^n(\epsilon_2))$. 

Let $M$ be an oriented hypersurface of $\bar{M}^m(\epsilon_1)\times \bar{M}^n(\epsilon_2)$ with a unit normal vector field $N$ and the induced 
metric, still denoted by $g$. Then, with respect
to the product structure $P$, the smooth product angle function $C:M\rightarrow\mathbb R$
and the vector field $X$ tangent to $M$ are defined by
%
\begin{align*}
	C:=g(PN,N ),\ \ 	X:=PN-CN.
\end{align*}
%
It is obvious that $-1\le C\le1$ and $\|X\|^2:=g(X,X)=1-C^2$.
For any tangential vector field $Y$ of M, acting by the product structure $P$, we have the following decomposition
$$
PY=TY+\mu(Y)N, 
$$
where $TY$ and $\mu(Y)N$ are the tangential and normal parts of $PY$. Thus $T$ is a tensorial field of type $(1,1)$, and $\mu$ is a $1$-form over $M$. Moreover, $\mu(Y)=g(PY,N)$.



Let $\nabla$ be the Levi-Civita connection of the induced metric
$g$ on $M$. Then the Gauss and Weingarten
formulae are given respectively as below:
%
\begin{equation}\label{eqn:APX2.5}
	\bar{\nabla}_YZ=\nabla_YZ+g(AY,Z)N, ~~
	\bar{\nabla}_YN=-AY,~~\hbox{for any}~ Y,Z\in TM,
\end{equation}
%
where $A$ is the shape operator of $M$.

Let $R$ be the Riemannian curvature tensor of $M$. Then, the Gauss
and Codazzi equations of $M$ are given respectively as follows:
%
\begin{equation}\label{eqn:APX2.6}
	\begin{aligned}
		g(R(U,Y)Z,W)=g(\bar{R}(U,Y)Z,W)+g(AY,Z)g(AU,W)-g(AU,Z)g(AY,W),
	\end{aligned}
\end{equation}
%
\begin{equation}\label{eqn:APX2.7}
	g((\nabla_YA)Z-(\nabla_ZA)Y,W)=g(\bar{R}(Y,Z)W,N),
\end{equation}
%
where $U, Y, Z, W\in TM$.

Similar to the proof of Lemma 2.1 in \cite{GMY}, we still have that 
the gradient of $C$ and the covariant derivative of $X$ are given by 
%
\begin{equation}\label{eqn:APX2.9}
	\nabla C=-2AX,~~ \nabla_YX=CAY-TAY, ~~\hbox{for any}~ Y\in TM.
\end{equation}
%

In the following, by using the so-called Tsinghua principle (see the description in Sect. \ref{sect:1}), one can establish a very useful identity on a hypersurface in any Riemannian manifold with constant sectional curvature. 
%Before we give our result, we first state a useful lemma about hypersurfaces in any Riemannian manifold with constant sectional curvature. 


\begin{lemma}[\cite{AD}]\label{thm:7.1}
	Let $M$ be a hypersurface of Riemannian manifold $(\bar{M},g)$ with constant sectional curvature. 
	Then, it holds
	\begin{equation}\label{eqn:4.1q1}
		\begin{aligned}
	\mathop\mathfrak{S}\limits_{W,U,Y}g(\bar{R}(U,Y)Z,AW)=0, \ \ \forall \ U,Y,Z,W\in TM, 
		\end{aligned}
	\end{equation}
	where the symbol $\mathfrak{S}$ stands for the cyclic summation, $\bar{R}$ denotes the curvature tensor of $\bar{M}$, 
	and $A$ denotes the shape operator of $M$.
\end{lemma}

%\begin{remark}
%The ideas of Lemma \ref{lemma:3.2aa} and Lemma \ref{thm:7.1} originated from \cite{A-L-V-W}. 
%The proof of Lemma \ref{thm:7.1} also used the Bianchi's %second identity in addition. 
%\end{remark}



%
Now, we can state our result as follows:


\begin{theorem}\label{thm:1.4}
	Let $M$ be a hypersurface of $M_1^m(\epsilon_1)\times M_2^n(\epsilon_2)$ with constant sectional curvature $\kappa$,  where $m,n\geq 2$ and $(\epsilon_1,\epsilon_2)\in\{(1,1),(1,-1),(-1,-1)\}$. For $\mathbb{S}^m\times \mathbb{S}^n$ 
	and $\mathbb{H}^m\times \mathbb{H}^n$, we further assume $m\geq 3$ and $n\geq 2$. 
	Then $\kappa=-\frac{1}{2}$, $\epsilon_1=\epsilon_2=-1$ and $M$ is either
	%\begin{enumerate}
		%\item[(1)]
		%$\epsilon_1=\epsilon_2=1$, $m=n=2$, then $M$ is classified by Theorem 1.1 of \cite{LVWY}.
		
		%\item[(2)]
		%$\epsilon_1=\epsilon_2=-1$, $m=n=2$, then $M$ is classified by Theorem \ref{thm:1.1}.
		
	%	\item[(3)]
	%	$\epsilon_1=\epsilon_2=-1$, $m\geq3$ and $n\geq2$, then $M$ is either 
		\begin{enumerate}
			\item[(1)]
		locally given by the immersion $\Phi_1$, see \eqref{eqn:ap3.1ss};  or 
		\item[(2)]
		locally given by the immersion $\Phi_2$, see \eqref{eqn:ap3.1sss}; or
		\item[(3)]
		locally given by the immersion $\Phi_3$, see \eqref{eqn:ap3.1ssss}. 
			\end{enumerate}
	%\end{enumerate}
\end{theorem}
\begin{proof}
	
%Since items (1) and (2) have been proved, next we only need to consider three types of ambient space: (i) $\mathbb{S}^m\times \mathbb{S}^n$ for $m\geq3$ and $n\geq2$; 
%(ii) $\mathbb{H}^m\times \mathbb{H}^n$ for $m\geq3$ and $n\geq2$;
%(iii) $\mathbb{S}^m\times \mathbb{H}^n$ for $m\geq2$ and $n\geq2$. 
Assume that the hypersurface $M$ has constant sectional curvature $\kappa$. We will discuss two cases depending on the value of the product angle function $C$ on $M$.

{\bf Case I}. $C\neq\pm1$.

In this case,  it holds that $PN=CN+X$. 
We can take a local orthonormal frame fields $\{E_\alpha\}_{\alpha=1}^{m+n-1}$ 
such that $E_i\in P(1)$ for $1\leq i\leq m-1$, $E_p\in P(-1)$ for $m\leq p\leq m+n-2$ and  $E_{m+n-1}=\tfrac{X}{\sqrt{1-C^2}}$, where $P(1)$ (resp. $P(-1)$) is the eigenspace of $P$ corresponding to $1$ (resp. $-1$). Then it holds that 
%
\begin{equation}\label{eqn:4.1}
	PN=CN+\sqrt{1-C^2}E_{m+n-1},\ \ PE_{m+n-1}=\sqrt{1-C^2}N-CE_{m+n-1}.
\end{equation}
Let us write $AE_{\alpha}=\sum_{\beta=1}^{m+n-1}a_{\alpha\beta}E_\beta$, $1\leq \alpha\leq m+n-1$, where $a_{\alpha\beta}=a_{\beta\alpha}$ for
$1\leq \alpha,\beta\leq m+n-1$. %Let $1\leq i,j\leq m-1$ and $m\leq p,q\leq m+n-2$.
%Now, by using the Tsinghua principal, we will get linear equations of the components $\{a_{\alpha\beta}\}$.

%For instance, by direct calculations, for $1\leq i,j,k\leq m-1$ and $m\leq p,q\leq m+n-2$,
%we have

Now, taking in \eqref{eqn:4.1q1}, respectively, $1\leq i, j\leq m-1$ and $m\leq p,q\leq m+n-2$,
%
$$
(U,Y,Z,W)=(E_p,E_i,E_j,E_j), (E_p,E_i,E_q,E_q), (E_p,E_i,E_{m+n-1},E_{m+n-1}),
$$
with the use of \eqref{eqn:apxql}, we obtain 
\begin{equation}\label{eqn:ap1}
\epsilon_1(1-\delta_{ij})a_{ip}=0,
\end{equation}
\begin{equation}\label{eqn:ap2}
	\epsilon_2(1-\delta_{pq})a_{ip}=0,
\end{equation}
\begin{equation}\label{eqn:ap3}
a_{ip}\big(( C-1)\epsilon_1 +(1+C)\epsilon_2\big)=0.
\end{equation}
%Note that \eqref{eqn:ap1} only appears when $m\geq3$, and 
%\eqref{eqn:ap2} only appears when $n\geq3$.

Taking in \eqref{eqn:4.1q1}, respectively, $1\leq i, j\leq m-1$ and $m\leq p,q\leq m+n-2$,
%
$$
(U,Y,Z,W)=(E_{m+n-1},E_i,E_p,E_p),(E_{m+n-1},E_p,E_i,E_i),
$$
with the use of \eqref{eqn:apxql}, we obtain  
\begin{equation}\label{eqn:ap4}
	a_{i,m+n-1}(1+C)\epsilon_2=0,
\end{equation}
\begin{equation}\label{eqn:ap5}
	a_{p,m+n-1}(C-1)\epsilon_1=0.
\end{equation}


For any $(\epsilon_1,\epsilon_2)\in \{(1,1),(1,-1),(-1,-1)\}$, $m,n\geq2$ and $m+n\geq5$, by 
\eqref{eqn:ap1} and \eqref{eqn:ap2}, it follows that $a_{ip}=0$ for $1\leq i\leq m-1$ and $m\leq p\leq m+n-2$. 
For $(\epsilon_1,\epsilon_2)=(1,-1)$ and $m=n=2$, by \eqref{eqn:ap3}, we still have $a_{12}=0$. 
For any $(\epsilon_1,\epsilon_2)\in \{(1,1),(1,-1),(-1,-1)\}$ and $m,n\geq2$, by using $C\neq\pm1$,  \eqref{eqn:ap4} and \eqref{eqn:ap5}, we have $a_{i,m+n-1}=a_{p,m+n-1}=0$ for $1\leq i\leq m-1$ and $m\leq p\leq m+n-2$. Thus, on hypersurface $M$, it always holds that the subdistributions 
${\rm Span}\,\{E_i\}_{i=1}^{m-1}$, ${\rm Span}\,\{E_p\}_{p=m}^{m+n-2}$ and ${\rm Span}\,\{E_{m+n-1}\}$ 
are $A$-invariant, respectively. This allows us to further require the frame $\{E_i\}_{i=1}^{m+n-1}$ satisfying 
$$
AE_i=\lambda_iE_i,  AE_p=\lambda_pE_p, AE_{m+n-1}=\lambda_{m+n-1}E_{m+n-1}, 
\ \ \forall \ 1\leq i\leq m-1,\ m\leq p\leq m+n-2.
$$

Now, we calculate the sectional curvatures $K(E_i,E_j)$, $K(E_i,E_p)$, $K(E_p,E_q)$, $K(E_i,E_{m+n-1})$ and 
$K(E_p,E_{m+n-1})$. By using the Gauss equation \eqref{eqn:APX2.6} and the fact that $M$ has constant sectional curvature $\kappa$, 
we obtain 
%
\begin{gather}
	K(E_i,E_j)=\epsilon_1+\lambda_i\lambda_j=\kappa,\label{eqn:4.7s}\\
	K(E_i,E_p)=\lambda_i\lambda_{p}=\kappa,\label{eqn:4.8s}\\
	K(E_p,E_q)=\epsilon_2+\lambda_p\lambda_q=\kappa,\label{eqn:4.7adss}\\
	K(E_i,E_{m+n-1})=\tfrac{1}{2}(1-C)\epsilon_1+\lambda_i\lambda_{m+n-1}=\kappa,\label{eqn:4.9s}\\
	K(E_p,E_{m+n-1})=\tfrac{1}{2}(1+C)\epsilon_2+\lambda_p\lambda_{m+n-1}=\kappa,\label{eqn:4.9as}
\end{gather}
where \eqref{eqn:4.7s} only appears for $m\geq3$, and 
\eqref{eqn:4.7adss} only appears for $n\geq3$.


In the following, according to the ambient space, we discuss three subcases. 

{\bf Case I-i}. $\mathbb{S}^m\times\mathbb{S}^n$ for $m\geq3$ and $n\geq2$. 

In this case, $\epsilon_1=\epsilon_2=1$. If $\kappa=0$, then by \eqref{eqn:4.8s}, it holds 
$\lambda_i=0$ or $\lambda_p=0$ which contradicts with \eqref{eqn:4.7s}, \eqref{eqn:4.9as} 
and $C\neq\pm1$. 
So $\kappa\neq0$ and $\lambda_i\lambda_p\neq0$. By \eqref{eqn:4.8s}, we have 
$\lambda_i=\lambda_j$ for $1\leq i\neq j\leq m-1$.
Now, by \eqref{eqn:4.7s} and
\eqref{eqn:4.8s}, we have $\lambda_p=\tfrac{1}{\lambda_i}+\lambda_i$. By \eqref{eqn:4.7s} 
and \eqref{eqn:4.9s},
we get $\lambda_{m+n-1}=\tfrac{1+2\lambda_i^2+C}{2\lambda_i}$. By substituting $\lambda_p=\tfrac{1}{\lambda_i}+\lambda_i$ and  $\lambda_{m+n-1}=\tfrac{1+2\lambda_i^2+C}{2\lambda_i}$ into \eqref{eqn:4.9as}, with the use of  $\kappa=1+\lambda_i^2$, 
it deduces that $(1+C)(1+2\lambda_i^2)=0$, which contradicts with $C\neq\pm1$. 
%Thus, there are such hypersurface in $\mathbb{S}^m\times\mathbb{S}^n$ with constant sectional curvature. 
Thus, {\bf Case I-i} does not occur. 

\vskip 2mm

{\bf Case I-ii}. $\mathbb{H}^m\times\mathbb{H}^n$ for $m\geq3$ and $n\geq2$. 

In this case, $\epsilon_1=\epsilon_2=-1$. 
If $\kappa=0$, then by \eqref{eqn:4.8s}, it holds 
$\lambda_i=0$ or $\lambda_p=0$ which contradicts with \eqref{eqn:4.7s}, \eqref{eqn:4.9as} 
and $C\neq\pm1$. 
So $\kappa\neq0$ and $\lambda_i\lambda_p\neq0$. 
Then by  \eqref{eqn:4.8s}, we have $\lambda_i=\lambda_j$ and $\lambda_p=\lambda_q$ for  $1\leq i,j\leq m-1$ and $m\leq p,q\leq m+n-2$. By \eqref{eqn:4.7s} and \eqref{eqn:4.8s}, we have $\lambda_{p}=\lambda_{i}-\frac{1}{\lambda_{i}}$. 
By \eqref{eqn:4.7s} and \eqref{eqn:4.9s}, we get  $\lambda_{m+n-1}=-\frac{1+C-2\lambda_i^2}{2\lambda_i}$. By substituting 
$\lambda_{p}=\lambda_{i}-\frac{1}{\lambda_{i}}$ and $\lambda_{m+n-1}=-\frac{1+C-2\lambda_i^2}{2\lambda_i}$ into 
\eqref{eqn:4.9as}, it deduces that  
$(1+C)(1-2\lambda_i^2)=0$. Since $C\neq\pm1$, we get $\lambda_i=\pm\frac{1}{\sqrt{2}}$ 
and $\kappa=-\frac{1}{2}$. 
Up to a sign of the normal vector field $N$, we can assume that 
$\lambda_i=\frac{1}{\sqrt{2}}$, $\lambda_p=\frac{-1}{\sqrt{2}}$ and 
$\lambda_{m+n-1}=\frac{-C}{\sqrt{2}}$.  By \eqref{eqn:APX2.9}, 
it holds $E_{m+n-1}C=C\sqrt{2(1-C^2)}$. 
Since the distributions ${\rm Span}\,\{E_i\}_{i=1}^{m-1}$, ${\rm Span}\,\{E_p\}_{p=m}^{m+n-2}$ and ${\rm Span}\,\{E_{m+n-1}\}$ are $A$-invariant, one can check that $M$ satisfies the condition of $AT=TA$. Note that, 
hypersurfaces satisfying $AT=TA$ are called class $\mathcal{A}$ hypersurfaces 
in \cite{CT}. 

In the following, we begin to find the special examples of class $\mathcal{A}$ in $\mathbb{H}^m\times\mathbb{H}^n$ with 
constant sectional curvature. 
Let $f_1:M_1\rightarrow \mathbb{H}^m$ and $f_2:M_2\rightarrow \mathbb{H}^n$ 
be two hypersurfaces with unit normal vector fields $N^{f_1}$ and $N^{f_2}$, respectively. 
By Theorem 3.1 of \cite{CT}, we know that $M$ is locally given by the immersion: 
%
\begin{equation}\label{eqn:APXap3.1}
	\begin{aligned}
		\Phi:I\times M_1\times M_2&\longrightarrow \mathbb{H}^m\times \mathbb{H}^n,\\
		(\tilde{t},x_1,x_2)&\longmapsto (p(\tilde{t},x_1),q(\tilde{t},x_2)),
	\end{aligned}
\end{equation}
%
where
%
$$
\left\{\begin{aligned}
	p(\tilde{t},x_1)=\cosh(u_1(\tilde{t}))f_1(x_1)+\sinh(u_1(\tilde{t}))N^{f_1}(x_1),\\
	%
	q(\tilde{t},x_2)=\cosh(u_2(\tilde{t}))f_2(x_2)+\sinh(u_2(\tilde{t}))N^{f_2}(x_2),  
\end{aligned}\right.
$$
%
where $u_1(\tilde{t})$, $u_2(\tilde{t})$ are two smooth functions satisfying
$(u'_{1}(\tilde{t}))^2+(u'_{2}(\tilde{t}))^2=1$ and $u'_{1}(\tilde{t})<0$, $u'_{2}(\tilde{t})<0$.
%In \cite{CT}, Carvalho-Tojeiro proved that any hypersurface in class $\mathcal{A}$ is locally given in this way in an open and dense subset of $M$. 
Moreover, it holds that $(\frac{\partial p}{\partial \tilde{t}},\frac{\partial q}{\partial \tilde{t}})=E_{m+n-1}$. By 
$PN=CN+\sqrt{1-C^2}E_{m+n-1}$, we have $N=(\sqrt{\frac{1+C}{1-C}}\frac{\partial p}{\partial \tilde{t}},-\sqrt{\frac{1-C}{1+C}}\frac{\partial q}{\partial \tilde{t}})$.

We integrate $E_{m+n-1}C=C\sqrt{2(1-C^2)}$ along the integral curves of $E_{m+n-1}$, and obtain
$$
C=\frac{1}{\cosh(\sqrt{2}\tilde{t}+d_1)},\ {\rm or}\ \  C=0,
\ {\rm or}\ \  C=\frac{-1}{\cosh(\sqrt{2}\tilde{t}+d_1)}, 
$$
where $d_1$ is constant and $\sqrt{2}\tilde{t}+d_1<0$. By taking the reparametrization $t=\tilde{t}+\frac{d_1}{\sqrt{2}}$, it follows that 
$$
C=\frac{1}{\cosh(\sqrt{2}t)},\ {\rm or}\ \  C=0,
\ {\rm or}\ \  C=\frac{-1}{\cosh(\sqrt{2}t)},\ \ t<0.  
$$
%Then the map $\Phi$ can be rewritten as 
%\begin{equation}\label{eqn:ap3.1}
%	\begin{aligned}
%		\Phi:I\times M_1\times M_2&\longrightarrow \mathbb{H}^m\times \mathbb{H}^n,\\
%		(t,x_1,x_2)&\longmapsto (p(t,r),q(t,s)). 
%	\end{aligned}
%\end{equation}


Let $\{X_i\}_{i=1}^{m-1}$ and $\{X_p\}_{p=m}^{m+n-2}$ be bases of $T_{x_1}M_1$ and $T_{x_2}M_2$ such that $\{df_1(X_i)\}_{i=1}^{m-1}$ and $\{df_2(X_p)\}_{p=m}^{m+n-2}$ are 
principal curvature vectors of $f_1(M_1)$ and $f_2(M_2)$, respectively. Assume that 
$$
A^{f_1}df_1(X_i)=\lambda^{f_1}_i(x_1)df_1(X_i),\ A^{f_2}df_2(X_p)=\mu^{f_2}_p(x_2)df_2(X_p),\ \forall \ 
1\leq i \leq m-1,\ m\leq p \leq m+n-2, 
$$
where $A^{f_1}$ and $A^{f_2}$ are the shape operators of $f_1(M_1)\hookrightarrow\mathbb{H}^m$ and $f_2(M_2)\hookrightarrow\mathbb{H}^n$,  respectively.
%the corresponding principal curvatures are 
%$\lambda^{f_1}_i(x_1)$ and $\lambda^{f_2}_p(x_2)$. 
Now, at point $\Phi(t,x_1,x_2)$, the basis $\{E_\alpha\}_{\alpha=1}^{m+n-1}$ of $\Phi$ can be given as follows:
%
\begin{equation*}
	\begin{aligned}
		&E_i=\left(df_1(X_i),0\right),\ 
		%
		E_p=\left(0,df_2(X_p)\right), \ 
		1\leq i \leq m-1,\ m\leq p \leq m+n-2, \\ 
		&E_{m+n-1}=\tfrac{\partial}{\partial t}(p(t,x_1),q(t,x_2))
		=\big(u'_{1}(t)\sinh(u_1(t))f_1(x_1)
		+u'_{1}(t)\cosh(u_1(t))N^{f_1}(x_1),\\
		&~~~~~~~~~~~~~~~~~~~~~~~~   u'_{2}(t)\sinh(u_2(t))f_2(x_2)
		+u'_{2}(t)\cosh(u_2(t))N^{f_2}(x_2)\big).
	\end{aligned}
\end{equation*}
The unit normal vector field $N$ of $M$ in $\mathbb{H}^m\times \mathbb{H}^n$ is given by
%
\begin{equation}\label{eqn:apx3.3}
	\begin{aligned}
		N=&\big(u'_{2}(t)\sinh(u_1(t))f_1(x_1)
		+u'_{2}(t)\cosh(u_1(t))N^{f_1}(x_1),\\
		%
		&\ \ \ -u'_{1}(t)\sinh(u_2(t))f_2(x_2)
		-u'_{1}(t)\cosh(u_2(t))N^{f_2}(x_2)\big).
	\end{aligned}
\end{equation}
%
It follows that
%
\begin{equation}\label{eqn:apxCC}
	C=(u'_{2}(t))^2-(u'_{1}(t))^2=1-2(u'_{1}(t))^2.
\end{equation}
%
%which satisfies $|C|<1$. 

Now, by direct calculation, at point $\Phi(t,x_1,x_2)$, we can get 
%
\begin{equation}\label{eqn:apx3.5}
	\left\{\begin{aligned}
		&AE_i=-\tfrac{\sinh(u_1(t))
			-\lambda^{f_1}_i(x_1)\cosh(u_1(t))}
		{\cosh(u_1(t))-\lambda^{f_1}_i(x_1)\sinh(u_1(t))}u'_{2}(t)E_i,\ \ 1\leq i\leq m-1,\\
		%
		&AE_p=\tfrac{\sinh(u_2(t))
			-\mu^{f_2}_p(x_2)\cosh(u_2(t))}
		{\cosh(u_2(t))-\mu^{f_2}_p(x_2)\sinh(u_2(t))}u'_{1}(t)E_p,\ \ m\leq p\leq m+n-2,\\
		%
		&AE_{m+n-1}=\left(u'_{2}(t)u''_{1}(t)-u'_{1}(t)u''_{2}(t)\right)E_{m+n-1}.
	\end{aligned}\right.
\end{equation}
Thus, it holds that 
%
\begin{equation}\label{eqn:apx3.5ss}
\begin{aligned}
&\lambda_i=-\tfrac{\sinh(u_1(t))
	-\lambda^{f_1}_i(x_1)\cosh(u_1(t))}
{\cosh(u_1(t))-\lambda^{f_1}_i(x_1)\sinh(u_1(t))}u'_{2}(t), \ \ \lambda_p=\tfrac{\sinh(u_2(t))
	-\mu^{f_2}_p(x_2)\cosh(u_2(t))}
{\cosh(u_2(t))-\mu^{f_2}_p(x_2)\sinh(u_2(t))}u'_{1}(t),\\
&\lambda_{m+n-1}=u'_{2}(t)u''_{1}(t)-u'_{1}(t)u''_{2}(t).
	\end{aligned}
\end{equation}

By \eqref{eqn:APX2.9} and $E_{m+n-1}=\tfrac{X}{\sqrt{1-C^2}}$, we can have 
 $\nabla_{E_{m+n-1}}{E_{m+n-1}}=0$. 
Then by using \eqref{eqn:4.1}, $AE_{m+n-1}=-\frac{C}{\sqrt{2}}E_{m+n-1}$ and $\nabla_{E_{m+n-1}}{E_{m+n-1}}=0$, 
we can get 
\begin{equation}\label{eqn:APX4.mm1}
	\begin{aligned}
		(\tfrac{\partial^2 p}{\partial t^2},\tfrac{\partial^2 q}{\partial t^2})&%=D_{E_3}{E_3}
		=\bar{\nabla}_{E_{m+n-1}}{E_{m+n-1}}-g( D_{E_{m+n-1}}{E_{m+n-1}},\tfrac{(p,q)}{\sqrt{2}})\tfrac{(p,q)}{\sqrt{2}}\\
		&~~~~~~~~~~~~~~~~    -g( D_{E_{m+n-1}}{E_{m+n-1}},\tfrac{(p,-q)}{\sqrt{2}})\tfrac{(p,-q)}{\sqrt{2}}\\
		&=-\tfrac{C}{\sqrt{2}}N+\tfrac{1}{2}(p,q)-\tfrac{C}{2}(p,-q)\\
		&=(-C\sqrt{\tfrac{1+C}{2(1-C)}}\tfrac{\partial p}{\partial t},C\sqrt{\tfrac{1-C}{2(1+C)}}\tfrac{\partial q}{\partial t})+(\tfrac{1-C}{2}p,\tfrac{1+C}{2}q),
	\end{aligned}
\end{equation}
where $D$ is the canonical Euclidean connection on $\mathbb{R}_2^{m+n+2}$. 
It can be checked that equation \eqref{eqn:APX4.mm1} is equivalent to the following 
ordinary differential equation system:
\begin{equation}\label{eqn:APX4.mm1sx}
	\left\{\begin{aligned}
	&u_1''(t)=-\frac{1}{\sqrt{2}}\Big((u_2'(t))^2-(u_1'(t))^2\Big)u_2'(t),\\
	&u_2''(t)=\frac{1}{\sqrt{2}}\Big((u_2'(t))^2-(u_1'(t))^2\Big)u_1'(t). 
\end{aligned}\right.
\end{equation}




In the following, we divide the discussion into three subcases, according to the expression of the function $C$ on $M$. 

{\bf Case I-ii-1}. $C=\frac{1}{\cosh(\sqrt{2}t)}$. 

In this subcase, from $(u'_{1}(t))^2+(u'_{2}(t))^2=1$, $(u'_{2}(t))^2-(u'_{1}(t))^2=C$ and $C=\frac{1}{\cosh(\sqrt{2}t)}$, $t<0$, we solve the equation \eqref{eqn:APX4.mm1sx} and obtain 
\begin{equation}\label{eqn:APXpq}
	\begin{aligned}
		u_1(t)=\ln\Big(\sqrt{2}\cosh(\tfrac{t}{\sqrt{2}})+\sqrt{\cosh(\sqrt{2}t)}\Big)-a,\ \
		u_2(t)={\rm arcsinh}\Big(-\sqrt{2}\sinh(\tfrac{t}{\sqrt{2}})\Big)-b, 
	\end{aligned}
\end{equation}
where $a,b$ are constants. 
From $\lambda_i=\frac{1}{\sqrt{2}}$ and $\lambda_p=\frac{-1}{\sqrt{2}}$ for $1\leq i\leq m-1$ and $m\leq p\leq m+n-2$, by using \eqref{eqn:apx3.5ss} and \eqref{eqn:APXpq}, we have 
$\lambda^{f_1}_i(x_1)=-\tanh(a)$ and $\mu^{f_2}_p(x_2)=-\coth(b)$. It means that 
$f_1(M_1)$ is a totally umbilical hypersurface of 
$\mathbb{H}^m$ with constant principal curvatures $-\tanh(a)$, and $f_2(M_2)$ is a totally umbilical hypersurface of 
$\mathbb{H}^n$ with constant principal curvatures $-\coth(b)$. 
Now, we have 
$$
\left\{
\begin{aligned}
	&p(t,x_1)=\sqrt{\cosh(\sqrt{2}t)+1}\big(\cosh(a)f_1(x_1)-\sinh(a)N^{f_1}(x_1)\big)\\
&\ \ \ \ \ \ \ \ \ \ \ \ \ \ \ \ \ \ \ \ \ \ \ +\sqrt{\cosh(\sqrt{2}t)}\big(-\sinh(a)f_1(x_1)+\cosh(a)N^{f_1}(x_1)\big),\\
	&q(t,x_2)=\sqrt{\cosh(\sqrt{2}t)}\big(\cosh(b)f_2(x_2)-\sinh(b)N^{f_2}(x_2)\big)\\
	&\ \ \ \ \ \ \ \ \ \ \ \ \ \ \ \ \ \ \ \ \ \ \
	+\sqrt{\cosh(\sqrt{2}t)-1}\big(-\sinh(b)f_2(x_2)+\cosh(b)N^{f_2}(x_2)\big).
\end{aligned}\right.
$$
One can check that $\cosh(a)f_1(x_1)-\sinh(a)N^{f_1}(x_1)$ is a totally geodesic hypersurface in $\mathbb{H}^{m}$, and $\cosh(b)f_2(x_2)-\sinh(b)N^{f_2}(x_2)$ is a point in $\mathbb{H}^{n}$. 

Since the immersion $\Phi$ is constructed by using the parallel hypersurfaces of $f_1(M_1)$ and $f_2(M_2)$,  
then up to an isometric of $\mathbb{H}^m\times \mathbb{H}^n$, we can rewrite $\Phi$ as follows:  
\begin{equation}\label{eqn:ap3.1ss}
	\begin{aligned}
		\Phi_1:I\times \mathbb{H}^{m-1}\times \mathbb{S}^{n-1}&\longrightarrow \mathbb{H}^m\times \mathbb{H}^n,\\
		(t,x,y)&\longmapsto (p(t,x),q(t,y)),
	\end{aligned}
\end{equation}
%
where
%
$$
\left\{\begin{aligned}
	&p(t,x)=\Big(\sqrt{\cosh(\sqrt{2}t)+1}\ x,\sqrt{\cosh(\sqrt{2}t)}\Big),\\
	&q(t,y)=\Big(\sqrt{\cosh(\sqrt{2}t)},\sqrt{\cosh(\sqrt{2}t)-1}\ y\Big).
\end{aligned}\right.
$$



{\bf Case I-ii-2}. $C=0$. 

In this subcase, we have $u'_{2}(t)^2-u'_{1}(t)^2=0$ and $u'_{1}(t)^2+u'_{2}(t)^2=1$. Thus, 
without loss of generality, we can assume that $u_{1}(t)=-\frac{t}{\sqrt{2}}-a$ and 
$u_{2}(t)=-\frac{t}{\sqrt{2}}-b$, where $a$ and $b$ are constant.
From $\lambda_i=\frac{1}{\sqrt{2}}$, $\lambda_p=\frac{-1}{\sqrt{2}}$ for $1\leq i\leq m-1$ and $m\leq p\leq m+n-2$, and \eqref{eqn:apx3.5ss}, it follows that  $\lambda^{f_1}_i(x_1)=\mu^{f_2}_p(x_2)=-1$, which means that $f_1(M_1)$ (resp. $f_2(M_2)$) is a horosphere of 
$\mathbb{H}^m$ (resp. $\mathbb{H}^n$) with constant principal curvatures $-1$. 

Since the immersion $\Phi$ is constructed by using the parallel hypersurfaces of $f_1(M_1)$ and $f_2(M_2)$, 
up to an isometric of $\mathbb{H}^m\times \mathbb{H}^n$, we can rewrite $\Phi$ as follows:  
\begin{equation}\label{eqn:ap3.1sss}
	\begin{aligned}
		\Phi_2:I\times \mathbb{R}^{m-1}\times \mathbb{R}^{n-1}&\longrightarrow \mathbb{H}^m\times \mathbb{H}^n,\\
		\big(t,(y_1,\ldots,y_{m-1}),(z_1,\ldots,z_{n-1})\big)&\longmapsto \big(p(t,y_1,\ldots,y_{m-1}),q(t,z_1,\ldots,z_{n-1})\big),
	\end{aligned}
\end{equation}
%
where
%
$$
\left\{\begin{aligned}
	&p(t,y_1,\ldots,y_{m-1})=\cosh(\frac{t}{\sqrt{2}})
	(1+\frac{1}{2}\sum_{i=1}^{m-1}y_i^2,y_1,\ldots,y_{m-1},\frac{1}{2}\sum_{i=1}^{m-1}y_i^2)\\
	&~~~~~~~~~~~~~~~~~~~~~~~~~~~~~~~~~~~                      +\sinh(\frac{t}{\sqrt{2}})(-\frac{1}{2}
	\sum_{i=1}^{m-1}y_i^2,-y_1,\ldots,-y_{m-1},1-\frac{1}{2}\sum_{i=1}^{m-1}y_i^2),\\
	&q(t,z_1,\ldots,z_{n-1})=\cosh(\frac{t}{\sqrt{2}})
	(1+\frac{1}{2}\sum_{i=1}^{n-1}z_i^2,z_1,\ldots,z_{n-1},\frac{1}{2}\sum_{i=1}^{n-1}z_i^2)\\
	&~~~~~~~~~~~~~~~~~~~~~~~~~~~~~~~~~~~           +\sinh(\frac{t}{\sqrt{2}})(-\frac{1}{2}
	\sum_{i=1}^{n-1}z_i^2,-z_1,\ldots,-z_{n-1},1-\frac{1}{2}\sum_{i=1}^{n-1}z_i^2).
\end{aligned}\right.
$$


{\bf Case I-ii-3}. $C=-\frac{1}{\cosh(\sqrt{2}t)}$. 

In this subcase, from $(u'_{1}(t))^2+(u'_{2}(t))^2=1$, $(u'_{2}(t))^2-(u'_{1}(t))^2=C$ and $C=\frac{-1}{\cosh(\sqrt{2}t)}$, $t<0$, we solve the equation \eqref{eqn:APX4.mm1sx} and obtain 
\begin{equation}\label{eqn:APXpqs}
		u_1(t)={\rm arcsinh}\Big(-\sqrt{2}\sinh(\tfrac{t}{\sqrt{2}})\Big)-a, \ \ 
		u_2(t)=\ln\Big(\sqrt{2}\cosh(\tfrac{t}{\sqrt{2}})+\sqrt{\cosh(\sqrt{2}t)}\Big)-b,
\end{equation}
where $a,b$ are constants.  
From $\lambda_i=\frac{1}{\sqrt{2}}$, $\lambda_p=\frac{-1}{\sqrt{2}}$ for $1\leq i\leq m-1$ and $m\leq p\leq m+n-2$, by using \eqref{eqn:apx3.5ss} and \eqref{eqn:APXpqs}, it follows that 
$\lambda^{f_1}_i(x_1)=-\coth(a)$ and $\mu^{f_2}_p(x_2)=-\tanh(b)$. It means that 
$f_1(M_1)$ is a totally umbilical hypersurface of 
$\mathbb{H}^m$ with constant principal curvatures $-\coth(a)$, and $f_2(M_2)$ is a totally umbilical hypersurface of 
$\mathbb{H}^n$ with constant principal curvatures $-\tanh(b)$. 
Now, we have 
$$
\left\{
\begin{aligned}
	&p(t,x_1)=\sqrt{\cosh(\sqrt{2}t)}\big(\cosh(a)f_1(x_1)-\sinh(a)N^{f_1}(x_1)\big)\\
	&\ \ \ \ \ \ \ \ \ \ \ \ \ \ \ \ \ \ \ \ \ \ \
	+\sqrt{\cosh(\sqrt{2}t)-1}\big(-\sinh(a)f_1(x_1)+\cosh(a)N^{f_1}(x_1)\big),\\
	&q(t,x_2)=\sqrt{\cosh(\sqrt{2}t)+1}\big(\cosh(b)f_2(x_2)-\sinh(b)N^{f_2}(x_2)\big)\\
	&\ \ \ \ \ \ \ \ \ \ \ \ \ \ \ \ \ \ \ \ \ \ \
	+\sqrt{\cosh(\sqrt{2}t)}\big(-\sinh(b)f_2(x_2)+\cosh(b)N^{f_2}(x_2)\big).
\end{aligned}\right.
$$
One can check that $\cosh(a)f_1(x_1)-\sinh(a)N^{f_1}(x_1)$ is a point in $\mathbb{H}^{m}$, and $\cosh(b)f_2(x_2)-\sinh(b)N^{f_2}(x_2)$ is a totally geodesic hypersurface in $\mathbb{H}^{n}$. 



Since the immersion $\Phi$ is constructed by using the parallel hypersurfaces of $f_1(M_1)$ and $f_2(M_2)$, 
up to an isometric of $\mathbb{H}^m\times \mathbb{H}^n$, we can rewrite $\Phi$ as follows: 
\begin{equation}\label{eqn:ap3.1ssss}
	\begin{aligned}
		\Phi_3:I\times \mathbb{S}^{m-1}\times \mathbb{H}^{n-1}&\longrightarrow \mathbb{H}^m\times \mathbb{H}^n,\\
		(t,x,y)&\longmapsto (p(t,x),q(t,y)),
	\end{aligned}
\end{equation}
%
where
%
$$
\left\{\begin{aligned}
	&p(t,x)=\Big(\sqrt{\cosh(\sqrt{2}t)},\sqrt{\cosh(\sqrt{2}t)-1}\ x\Big),\\
	&q(t,y)=\Big(\sqrt{\cosh(\sqrt{2}t)+1}\ y,\sqrt{\cosh(\sqrt{2}t)}\Big).
\end{aligned}\right.
$$

\vskip 2mm

{\bf Case I-iii}. $\mathbb{S}^m\times\mathbb{H}^n$ for $m\geq2$ and $n\geq2$. 

In this case, $\epsilon_1=1$, $\epsilon_2=-1$.  
If $\kappa=0$, then by \eqref{eqn:4.8s}, it holds 
$\lambda_i=0$ or $\lambda_p=0$ which contradicts with \eqref{eqn:4.7s}, \eqref{eqn:4.9as} 
and $C\neq\pm1$. 
So $\kappa\neq0$ and $\lambda_i\lambda_p\neq0$. 
Then by  \eqref{eqn:4.8s}, we have $\lambda_i=\lambda_j$ and $\lambda_p=\lambda_q$ for  $1\leq i,j\leq m-1$ and $m\leq p,q\leq m+n-2$. 

In the following, according to the values of $m$ and $n$, we divide the discussion into three subcases. 

{\bf Case I-iii-1}. $m\geq 3$. 

In this subcase, by \eqref{eqn:4.7s} and
\eqref{eqn:4.8s}, we have $\lambda_p=\tfrac{1}{\lambda_i}+\lambda_i$ for any $m\leq p\leq m+n-2$. 
By \eqref{eqn:4.7s} and \eqref{eqn:4.9s}, we have 
$\lambda_{m+n-1}=\frac{1+2\lambda_i^2+C}{2\lambda_i}$.  By substituting 
$\lambda_p=\tfrac{1}{\lambda_i}+\lambda_i$ and  $\lambda_{m+n-1}=\frac{1+2\lambda_i^2+C}{2\lambda_i}$ into \eqref{eqn:4.9as}, 
we get 
$\frac{1+2(\lambda_i^2+\lambda_i^4)+C}{2\lambda_i^2}=\kappa$, which combines with 
\eqref{eqn:4.7s}, it deduces that $C=-1$. It is a contradiction. 

{\bf Case I-iii-2}. $m=2$ and $n\geq 3$. 

In this subcase, by \eqref{eqn:4.8s} and \eqref{eqn:4.7adss}, we have $\lambda_i=\tfrac{-1}{\lambda_p}+\lambda_p$. 
By \eqref{eqn:4.9as} and \eqref{eqn:4.7adss}, we have 
$\lambda_{m+n-1}=\frac{C-1+2\lambda_p^2}{2\lambda_p}$.  We substitute $\lambda_i=\tfrac{-1}{\lambda_p}+\lambda_p$ and $\lambda_{m+n-1}=\frac{C-1+2\lambda_p^2}{2\lambda_p}$ into \eqref{eqn:4.9s}, and using $\kappa=-1+\lambda_p^2$, 
we can have $1-C=0$, which is a contradiction. 

{\bf Case I-iii-3}. $m=n=2$. 

In this subcase, if $\lambda_3=0$, then \eqref{eqn:4.9s} and \eqref{eqn:4.9as} become to 
$\frac{1-C}{2}=\frac{-1-C}{2}=\kappa$, which is a contradiction. So $\lambda_3\neq0$. 
By \eqref{eqn:4.9s} and \eqref{eqn:4.9as}, we have $\lambda_1=\frac{-1+C+2\kappa}{2\lambda_3}$ and $\lambda_2=\frac{1+C+2\kappa}{2\lambda_3}$. These combine 
with \eqref{eqn:4.8s}, we have $\frac{(-1+C+2\kappa)(1+C+2\kappa)}{4\lambda_3^2}=\kappa$. 
It follows from $C\neq\pm1$ that $\kappa\neq0$ and $\kappa=\frac{1}{2}(\lambda_3^2-C-\sqrt{1+\lambda_3^4-2\lambda_3^2C})$ or $\kappa=\frac{1}{2}(\lambda_3^2-C+\sqrt{1+\lambda_3^4-2\lambda_3^2C})$. 

If $\kappa=\frac{1}{2}(\lambda_3^2-C-\sqrt{1+\lambda_3^4-2\lambda_3^2C})$, then by the fact that $E_3\kappa=0$ and $E_3C=-2\lambda_3\sqrt{1-C^2}$, we have 
$$
E_3\lambda_3=\frac{\sqrt{1-C^2}(\lambda_3^2-\sqrt{1+\lambda_3^4-2\lambda_3^2C})}{-\lambda_3^2+C+\sqrt{1+\lambda_3^4-2\lambda_3^2C}}. 
$$
By $\lambda_1=\frac{-1+C+2\kappa}{2\lambda_3}$,  $\lambda_2=\frac{1+C+2\kappa}{2\lambda_3}$ and $E_3C=-2\lambda_3\sqrt{1-C^2}$, 
we have 
\begin{equation}\label{eqn:E312}
E_3\lambda_1=\frac{\sqrt{1-C^2}(1-\lambda_3^2+\sqrt{1+\lambda_3^4-2\lambda_3^2C})}
{2\lambda_3^2(\lambda_3^2-C-\sqrt{1+\lambda_3^4-2\lambda_3^2C})},\ \ 
E_3\lambda_2=\frac{\sqrt{1-C^2}(1+\lambda_3^2-\sqrt{1+\lambda_3^4-2\lambda_3^2C})}
{2\lambda_3^2(\lambda_3^2-C-\sqrt{1+\lambda_3^4-2\lambda_3^2C})}.
\end{equation}

On the other hand, we calculate 
$g\big((\nabla_{E_3}A)E_{1}-(\nabla_{E_1}A)E_{3},E_1\big)$ and $g\big((\nabla_{E_3}A)E_{2}-(\nabla_{E_2}A)E_{3},E_2\big)$. 
By direct calculation, with the use of \eqref{eqn:APX2.9} and Codazzi equation \eqref{eqn:APX2.7}, it holds that  
\begin{equation}\label{eqn:apxe11}
E_3\lambda_1=\frac{(1-C)(\lambda_1^2-\lambda_1\lambda_3)}{\sqrt{1-C^2}}+\frac{\sqrt{1-C^2}}{2},\ 
E_3\lambda_2=\frac{(1+C)(-\lambda_2^2+\lambda_2\lambda_3)}{\sqrt{1-C^2}}+\frac{\sqrt{1-C^2}}{2}.
\end{equation}
By using $\lambda_1=\frac{-1+C+2\kappa}{2\lambda_3}$ and $\lambda_2=\frac{1+C+2\kappa}{2\lambda_3}$, we  compare \eqref{eqn:apxe11} with \eqref{eqn:E312}, it follows that 
$$
\frac{-\sqrt{1-C^2}(1-\lambda_3^2+\sqrt{1+\lambda_3^4-2\lambda_3^2C})}
{\lambda_3^2(-\lambda_3^2+C+\sqrt{1+\lambda_3^4-2\lambda_3^2C})}=0,\ \ 
\frac{\sqrt{1-C^2}(-1-\lambda_3^2+\sqrt{1+\lambda_3^4-2\lambda_3^2C})}
{\lambda_3^2(-\lambda_3^2+C+\sqrt{1+\lambda_3^4-2\lambda_3^2C})}=0. 
$$
Since $C\neq\pm1$, we get a contradiction.



If $\kappa=\frac{1}{2}(\lambda_3^2-C+\sqrt{1+\lambda_3^4-2\lambda_3^2C})$, then by the fact that $E_3\kappa=0$ and $E_3C=-2\lambda_3\sqrt{1-C^2}$, we have 
$$
E_3\lambda_3=-\frac{\sqrt{1-C^2}(\lambda_3^2+\sqrt{1+\lambda_3^4-2\lambda_3^2C})}{\lambda_3^2-C+\sqrt{1+\lambda_3^4-2\lambda_3^2C}}. 
$$
By $\lambda_1=\frac{-1+C+2\kappa}{2\lambda_3}$ and $\lambda_2=\frac{1+C+2\kappa}{2\lambda_3}$, we have 
\begin{equation}\label{eqn:E312s}
	E_3\lambda_1=-\frac{\sqrt{1-C^2}(-1+\lambda_3^2+\sqrt{1+\lambda_3^4-2\lambda_3^2C})}
	{2\lambda_3^2(\lambda_3^2-C+\sqrt{1+\lambda_3^4-2\lambda_3^2C})},\ \ 
	E_3\lambda_2=\frac{\sqrt{1-C^2}(1+\lambda_3^2+\sqrt{1+\lambda_3^4-2\lambda_3^2C})}
	{2\lambda_3^2(\lambda_3^2-C+\sqrt{1+\lambda_3^4-2\lambda_3^2C})}.
\end{equation}

On the other hand, we calculate 
$g\big((\nabla_{E_3}A)E_{1}-(\nabla_{E_1}A)E_{3},E_1\big)$ and $g\big((\nabla_{E_3}A)E_{2}-(\nabla_{E_2}A)E_{3},E_2\big)$. 
By direct calculation, with the ue of \eqref{eqn:APX2.9} and Codazzi equation \eqref{eqn:APX2.7}, we get 
\begin{equation}\label{eqn:apxE312s}
E_3\lambda_1=\frac{(1-C)(\lambda_1^2-\lambda_1\lambda_3)}{\sqrt{1-C^2}}+\frac{\sqrt{1-C^2}}{2},\  
E_3\lambda_2=\frac{(1+C)(-\lambda_2^2+\lambda_2\lambda_3)}{\sqrt{1-C^2}}+\frac{\sqrt{1-C^2}}{2}.
\end{equation}

By substituting $\lambda_1=\frac{-1+C+2\kappa}{2\lambda_3}$ and $\lambda_2=\frac{1+C+2\kappa}{2\lambda_3}$ into the right hands of \eqref{eqn:apxE312s}, 
we compare \eqref{eqn:apxE312s} with \eqref{eqn:E312s}, it follows that  
$$
\frac{-\sqrt{1-C^2}(-1+\lambda_3^2+\sqrt{1+\lambda_3^4-2\lambda_3^2C})}
{\lambda_3^2(\lambda_3^2-C+\sqrt{1+\lambda_3^4-2\lambda_3^2C})}=0,\ \ 
\frac{\sqrt{1-C^2}(1+\lambda_3^2+\sqrt{1+\lambda_3^4-2\lambda_3^2C})}
{\lambda_3^2(\lambda_3^2-C+\sqrt{1+\lambda_3^4-2\lambda_3^2C})}=0. 
$$
Since $C\neq\pm1$, we get a contradiction. 
Thus {\bf Case I-iii} does not occur. 


\vskip 2mm

{\bf Case II}. $C^2=1$.

In this case, we know that $M$ is an open part of the Riemannian product manifold 
$\Gamma_1\times M_2^n(\epsilon_2)$  or $M_1^m(\epsilon_1)\times\Gamma_2$, where $\Gamma_1$ 
(resp. $\Gamma_2$) is a hypersurface of $M_1^m(\epsilon_1)$ (resp. $M_2^n(\epsilon_2)$). 
Since $\epsilon_1\epsilon_2\neq0$, it follows that the Riemannian product manifolds 
$\Gamma_1\times M_2^n(\epsilon_2)$ and $M_1^m(\epsilon_1)\times\Gamma_2$ 
can not have constant sectional curvature. 

We have completed the proof of Theorem \ref{thm:1.4}. 
\end{proof}


\begin{remark}
The hypersurface $\Phi_1(I\times \mathbb{H}^{m-1}\times \mathbb{S}^{n-1})$ is the high-dimensional version of  Example \ref{exam:4.1ss1ss11ab1}. 
The hypersurface $\Phi_2(I\times \mathbb{R}^{m-1}\times \mathbb{R}^{n-1})$ is the high-dimensional version of  hypersurface $M_{1,1}^{1/2}$, which appears in Example \ref{exam:3.4fg1}. 
When $m=n\geq3$, up to the isometry map
$$
\mathbb{F}:\mathbb{H}^m\times \mathbb{H}^m\rightarrow\mathbb{H}^m\times \mathbb{H}^m,\ \
\mathbb{F}(p,q)=(q,p),
$$
the hypersurface $\Phi_1(I\times \mathbb{H}^{m-1}\times \mathbb{S}^{m-1})$ is congruent to 
the hypersurface $\Phi_3(I\times \mathbb{S}^{m-1}\times \mathbb{H}^{m-1})$. 
\end{remark}



\begin{remark}
Combining Theorem 1.1 of \cite{LVWY}, Theorem \ref{thm:1.1}  and Theorem \ref{thm:1.4} of the present paper, we 
obtain a complete classification of the hypersurfaces in $\mathbb{S}^m\times \mathbb{S}^n$, 
$\mathbb{H}^m\times \mathbb{H}^n$ and $\mathbb{S}^m\times \mathbb{H}^n$
with constant sectional curvature, for $m,n\geq2$. 
Together with the results of \cite{AEG,CT2025,M-T}, one can get a complete classification of the hypersurfaces with constant sectional curvature in a product of two space forms. Note that the authors in \cite{CT2025} also claimed the related results (cf. \cite{KNT}). 
\end{remark} 














%%%%%%%%%%%%%%%%%%%%%%%%%%%%%%%%%%%%%%%%%%%%%%%%%%%%%%%%%%%%%%%%%%%%
\normalsize\noindent
\begin{thebibliography}{11}

\bibitem{AEG}
Juan A. Aledo, Jos\'{e} M. Espinar and Jos\'{e} A. G\'{a}lvez,
\emph{Complete surfaces of constant curvature in $\mathbb{S}^2\times \mathbb{R}$ and $\mathbb{H}^2\times \mathbb{R}$}, Calc. Var. Partial Differential Equations \textbf{29} (2007), no.~3, 347--363. \MR{2321892}

\bibitem{AD}
Mateo Anarella and Marie D'haene, 
\emph{Hypersurfaces of six-dimensional nearly K\"{a}hler manifolds}, 
arXiv:2507.22526

\bibitem{A-L-V-W}
Miroslava Anti\'{c}, Haizhong Li, Luc Vrancken and Xianfeng Wang, \emph{Affine
hypersurfaces with constant sectional curvature}, Pacific J. Math.
\textbf{310} (2021), no.~2, 275--302. \MR{4229240}

%\bibitem{B-C-O}
%J\"{u}rgen Berndt, Sergio Console and %Carlos~Enrique Olmos, \emph{Submanifolds
%and holonomy}, second ed., Monographs and %Research Notes in Mathematics, CRC
%Press, Boca Raton, FL, 2016. \MR{3468790}

%\bibitem{GHMY}
%Dong Gao, Zejun Hu, Hui Ma and Zeke Yao:
%On real hypersurfaces of $\mathbb{S}^2\times\mathbb{S}^2$.
%Proc. Amer. Math. Soc. to appear

\bibitem{CT}
Arnando Carvalho and Ruy Tojeiro, \emph{On a class of hypersurfaces of a product of two space forms}, arXiv:2602.16087v4 



\bibitem{C-U}
Ildefonso Castro and Francisco Urbano, \emph{Minimal {L}agrangian surfaces in
{$\mathbb{S}^2\times\mathbb{S}^2$}}, Comm. Anal. Geom. \textbf{15} (2007), no.~2,
217--248. \MR{2344322}

\bibitem{CR1982}
Thomas~E. Cecil and Patrick~J. Ryan, \emph{Focal sets and real hypersurfaces in
complex projective space}, Trans. Amer. Math. Soc. \textbf{269} (1982),
no.~2, 481--499. \MR{637703}


\bibitem{CT2025}
Arnando Carvalho and Ruy Tojeiro, \emph{Constant curvature hypersurfaces of cylinders
	over space forms}, 
arXiv:2510.17030.


\bibitem{DT2019}
Marcos Dajczer and Ruy Tojeiro,
\emph{Submanifold theory beyond an introduction},
Universitext, Springer, New York, 2019.
\MR{3969932}



\bibitem{DT2002}
Marcos Dajczer and Ruy Tojeiro,
\emph{An extension of the classical Ribaucour transformation},
Proc. London Math. Soc. (3) \textbf{85} (2002), no. 1, 211--232. \MR{1901374}


%\bibitem{D-V-W}
%Bart Dioos, Luc Vrancken and Xianfeng Wang, \emph{Lagrangian submanifolds in
%the homogeneous nearly {K}\"{a}hler {$\Bbb S^3\times\Bbb S^3$}}, Ann. Global
%Anal. Geom. \textbf{53} (2018), no.~1, 39--66. \MR{3746515}

\bibitem{GHMY}
Dong Gao, Zejun Hu, Hui Ma and Zeke Yao, \emph{On real hypersurfaces of
{$\mathbb{S}^2\times\mathbb{S}^2$}}, Proc. Amer. Math. Soc. \textbf{150} (2022),
no.~10, 4447--4461. \MR{4470187}

\bibitem{GMY}
Dong Gao, Hui Ma and Zeke Yao, \emph{On Hypersurfaces of $\mathbb{H}^2\times \mathbb{H}^2$},
Sci. China Math. \textbf{67} (2024), no.~2, 339--366.

\bibitem{GVWX}
Dong Gao, Joeri Van der Veken, Anne Wijffels and Botong Xu,
\emph{Lagrangian surfaces in $\mathbb{H}^2\times \mathbb{H}^2$},
Comm. Anal. Geom. \textbf{33} (2025), no.~4, 957--980.

\bibitem{HLYZ}
Zejun Hu, Xiaoge Lu, Zeke Yao and Xi Zhang, \emph{On curvature-adapted hypersurfaces of $\mathbb{S}^2\times \mathbb{S}^2$ and $\mathbb{H}^2\times \mathbb{H}^2$},
Internat. J. Math. \textbf{36} (2025), no.~2, Paper No. 2450072, 22 pp.

\bibitem{I-R}
Thomas~A. Ivey and Patrick~J. Ryan, \emph{Hopf hypersurfaces of small {H}opf
principal curvature in {$\Bbb CH^2$}}, Geom. Dedicata \textbf{141} (2009),
147--161. \MR{2520069}

\bibitem{K-R}
Hyang~Sook Kim and Patrick~J. Ryan, \emph{A classification of pseudo-{E}instein
hypersurfaces in {$\bold C{\rm P}^2$}}, Differential Geom. Appl. \textbf{26}
(2008), no.~1, 106--112. \MR{2393977}

\bibitem{KNT}
Jihun Kim, Yuri Nikolayevsky, and Ruy Tojeiro, \emph{Locally symmetric hypersurfacesin globally symmetric spaces}, Preprint.


\bibitem{Kon}
Masahiro Kon, \emph{Pseudo-{E}instein real hypersurfaces in complex space
forms}, J. Differential Geometry \textbf{14} (1979), no.~3, 339--354 (1980).
\MR{594705}

\bibitem{K}
Daniel Kowalczyk, \emph{Isometric immersions into products of space forms},
Geom. Dedicata \textbf{151} (2011), 1--8. \MR{2780734}

%\bibitem{K-V}
%Marcus Kriele and Luc Vrancken, \emph{An extremal class of three-dimensional
%hyperbolic affine spheres}, Geom. Dedicata \textbf{77} (1999), no.~3,
%239--252. \MR{1716673}

\bibitem{LMVVW}
Haizhong Li, Hui Ma, Joeri Van~der Veken, Luc Vrancken and Xianfeng Wang,
\emph{Minimal {L}agrangian submanifolds of the complex hyperquadric}, Sci.
China Math. \textbf{63} (2020), no.~8, 1441--1462. \MR{4125728}

\bibitem{LVWY}
Haizhong Li, Luc Vrancken, Xianfeng Wang and Zeke Yao,
\emph{Hypersurfaces of $\mathbb{S}^2\times\mathbb{S}^2$ with constant sectional curvature},
Calc. Var. Partial Differential Equations \textbf{63} (2024), no.~7, Paper No. 167, 33 pp.
\MR{4768511}

\bibitem{L-T-V}
Jorge H. Lira, Ruy Tojeiro and Feliciano A. Vit\'{o}rio,
\emph{A {B}onnet theorem for
isometric immersions into products of space forms}, Arch. Math. (Basel)
\textbf{95} (2010), no.~5, 469--479. \MR{2738866}

%\bibitem{L-W-W}
%Xiaoge Lu, Peng Wang and Xiaozhen Wang, %\emph{On the {G}eometry of
%{H}ypersurfaces in {$\mathbb S^2\times\mathbb %S^2$}}, Results Math. \textbf{78}
%(2023), no.~1, Paper No. 10. \MR{4510548}

%\bibitem{To}
%Francisco Torralbo, Francisco Urbano:
%Surfaces with parallel mean curvature vector in $\mathbb{S}^2\times\mathbb{S}^2$ and
%$\mathbb{H}^2\times\mathbb{H}^2$. Trans. Amer. Math. Soc. 364 (2012), no. 2, 785--813.

\bibitem{M-T}
Fernando Manfio and Ruy Tojeiro, \emph{Hypersurfaces with constant sectional curvature of
$\mathbb{S}^n\times \mathbb{R}$ and $\mathbb{H}^n\times \mathbb{R}$},
Illinois J. Math. \textbf{55} (2011), no.~1, 397--415. \MR{3006694}

\bibitem{Me-T}
Bruno Mendon\c{c}a and Ruy Tojeiro, \emph{Submanifolds of products of space forms},
Indiana Univ. Math. J. \textbf{62} (2013), no.~4, 1283--1314. \MR{3179691}

\bibitem{M}
Sebasti\'{a}n Montiel, \emph{Real hypersurfaces of a complex hyperbolic space},
J. Math. Soc. Japan \textbf{37} (1985), no.~3, 515--535. \MR{792990}

\bibitem{N-P}
Yuri Nikolayevsky and JeongHyeong Park, \emph{Einstein hypersurfaces in irreducible symmetric spaces},
Ann. Mat. Pura Appl. (4) \textbf{202} (2023), no.~4, 1719--1751. \MR{4597599}

\bibitem{T-U1}
Francisco Torralbo and Francisco Urbano, \emph{Surfaces with parallel mean
curvature vector in {${\Bbb S}^2\times{\Bbb S}^2$} and
{${\Bbb H}^2\times{\Bbb H}^2$}}, Trans. Amer. Math. Soc. \textbf{364} (2012), no.~2,
785--813. \MR{2846353}

\bibitem{T-U2}
Francisco Torralbo and Francisco Urbano, \emph{Minimal surfaces in {$\mathbb{S}^2\times\mathbb{S}^2$}}, J. Geom. Anal. \textbf{25} (2015), no.~2, 1132--1156. \MR{3319965}

\bibitem{Ur}
Francisco Urbano, \emph{On hypersurfaces of {$\mathbb S^2\times\mathbb S^2$}}, Comm.
Anal. Geom. \textbf{27} (2019), no.~6, 1381--1416. \MR{4052731}

\bibitem{ZGHY}
Xi Zhang, Dong Gao, Zejun Hu and Zeke Yao, \emph{On Hopf hypersurfaces of  {$\mathbb S^2\times\mathbb S^2$} and  {$\mathbb H^2\times\mathbb H^2$}}, J. Geom. Phys. \textbf{194} (2023), Paper No. 104996, 14 pp. \MR{4645631}



%\bibitem{Y-Y-H}
%Zeke Yao, Bangchao Yin and Zejun Hu, \emph{Non-existence of conformally flat
%real hypersurfaces in both the complex quadric and the complex hyperbolic
%quadric}, Canad. Math. Bull. \textbf{65} (2022), no.~1, 68--83. \MR{4396624}
\end{thebibliography}

\vskip 10mm

\begin{flushleft}
Haizhong Li\\
{\sc Department of Mathematical Sciences, Tsinghua University,\\
Beijing, 100084, P.R. China}\\
E-mail: lihz@tsinghua.edu.cn

\vskip 1mm

Luc Vrancken\\
{\sc
Univ. Polytechnique Hauts-de-France,
CERAMATHS-Laboratoire de Mat\'eriaux C\'{e}ramiques et de Math\'{e}matiques,
F-59313 Valenciennes, France;\\
INSA Hauts-de-France, CERAMATHS, F-59313 Valenciennes, France;\\
Department of Mathematics, KU Leuven, Celestijnenlaan 200B, Box 2400, BE-3001 Leuven, Belgium}\\
E-mail: luc.vrancken@uphf.fr

\vskip 1mm

Xianfeng Wang\\
{\sc School of Mathematical Sciences and LPMC, Nankai University,\\
Tianjin 300071, P.R. China}\\
E-mail: wangxianfeng@nankai.edu.cn

\vskip 1mm

Zeke Yao\\
{\sc School of Mathematical Sciences, South China Normal University,\\
Guangzhou 510631, P.R. China}\\
E-mail: yaozkleon@163.com

\end{flushleft}



\end{document}	
